% Original research notes for integration by the coordinating agent.
% No source-paper text is copied into this file.
\section{Exponent audit and finite-field descent}
\label{sec:exponent-audit}

The statements in this section distinguish an optimization of a specified
argument from a limitation on every conceivable direct argument.  The latter
does not follow from the two attached manuscripts.  Source locators below refer
to the attached versions, treated throughout as unrefereed claims.

\subsection{The two exponents actually established by the stated estimates}
\label{sec:ex-stated}

\paragraph{EX-01: optimization of the complete-network estimates.}
Anthropic, Appendix E, Proposition E.10 (pp.~36--37), supplies, under its flat
hypothesis, the estimates
\[
 \log L\geq c\ell k^{3/2},\qquad \ell\geq C\sqrt{k}.
\]
Lemma E.3 (p.~33) embeds the network polynomial in a permanent whose logarithmic
size is $k\ell+O(\log k+\log\ell)$.  Put $B=k\ell$.  Then
$B\geq Ck^{3/2}$, and consequently
\[
 \ell k^{3/2}=B\sqrt{k}\leq C^{-1/3}B^{4/3}.
\]
Conversely, taking $k=r^2$ and $\ell=Cr$ makes the left-hand expression a
constant multiple of $B^{4/3}$.  Integer rounding changes only constants.
Thus $4/3$ is the exact optimum of these parameter inequalities.  This is a
parameter-optimization statement, not an upper bound on the actual circuit
complexity and not a universal obstruction to improving the method.  Anthropic,
``Remarks and open questions,'' p.~61, explicitly leaves improvement open.

\paragraph{EX-02: optimization of the expander estimates.}
In OpenAI, Section~2.1, source lines 233--275, the bundle width is
$m=\Theta(\log n)$, the number of vertices is $d$, and a bounded-degree expander
has $D=\Theta(d)$ edges.  The first-order resultant has size
$\Omega(\sqrt d)$, so a bridge with $h=O(m\sqrt d)$ permits a saving of
$\Omega(m)$ by taking a displacement $u=\Theta(d^{-1/2})$.  The charged selector
budget is $2D\log_2 A=O(d^2)$ with $A\leq2^d$.  Hence $d=\Theta(\sqrt m)$
gives $h=\Theta(m^{5/4})$.  This checks the exponent asserted in prompt Fact~6.
It does not identify $5/4$ as an optimum among more general geometries.

\subsection{An exact limitation for the orthogonal two-hub model}
\label{sec:ex-orthogonal}

\paragraph{Model.}
Let $G$ be a finite graph on $k$ vertices, with $D$ edges and maximum degree
$\Delta$.  Every edge is an independent coordinate diameter with equal width
$m$.  For a balanced cut, all terminals on its first side are connected to one
hub and all terminals on its other side to a second hub.  A width-$h$ bridge
joins the hubs.  The selector charge is the definition's $2D\log_2 A$, and we
assume $\log_2 A\geq ak$ for a fixed $a>0$; this includes paying for all balanced
cuts.  The hubs may move arbitrarily in the unit ball.  No Taylor approximation
is built into this model.

\paragraph{EX-03: exact two-hub geometric bounds.}
For a cut with $c$ crossing edges, both fan-ins together save at most $2m$,
even if the bridge is free.  Moreover, if $h\geq m\sqrt c$, the entire two-hub
layout has cost at least $2Dm$ for every pair of hub positions.

\emph{Proof.}
Every edge internal to a cut side contributes both endpoints $e_j,-e_j$ to
that side.  Their distances to any hub add to at least $2$, by the triangle
inequality.  The $c$ crossing endpoints on a given side are mutually orthogonal
signed unit vectors.  For $c\geq1$, their minimum sum of distances to a point
is $\sqrt{c(c-1)}$.  Indeed, projection onto their span cannot increase any
distance, and convexity permits averaging over all permutations of the
orthogonal coordinates.  An averaged minimizer has all coordinates equal to
$t$, where its objective is
$c\sqrt{1-2t+ct^2}$, minimized at $t=1/c$.  Thus one side saves at most
$m(c-\sqrt{c(c-1)})\leq m$.  A cut with no crossing edge saves nothing.

For the bridge assertion, let $r$ be the sum of the unit endpoint vectors on
the first side.  Internal pairs cancel, so $\|r\|=\sqrt c$, and the sum on the
second side is $-r$.  For any unit vector $p$, convexity gives
$\|p-x\|\geq1-p\cdot x$.  Applying this at hub positions $x,y$ and adding the
bridge gives total cost at least
\[
 2Dm-mr\cdot(x-y)+h\|x-y\|
 \geq 2Dm+(h-m\sqrt c)\|x-y\|.
\]
This proves both assertions.\hfill$\square$

\paragraph{EX-04: exact $5/4$ exponent barrier in this model.}
Let $B=m\Delta$ be the largest local bit load.  A strict shortening through this
model requires $m>D\log_2 A\geq aDk$, because its total fan-in saving is at most
$2m$.  Since $\Delta k\geq2D$,
\[
 B=m\Delta> aDk\Delta\geq2aD^2.
\]
EX-03 says that the bridge width up to which a geometric contradiction can be
obtained is no larger than $m\sqrt D$.  Therefore
\[
 m\sqrt D\leq B\left(\frac{B}{2a}\right)^{1/4}
 =O(B^{5/4}).
\]
The graph polynomial has at least $2^B$ local variables, so $B$ is at most the
logarithm of any explicit ambient variable budget.  Bounded-degree expanders
and the displacement calculation in EX-02 attain this exponent.  Consequently
$5/4$ is optimal for this stated model.  The result does not cover eliminating
most cut layouts, using more than two hubs, changing the target diameters, or
using a different selector accounting rule.

\subsection{The midpoint bound is optimal among binary-seed certificates}
\label{sec:ex-midpoint}

Let $P_k=\{\pm e_1/2,\ldots,\pm e_k/2\}$.  A binary-seed convexity certificate
for $P_k$ is a finite probability distribution on maps $X:P_k\to\mathbb R^k$
whose images have at most two points, such that $\mathbb E X(p)=p$ for every
$p\in P_k$.  Its cost is $\mathbb E\operatorname{diam}X(P_k)$.  Iterating
midpoint convexity from two-point configurations is a special case (with dyadic
weights); limits can be included without changing the following infimum.

\paragraph{EX-05: the minimum binary-seed cost is exactly $\sqrt k$.}
\emph{Proof.}
For a particular two-point image, write
$X(e_i/2)-X(-e_i/2)=t_i v$, where $t_i\in\{-1,0,1\}$ and $\|v\|$ is its
diameter.  The barycenter identities yield
\[
 k=\sum_{i=1}^k\left\langle e_i,
   \mathbb E[X(e_i/2)-X(-e_i/2)]\right\rangle
 =\mathbb E\sum_i t_i v_i
 \leq\sqrt{k}\,\mathbb E\|v\|.
\]
This is the lower bound.  For equality, take $\varepsilon$ uniformly from
$\{-1,1\}^k$ and set
\[
 X_\varepsilon(\pm e_i/2)=\pm\varepsilon_i\varepsilon/2.
\]
Its mean is $\pm e_i/2$ and its image diameter is $\sqrt k$.\hfill$\square$

Anthropic's Lemma E.8 and Proposition E.9 (pp.~35--36) use precisely such
midpoint convexity to upper-bound the scaled multipartite unit-tensor exponent.
EX-05 proves the optimality of that certificate route.  It does \emph{not}
prove that every spectral point realizing these wire distances has scaled
unit-tensor exponent at least $\sqrt k$.  Conflating these statements would
leave a gap in an asserted universal $4/3$ barrier.

\subsection{Strict shortening depends only on the characteristic}
\label{sec:ex-field}

\paragraph{EX-06: descent across a finite field extension.}
Let $E/F$ be a finite extension.  If a strict shortening exists over $E$, then
one exists over $F$ (possibly at a larger cable width and with more summands).

\emph{Proof.}
Fix a shortening of $C_{D,m}$ with $A$ summands, cost at most $w$, and gap
\[
 \delta=2Dm-w-2D\log_2 A>0.
\]
Take the finite union of the layout vertices of all summands, adding trivial
one-entry local tables at otherwise unused vertices.  Distinct coincident
vertices may be identified, since local computation is free.  Fix a spanning
tree on this common vertex set, including the terminal vertices.  All its
segments lie in the ball.  Write $r=[E:F]$, choose an $F$-basis of $E$ containing
$1$, and an $F$-linear functional $\ell:E\to F$ with $\ell(1)=1$.

The $t$-th tensor power of the identity has $A^t$ summands.  Each is a network
on the common vertex set with the $t$ copies' incident ports reblocked; its cost
is at most $tw$.  At every vertex its local table entry is a single element of
$E$, namely the product of the $t$ entries assigned there.  Express this element
in the chosen basis.  On the fixed auxiliary tree, use the multiplication
structure constants of $E$ over $F$ to transmit the product of these local
elements from each subtree toward the root.  The state space of each auxiliary
edge is $E$ as an $r$-dimensional $F$-vector space.  Realize it with
$\lceil\log_2 r\rceil$ bits and set the unused coordinate entries to zero.
At the root apply $\ell$.

Contracting this auxiliary network computes the $F$-linear functional $\ell$
of the original coefficient product.  Because $\ell$ commutes with addition,
the resulting sum of $A^t$ networks is
$\ell(C_{D,m}^{\otimes t})=C_{D,tm}$: target entries are $0$ or $1$ and
$\ell(1)=1$.  The extra cost is the fixed number
\[
 K=\lceil\log_2 r\rceil
       \sum_{e\ \mathrm{in\ the\ auxiliary\ tree}}\operatorname{length}(e),
\]
independent of $t$ and of the term word.  Thus each new summand costs at most
$tw+K$.  For an integer $t>K/\delta$,
\[
 tw+K+2D\log_2(A^t)<2Dtm.
\]
This is a strict shortening over $F$.  The argument uses an arbitrary linear
functional with $\ell(1)=1$, not the field trace, and therefore also covers
inseparable extensions and degrees divisible by the characteristic.
\hfill$\square$

\paragraph{EX-07: characteristic invariance.}
For arbitrary fields $F$ and $K$ of the same characteristic, a strict shortening
exists over $F$ if and only if one exists over $K$.

\emph{Proof.}
One direction of scalar extension is immediate.  For the converse, let $P$ be
the prime field and suppose an identity exists over $F$.  Its finite layouts
and bit widths are fixed.  Introduce one indeterminate for every local table
entry.  Equating each tensor entry of the sum to the corresponding $0$ or $1$
target entry gives finitely many polynomial equations over $P$.  They have an
$F$-solution, so the generated ideal is proper.  Choose a maximal ideal above
it.  Zariski's lemma says the quotient is a finite algebraic extension $E/P$,
yielding the same identity over $E$.  Layout geometry and the strict cost gap
have not changed.  Apply EX-06 to descend to $P$, then extend scalars to $K$.
\hfill$\square$

\paragraph{EX-08: characteristic-zero shortening transfers to almost all
primes.}
If a strict shortening exists over $\mathbb C$, it exists over $\mathbb Q$ by
EX-07.  Fix that rational identity and multiply together the finitely many
nonzero denominators of its local table entries, obtaining $N\neq0$.
For every prime $p\nmid N$, reduce the tables in $\mathbb Z[1/N]$ modulo $p$.
Every tensor identity is preserved, while all geometric lengths, widths, and
the strict margin remain unchanged.  This produces a shortening over
$\mathbb F_p$.  The assertion gives no conclusion at a specified divisor of
$N$, in particular at $2$.  It does not assert that a shortening modulo $2$
lifts to characteristic zero.

\subsection{Audit of the supplied separation bridges}
\label{sec:ex-bridges}

\paragraph{EX-09: B7.1--B7.3 have a fixed-instance scope.}
The attached \path{ALPOGE_GCT_bridges.tex}, B7.1 (lines 359--401), gives a
single integral separating polynomial for a fixed pair of integral points in
an integral representation.  Its nonzero value at the excluded point controls
the exceptional primes.  B7.2 (lines 403--447) uses two fixed integer matrices
at a fixed degree; preserving chosen nonzero rank-size minors preserves the
separation criterion.  The linear-algebra proof and the torsion-free
localization argument have been checked and no gap was found in those stated
claims.  B7.3 (lines 449--457) correctly exhibits failure to lift: the scaling
orbit of $2$ in $\mathbb A^1$ is dense in characteristic zero but becomes the
zero orbit in characteristic two, excluding the point $1$.

These statements do not identify a good-reduction certificate at $p=2$, do not
bound exceptional primes uniformly in a growing family, and do not themselves
transfer universal assertions about all spectral points.  EX-06--EX-08 concern
finite shortening identities and avoid claiming any of these stronger
conclusions.

\subsection{Prices also depend only on characteristic}
\label{sec:ex-price-field}

Here a normalized spectral point means a restriction-monotone, additive and
multiplicative map on the tensor semiring, taking the rank-$r$ full unit tensor
to $r$.  We use the normalization lemma proved in the shortening section:
every nontrivial spectral point of Anthropic Definition A.1 is a positive
power of such a normalized point.  Thus replacing generalized points by
normalized points does not change the closed cone of scaled wire metrics.

\paragraph{EX-13: constant-cover descent of tensor restrictions.}
Let $F\subset E$ be fields, and let $T,S$ be $k$-party tensors over $F$.
If $T\leq_E S$, then there is a constant $R\geq1$ such that, for every
$n\geq1$, $T^{\otimes n}$ is covered over $F$ by $R$ copies of
$S^{\otimes n}$.

\emph{Proof.}
The local matrices witnessing the restriction solve a finite system of
polynomial equations over $F$.  A solution over $E$ makes its ideal proper;
a maximal ideal and Zariski's lemma produce a solution over a finite
extension $K/F$.  Write $r=[K:F]$, choose an $F$-basis $b_1,\ldots,b_r$ of $K$,
and an $F$-linear functional $\ell:K\to F$ with $\ell(1)=1$.
If $A_1,\ldots,A_k$ give the restriction over $K$, expand each whole matrix
$A_i^{\otimes n}$, not each tensor-power factor separately, in that basis:
\[
 A_i^{\otimes n}=\sum_{a=1}^r b_a B_{i,a}^{(n)},
\]
where all $B_{i,a}^{(n)}$ have entries in $F$.  Applying $\ell$ entrywise to
the restriction identity, and using that $T,S$ have entries in $F$, gives
\[
 T^{\otimes n}=\sum_{a_1,\ldots,a_k=1}^r
 \ell(b_{a_1}\cdots b_{a_k})
 (B_{1,a_1}^{(n)}\otimes\cdots\otimes B_{k,a_k}^{(n)})S^{\otimes n}.
\]
Every summand is an $F$-restriction of $S^{\otimes n}$, with its scalar
absorbed at one party.  Hence $R=r^k$ works, independently of $n$.
\hfill$\square$

\paragraph{EX-14: surjectivity of scalar-extension restriction on spectra.}
For every field extension $E/F$, every normalized spectral point on all
$F$-tensors extends, in the sense of equal values on scalar-extended
$F$-tensors, to a normalized spectral point on all $E$-tensors.

\emph{Proof.}
Let $\psi$ be the given point.  EX-13 and its cover bound imply
\[
 \psi(T)^n\leq R\psi(S)^n
 \quad\Longrightarrow\quad \psi(T)\leq\psi(S)
\]
whenever $T\leq_E S$ are defined over $F$.  Therefore $\psi$ is monotone for
the inherited $E$-restriction preorder on the image of the $F$-tensor
semiring.  Mutual $E$-restriction forces equality of its values, so the map
is well defined even after passing to restriction-equivalence classes.
This image is a unital subsemiring: direct sums, tensor products, and unit
tensors remain defined over $F$.  Its inherited preorder is a Strassen
preorder; in particular, a nonzero tensor has the scalar unit as a restriction,
and is bounded above by a finite full unit tensor using its finite tensor
rank.  Natural-number comparisons are the usual ones by flattening rank.

Apply the spectral extension theorem, in the precise form of Wigderson and
Zuiddam (2025), Section~4.3, Theorem~4.16, p.~53 (proof pp.~56--57): a
monotone semiring homomorphism on a unital subsemiring with the inherited
Strassen preorder extends to the full semiring.  It supplies the desired
$E$-spectral point.\hfill$\square$

\paragraph{EX-15: characteristic invariance of price cones.}
For every fixed finite party set, the sets of wire metrics of normalized
spectral points are identical over fields of the same characteristic.
Consequently the price cones, the flat hypothesis, and the contracting
hypothesis depend only on the characteristic.

\emph{Proof.}
Restriction of an $E$-spectral point to scalar-extended $F$-tensors is a
normalized $F$-spectral point.  EX-14 gives the reverse surjectivity, preserving
every binary wire, whose entries belong to the prime field.  Thus the wire
metric sets agree across $E/F$.  Any two fields of the same characteristic
contain their common prime field, so the equality holds between them.
The normalization lemma, scaling, and pointwise closure give the assertion
for prices.  Flatness and contraction quantify solely over these price cones.
\hfill$\square$

This does not identify the price cones in characteristic zero and
characteristic two.  It establishes that adjoining algebraic or
transcendental elements within either characteristic cannot change them.

\subsection{A broader exact barrier for opposite-hub constructions}
\label{sec:ex-broad-barrier}

The next three statements permit arbitrary finite weighted graphs and arbitrary
antipodal unit endpoint directions.  They do not assume orthogonal terminals,
equal cable widths, regularity, or infinitesimal hub displacement.  Their
restriction is the selector certificate and the use of opposite hubs, which
are specified below rather than tacitly identified with every possible
``hub-displacement method.''

We use the following classical vector-balancing theorem: if vectors have
Euclidean norms at most $1/5$ and a convex body has standard Gaussian measure
at least $1/2$, then some signed sum of those vectors lies in the body
(Banaszczyk, 1998; exact formulation in Dadush, Garg, Lovett, and Nikolov,
2019, Theorem~1.1, p.~2).  This is an external established theorem, not a
consequence of either attached manuscript.

\paragraph{Model and notation.}
Let $G$ have $k\geq2$ vertices.  An edge $e$ has a bit width $b_e\geq0$ and
endpoints placed at $x_e,-x_e$ in a real Euclidean space, where $\|x_e\|=1$.
Choose either orientation of each edge, assigning $x_e$ to its tail.  Write
\[
 d_a=\sum_{e\ni a}b_e,\qquad
 B=\max_a d_a>0,\qquad W=\sum_e b_e>0,
 \qquad v_a=\sum_{e\ni a}\sigma_{a,e}b_e x_e.
\]
Thus $\sum_a v_a=0$, $\|v_a\|\leq d_a\leq B$, and $2W\leq kB$.
A cut $S$ is balanced if $k/3\leq|S|\leq2k/3$, and its resultant is
$R_S=\sum_{a\in S}v_a$.  All terminals belonging to vertices in $S$ connect
to a hub $z$, and all other terminals to $-z$.  The baseline cable cost is
$2W$.  A bridge of width $h$ adds $2h\|z\|$.

The endpoint set $P=\{x_e,-x_e:e\in E(G)\}$ has a binary-seed certificate
if there is a finite probability distribution on maps $X:P\to\mathbb R^d$
such that $|X(P)|\leq2$ and $\mathbb EX(p)=p$.  Its cost is
$C=\mathbb E\operatorname{diam}X(P)$.  This is the same certificate notion
as EX-05, with a different radius normalization.  The model charges a selector
debit at least $akC$, for a fixed $a>0$, and requires a strict cost improvement
after that debit for every balanced cut.  This describes the certificate
accounting in which all balanced cuts are paid for and the scaled selector
exponent is bounded through binary averaging.  It is not an assertion that
the actual selector exponent is bounded below by $C$.

\paragraph{EX-16: a small balanced resultant from a binary certificate.}
For $k\geq43$ there is a balanced cut $S$ such that
\[
 \|R_S\|\leq \frac{5\sqrt2}{2}BC.
\]
For every $k\geq2$ there is also a balanced cut, possibly different, satisfying
\[
 \|R_S\|^2\leq BW\leq \tfrac12 kB^2.
\]

\emph{Proof.}
In sample $\omega$ let the two image points differ by $q_\omega$, with
probability $p_\omega$.  A one-point image has $q_\omega=0$.  The equality
$\mathbb E[X(x_e)-X(-x_e)]=2x_e$ gives
\[
 x_e=\sum_\omega t_{e,\omega}g_\omega,
 \qquad t_{e,\omega}\in\{-1,0,1\},
 \qquad g_\omega=p_\omega q_\omega/2.
\]
Put $L=\sum_\omega\|g_\omega\|=C/2$.  Since $d_a\leq B$, each $v_a/B$
belongs to the zonotope $Z=\sum_\omega[-g_\omega,g_\omega]$.  Define, on the
span of its nonzero generators,
\[
 Q=L\sum_{g_\omega\ne0}
       \frac{g_\omega g_\omega^{\mathsf T}}{\|g_\omega\|}.
\]
Then $Q$ is positive definite on that span and $\operatorname{tr}Q=L^2$.
For every vector $u$, Cauchy--Schwarz gives
\[
 h_Z(u)^2=\left(\sum_\omega|u\cdot g_\omega|\right)^2
 \leq L\sum_\omega\frac{(u\cdot g_\omega)^2}{\|g_\omega\|}
 =u^{\mathsf T}Qu.
\]
Comparison of support functions therefore gives
$Z\subseteq Q^{1/2}\mathbb B_2$.

The vectors
\[
 \frac{1}{5\sqrt2}\left(Q^{-1/2}v_a/B,1\right)
\]
have norm at most $1/5$.  The convex body
\[
 K=\{(y,t):\|Q^{1/2}y\|\leq2L,\ |t|\leq2\}
\]
has Gaussian measure at least $(3/4)^2>1/2$: the two Gaussian coordinates
are independent, and Markov's inequality applies to
$\mathbb E\|Q^{1/2}Y\|^2=L^2$ and $\mathbb ET^2=1$.
Banaszczyk's theorem supplies signs $\epsilon_a$ with
\[
 \left\|\sum_a\epsilon_a v_a\right\|
       \leq10\sqrt2 BL=5\sqrt2 BC,
 \qquad \left|\sum_a\epsilon_a\right|\leq10\sqrt2.
\]
For $k\geq43$ the positive-sign set is balanced, and its resultant is half
the signed sum because $\sum_a v_a=0$.  This proves the first bound.

For the second, choose $S$ uniformly among sets of size
$s=\lfloor k/2\rfloor$.  Direct expansion, using $\sum_a v_a=0$, yields
\[
 \mathbb E\|R_S\|^2
 =\frac{s(k-s)}{k(k-1)}\sum_a\|v_a\|^2
 \leq\frac{k}{4(k-1)}\sum_a B d_a
 =\frac{k}{2(k-1)}BW\leq BW.
\]
Every such set is balanced.\hfill$\square$

\paragraph{EX-17: some balanced cut has only $O(B)$ exact hub saving.}
For every graph and endpoint placement above, there is a balanced cut for
which the total fan-in saving is at most $472B$, for every opposite pair of
hub positions, even with a free bridge.  The constant is not optimized.

\emph{Proof.}
Define the positive semidefinite matrix
\[
 A=\sum_e b_e x_e x_e^{\mathsf T},\qquad \operatorname{tr}A=W.
\]
Let $u$ be a unit eigenvector corresponding to its largest eigenvalue, and
write $v_a=\alpha_a u+w_a$, where $w_a\perp u$.  Let $A_\perp$ denote the
compression of $A$ to $u^\perp$, restricted to its range.  Weighted
Cauchy--Schwarz at vertex $a$ gives the matrix inequality
\[
 w_aw_a^{\mathsf T}
 \preceq d_a\sum_{e\ni a}b_e
             (\operatorname{proj}_{u^\perp}x_e)
             (\operatorname{proj}_{u^\perp}x_e)^{\mathsf T}
 \preceq B A_\perp.
\]
Consequently $\|A_\perp^{-1/2}w_a\|\leq\sqrt B$, and
$|\alpha_a|\leq B$.  If $A_\perp=0$, all $w_a$ vanish and the corresponding
coordinates below are omitted.  Otherwise put $T=\operatorname{tr}A_\perp$.

Apply Banaszczyk's theorem to
\[
 z_a=\frac{1}{5\sqrt3}
       \left(\alpha_a/B,A_\perp^{-1/2}w_a/\sqrt B,1\right),
\]
whose norms are at most $1/5$, and to the ellipsoid
\[
 K=\left\{(s,y,t):
 s^2+\frac{\|A_\perp^{1/2}y\|^2}{T}+t^2\leq6\right\}.
\]
The middle term is omitted when $T=0$.  The Gaussian expectation of the
left-hand side is at most $3$, so $K$ has Gaussian measure at least $1/2$.
There are signs satisfying
\[
 \left|\sum_a\epsilon_a\right|\leq15\sqrt2,
 \quad \left|\sum_a\epsilon_a\alpha_a\right|\leq15\sqrt2 B,
 \quad \left\|\sum_a\epsilon_a w_a\right\|^2\leq450BT.
\]
For $k\geq64$ their positive-sign set is balanced.  Its resultant
$R=R_S=R_\parallel u+R_\perp$ therefore satisfies
\[
 |R_\parallel|\leq\frac{15\sqrt2}{2}B,
 \qquad \|R_\perp\|^2\leq\frac{225}{2}BW.
\]

It remains to turn these resultant bounds into an exact, global bound on
hub savings.  Projection onto the unit ball does not increase distances to
the endpoints; radial projection also preserves the form of an opposite
hub pair.  Thus it suffices to take $\|z\|\leq1$.  For a unit vector $p$,
\[
 \|p-z\|-(1-p\cdot z)
 =\frac{\|z\|^2-(p\cdot z)^2}
        {\|p-z\|+1-p\cdot z}
 \geq\frac{\|z\|^2-(p\cdot z)^2}{4}.
\]
The equality is interpreted by continuity if numerator and denominator are
both zero.  The denominator is at most $4$, and the numerator is nonnegative.
Summing this lower bound over the two fan-ins gives their exact cost $F_S(z)$
the bound
\[
 F_S(z)\geq2W-2R\cdot z+\tfrac12z^{\mathsf T}(WI-A)z.
\]
Every eigenvalue of $A$ except its largest is at most $W/2$, since two
eigenvalues larger than $W/2$ would contradict $\operatorname{tr}A=W$.
Writing $z=z_\parallel u+z_\perp$ and discarding the nonnegative quadratic
term in the $u$ direction, one obtains
\[
 \begin{split}
 2W-F_S(z)
 &\leq2|R_\parallel|
       +2R_\perp\cdot z_\perp-\frac W4\|z_\perp\|^2\\
 &\leq2|R_\parallel|+\frac4W\|R_\perp\|^2
 \leq(15\sqrt2+450)B<472B.
 \end{split}
\]
For $2\leq k<64$, the trivial upper bound on any saving is
$2W\leq kB<64B$, so any balanced cut suffices.\hfill$\square$

\paragraph{EX-18: exact $4/3$ barrier in the stated binary-selector model.}
Suppose every balanced cut admits opposite hubs for which the fan-ins,
a width-$h$ bridge, and a selector debit at least $akC$ together cost less
than $2W$.  Then, as $B\to\infty$,
\[
 h=O_a(B^{4/3}).
\]
This exponent is attained, conditional on the flat hypothesis and the
attached circuit-to-layout reductions, by the complete-network construction.

\emph{Proof.}
EX-17 gives the necessary condition $akC<472B$.  For any cut and any hubs
$x,y$, not necessarily opposite, the supporting-hyperplane inequality
$\|p-x\|\geq1-p\cdot x$ gives the lower bound
\[
 \text{fan-in cost}+h\|x-y\|
 \geq2W-R_S\cdot(x-y)+h\|x-y\|
 \geq2W+(h-\|R_S\|)\|x-y\|.
\]
Thus strict improvement for every balanced cut requires
$h<\min_S\|R_S\|$.  For $k\geq43$, EX-16 gives
\[
 h\leq B\min\left\{\sqrt{k/2},\frac{5\sqrt2}{2}C\right\}
 \leq B\min\left\{\sqrt{k/2},
                  \frac{1180\sqrt2}{a}\frac Bk\right\}.
\]
If $k\leq B^{2/3}$ the first term is $O(B^{1/3})$; if
$k>B^{2/3}$ the second is $O_a(B^{1/3})$.  This proves the stated bound.
For $k<43$ the random-cut bound alone gives $h=O(B)$.

For attainment take the complete graph, common width $\ell$, and endpoint
direction $x_{ab}=e_a$ for $a<b$, with the positive endpoint assigned to
vertex $a$.  EX-05, rescaled by $2$, supplies $C\leq2\sqrt k$.
Every balanced cut has $\|R_S\|=\Omega(\ell k^{3/2})$, by the exact
Cauchy--Schwarz calculation in Anthropic Proposition~E.10, pp.~36--37.
The quadratic upper bound on the fan-in cost, now used only to establish
attainment, gives a saving $\Omega(\ell k)$ at displacement
$\Theta(k^{-1/2})$ while accommodating bridge width
$h=\Theta(\ell k^{3/2})$.  Taking $\ell$ a sufficiently large fixed
multiple of $\sqrt k$ pays $akC$, while
$B=(k-1)\ell=\Theta(k^{3/2})$ and $h=\Theta(k^2)=\Theta(B^{4/3})$.
Metric approximation errors can be made smaller than the fixed positive
margins for each finite configuration.\hfill$\square$

The upper bound in EX-18 is purely geometric once its model is specified.
Its lower-bound application still uses the unrefereed flat-case machinery.
The next result also removes the opposite-hub condition, at the expense of
a logarithmic factor that does not affect the supremum of power exponents.

\paragraph{EX-19: arbitrary two hubs have the same exponent supremum.}
Keep the weighted graph, antipodal unit endpoints, and binary-seed selector
debit at least $akC$ from the model above.  Allow the two hubs for each
balanced cut to be placed independently.  Assume every nonzero edge width
$b_e$ is a positive integer.  If every balanced cut can be shortened after
paying the selector and a width-$h$ bridge, then
\[
 h=O_a\!\left(B^{4/3}\bigl(\log(2+B)\bigr)^{1/3}\right).
\]
Consequently the supremum of achievable fixed power exponents in this
precisely specified two-hub certificate model is $4/3$.  Attainment at
$4/3$ is the conditional flat-case construction of EX-18.  The upper bound
does not assert an exact constant-times-$B^{4/3}$ bound for independently
placed hubs.

\emph{Proof.}
We first establish a geometric lemma, allowing arbitrary positive real
widths: with $d$ the dimension of the span of all endpoint directions, some
balanced cut has fan-in saving at most
\[
 2200B\log(16d)
\]
uniformly over all independently placed hubs.  All logarithms in this
proof are natural; changing their base changes constants only.

Retain $A=\sum_e b_ex_ex_e^{\mathsf T}$ and put $Q=WI-A$.
Choose a principal eigenvector $u$ of $A$, write
$v_a=\alpha_a u+w_a$, and define
\[
 Q_a=\sum_{e\ni a}b_e(I-x_ex_e^{\mathsf T}).
\]
Then $0\preceq Q_a\preceq BI$ and $\sum_aQ_a=2Q$.
If $k$ is odd, adjoin one dummy vertex with $v_a=0$ and $Q_a=0$ for the
following pairing argument.  Order the vertices by increasing $\alpha_a$,
pair consecutive vertices, and give the two vertices of pair $j$
opposite signs $\epsilon_j,-\epsilon_j$, where the $\epsilon_j$ are
independent uniform signs.  On the original vertices the positive-sign
set always has size $\lfloor k/2\rfloor$ or $\lceil k/2\rceil$, hence is
balanced.  Since the sum of the adjacent gaps is no larger than the full
range of the ordered numbers,
\[
 \left|\sum_a\sigma_a\alpha_a\right|\leq
       \max_a\alpha_a-\min_a\alpha_a\leq2B.
\]
Thus its resultant satisfies $|R_\parallel|\leq B$ deterministically.
Moreover,
\[
 \mathbb E\|R_\perp\|^2
 =\tfrac14\sum_{\text{pairs }(a,b)}\|w_a-w_b\|^2
 \leq\tfrac12\sum_a\|w_a\|^2
 \leq BW.
\]
The last inequality follows from weighted Cauchy--Schwarz at every vertex,
as in EX-17.  Therefore
$\|R_\perp\|^2\leq4BW$ with probability at least $3/4$.

Let $D_j=Q_a-Q_b$ for pair $j$ and $Y=\sum_j\epsilon_jD_j$.
Writing $Q_S=\sum_{a\in S}Q_a$, one has
\[
 Q_S=Q+\tfrac12Y,\qquad Q_{S^c}=Q-\tfrac12Y.
\]
The elementary inequality $(H-K)^2\preceq2(H^2+K^2)$ for self-adjoint
matrices, together with $Q_a^2\preceq BQ_a$, gives
\[
 \sum_jD_j^2\preceq2B\sum_aQ_a=4BQ.
\]
We use the standard matrix Rademacher-series tail estimate
\[
 \Pr\!\left(\left\|\sum_j\epsilon_jD_j\right\|\geq t\right)
 \leq2d\exp\!\left(-\frac{t^2}{2\|\sum_jD_j^2\|}\right).
\]
The exact cited formulation is Tropp (2012), Theorem~4.1,
inequality~(4.4), author preprint p.~14; Theorem~1.2, p.~4,
is its one-sided version.

Put $L=\log(16d)$.  If $d=1$, the pairing above gives $\|R\|\leq B$,
and the supporting-hyperplane lower bound implies that the free fan-in
saving is at most $2B$ for hubs in the unit ball.  If $W<128BL$, the
trivial bound $2W<256BL$ suffices.  We may therefore assume $d\geq2$
and $W\geq128BL$.

On $u^\perp$, every eigenvalue of $Q$ is at least $W/2$.
Let $P$ be the orthogonal projection onto this subspace.
Compression satisfies $(PD_jP)^2\preceq PD_j^2P$, so the variance
parameter of $PYP$ is at most $4BW$.  The tail estimate gives
\[
 \Pr(\|PYP\|>W/2)
 \leq2d\exp(-W/(32B))<1/8.
\]
On the complementary event,
\[
 Q_S|_{u^\perp}\succeq\tfrac14WI,\qquad
 Q_{S^c}|_{u^\perp}\succeq\tfrac14WI.
\]

Write $q=u^{\mathsf T}Qu=W-\lambda_{\max}(A)$, which is the smallest
eigenvalue of $Q$.  If $q\geq128BL$, apply the same tail estimate to
$Q^{-1/2}D_jQ^{-1/2}$.  Its variance parameter is at most $4B/q$:
\[
 \sum_j(Q^{-1/2}D_jQ^{-1/2})^2
 \preceq q^{-1}Q^{-1/2}\Bigl(\sum_jD_j^2\Bigr)Q^{-1/2}
 \preceq(4B/q)I.
\]
It follows that
\[
 \Pr(\|Q^{-1/2}YQ^{-1/2}\|>1)
 \leq2d\exp(-q/(8B))<1/8.
\]
Hence a cut exists for which
$Q_S,Q_{S^c}\succeq Q/2$ and simultaneously
$|R_\parallel|\leq B$, $\|R_\perp\|^2\leq4BW$.

Project the two hubs independently onto the unit ball.  This cannot
increase a fan-in distance, and metric projection is nonexpansive, so
it cannot increase the bridge distance either.  For their projected
positions $x,y$, the exact norm inequality from EX-17 yields
\[
 F_S(x,y)\geq2W-R\cdot(x-y)
       +\tfrac14(x^{\mathsf T}Q_Sx+y^{\mathsf T}Q_{S^c}y).
\]
In the present case, discard the nonnegative quadratic terms in the
$u$ direction and use $Q|_{u^\perp}\succeq WI/2$ to obtain
\[
 2W-F_S(x,y)
 \leq2|R_\parallel|+\frac8W\|R_\perp\|^2
 \leq34B.
\]

It remains to consider $q<128BL$.  Choose a cut satisfying the
perpendicular concentration event and
$\|R_\perp\|^2\leq4BW$; their probabilities ensure a nonempty
intersection.  For $H=Q_S$ or $Q_{S^c}$,
\[
 0\preceq H\preceq2Q,\quad
 H|_{u^\perp}\succeq WI/4,\quad
 u^{\mathsf T}Hu\leq2q,\quad \|H\|\leq2W.
\]
Positivity of the block matrix gives
$\|PHu\|^2\leq(u^{\mathsf T}Hu)\|H|_{u^\perp}\|\leq4qW$.
Consequently, writing $x=x_\parallel u+x_\perp$ and applying Young's
inequality to the mixed term,
\[
 x^{\mathsf T}Hx
 \geq\frac W8\|x_\perp\|^2-32q\,x_\parallel^2.
\]
The same holds for $y$.  Since the projected hubs have norm at most one,
the exact fan-in lower bound therefore implies
\[
 \begin{split}
 2W-F_S(x,y)
 &\leq2|R_\parallel|+\frac{16}{W}\|R_\perp\|^2+16q\\
 &\leq66B+16q
 <66B+2048BL
 <2200BL.
 \end{split}
\]
This proves the geometric lemma in all cases.  The mixed-block estimate
is essential: a lower bound on the restriction of $Q_S$ to $u^\perp$
alone would not justify the conclusion for arbitrary hubs.

Strict improvement after the selector debit now requires
$akC<2200B\log(16d)$.  The bridge obstruction and EX-16 still apply
without the opposite-hub restriction.  For $k\geq43$ they yield
\[
 h\leq O\!\left(B\min\{\sqrt k,C\}\right)
 \leq O_a\!\left(
 B\min\left\{\sqrt k,\frac{B\log(16d)}k\right\}\right)
 \leq O_a\!\left(B^{4/3}(\log(16d))^{1/3}\right).
\]
For $k<43$ the bridge obstruction gives $h=O(B)$.

It is here that positive integer widths enter.  They give
$d\leq |E(G)|\leq W\leq kB/2$, even when the graph has parallel edges
with different endpoint directions.  Every binary certificate has
$C\geq2$, because its mean difference across any antipodal pair is a
vector of norm $2$.  Thus the selector condition implies
\[
 k\leq O_a(B\log(16kB)).
\]
For sufficiently large $B$ this forces $k<B^2$: otherwise the increasing
function $x/\log(16xB)$, evaluated at $x=B^2$, already exceeds every
fixed multiple of $B$.  Substitution then gives $k=O_a(B\log B)$ and
$\log(16d)=O_a(\log B)$.  This completes the upper bound.

Finally, for every fixed $\varepsilon>0$ the upper bound is
$o(B^{4/3+\varepsilon})$, while EX-18 attains
$h=\Theta(B^{4/3})$.  The supremum of fixed power exponents is therefore
$4/3$ in this model. \hfill$\square$

The selector assumptions remain substantive.  EX-19 does not lower-bound
the actual spectral unit-tensor exponent by a binary-seed cost.  Its
conclusion can be bypassed by a stronger spectral selector estimate,
more than two hubs, a circuit decomposition that avoids a linear
selector budget over balanced cuts, or target layouts beyond antipodal
equal-radius pairs.  For arbitrary positive real widths tending to zero,
the stated proof retains $\log(16d)$ instead of $\log B$ unless another
bound on the number of endpoint directions is supplied.

\paragraph{EX-20: bounded-size seed gadgets do not change the exponent.}
Replace binary seeds by seeds with at most $q$ image points, where
$q\geq2$ is fixed. Suppose the certificate still obtains the endpoint
positions by pointwise averaging, and charges a seed of image $A_\omega$
a number $c_\omega\geq\operatorname{diam}A_\omega$. Put
$K=\mathbb E c_\omega$ and charge the selector at least $akK$.
The exponent supremum in EX-19 remains $4/3$.

\emph{Proof.}
For a particular seed map $X$ with image
$\{a_1,\ldots,a_r\}$, $2\leq r\leq q$, form $r-1$ binary-image maps
\[
 Y_j(p)=a_1+(r-1)\mathbf 1_{\{X(p)=a_j\}}(a_j-a_1),
 \qquad 2\leq j\leq r.
\]
Their uniform mean is $X$, and their average diameter is
\[
 \sum_{j=2}^r\|a_j-a_1\|
 \leq(r-1)\operatorname{diam}X(P).
\]
A one-point seed is already binary. Replacing each seed in the original
distribution this way produces a binary certificate of cost
$C\leq(q-1)K$. Thus the original selector debit is at least
$a kC/(q-1)$, and EX-19 applies with the fixed constant $a/(q-1)$.
The lower bound on seed cost is automatic for an actual scaled
unit-tensor exponent, since every wire in the seed is a restriction of
that unit tensor. This argument only concerns estimates propagated by
pointwise averaging; it does not exhaust possible spectral constructions.
\hfill$\square$

Consequently, stronger estimates for three-party or any other fixed-size
seed cannot by themselves improve the power exponent through averaging.
An improvement along this route needs seed sizes growing with the
configuration, or a mechanism beyond this form of averaging.

The same argument gives a quantitative statement when $q=q(B)$ grows.
Keep $a>0$ fixed, independently of $q$. Since the mean endpoint map
contains antipodes at distance two, Jensen's inequality gives $K\geq2$.
The free-saving bound in EX-19 therefore implies
$kK=O_a(B\log(16d))$ and hence, exactly as there,
$\log(16d)=O_a(\log(2+B))$, with constants independent of $q$.
The binary refinement and EX-16 now give
\[
 h\leq O_a\!\left(B\min\left\{\sqrt{k},
       \frac{(q-1)B\log(2+B)}{k}\right\}\right)
 \leq O_a\!\left((q-1)^{1/3}B^{4/3}
                       (\log(2+B))^{1/3}\right).
\]
Thus even seeds of size $q=B^{o(1)}$ cannot improve the fixed power
exponent. A claimed bridge width $h\geq cB^{4/3+\varepsilon}$ in this
model would require
$q-1\geq\Omega_{a,c}(B^{3\varepsilon}/\log(2+B))$.

\subsection{Remaining gate for Q3}
\label{sec:ex-open}

\paragraph{EX-10: an improved selector bound would improve the exponent.}
Keep Anthropic's complete-network endpoint placement and two-hub geometry.
Suppose that, for all sufficiently large $k$, the simultaneous flat
realizations of all required endpoint and hub distances can be chosen with
\[
 c\,g_\varphi(P_k)=O(k^\beta),\qquad\beta\geq0.
\]
Then the same comparison proves a lower-bound exponent
\[
 \alpha=1+\frac{1}{2(1+\beta)}.
\]
Indeed the total hub saving is a constant multiple of $\ell k$, whereas the
selector cost is at most a constant multiple of $k^{1+\beta}$ and the bridge
length is $\Theta(k^{-1/2})$.  Taking $\ell$ a sufficiently large constant
multiple of $k^\beta$, and $B=k\ell$, gives
$\log L=\Omega(B\sqrt k)=\Omega(B^{1+1/(2(1+\beta))})$.
Rounding $\ell$ and choosing a suitable integer subsequence of $k$ changes
only constants.  Metric approximation errors can be made arbitrarily smaller
than these margins because each configuration is finite.  The known bound
$\beta=1/2$ gives $4/3$; a uniform bound $\beta=0$ would give $3/2$.

Since a wire between two endpoints is a restriction of the multipartite unit
tensor, $g_\varphi(P_k)\geq d_\varphi(p,q)$.  Opposite endpoints have Euclidean
distance $1$, so any such approximate realization has
$c g_\varphi(P_k)\geq1-o(1)$.  This rules out a negative power $\beta$ and limits
this particular parameter balance to $3/2$, without proving that $3/2$ can be
attained.

\paragraph{EX-11: a constant-mass cut representation of the endpoint metric.}
Let $p=2-\sqrt2$ and $M=p^{-1}=1+1/\sqrt2$.  On
$P_k\cup\{0\}$ form a random cut as follows.  The center lies on side zero.
Independently for each antipodal pair, activate the pair with probability $p$.
For an active pair put one uniformly chosen endpoint on side one, leaving its
partner on side zero.  Inactive pairs have both endpoints on side zero.  Let
$\delta_C$ be this cut's zero-one pseudometric.  Then
\[
  \|x-y\|=M\,\mathbb E\delta_C(x,y)
  \qquad(x,y\in P_k\cup\{0\}).
\]
For a center-endpoint pair the right-hand side is $Mp/2=1/2$; for antipodes it
is $Mp=1$; and for endpoints from distinct pairs it is
$M(p-p^2/2)=1-p/2=1/\sqrt2$.  These are exactly the required distances.
The sample space is finite, so this is a finite cut decomposition with total
coefficient mass $M$ independent of $k$.  Empty cuts can be discarded, making
the total mass no larger.

\paragraph{EX-12: ordinary graph-rank inequalities cannot supply a growing
lower bound on the selector exponent.}
Consider any asymptotic restriction
\[
 T_a\preceq \langle q^\rho\rangle_{P_k}\otimes T_b,
\]
where $T_a,T_b$ are products of wires on $P_k\cup\{0\}$, with respective
nonnegative logarithmic loads $a_{xy},b_{xy}$, and $\rho\geq0$.  For every cut
in the support of EX-11 that is nonempty, the cut splits $P_k$.  Flattening ranks
therefore require
\[
 \rho\geq\sum_{xy}(a_{xy}-b_{xy})\delta_C(x,y).
\]
Using EX-11 gives
\[
 \sum_{xy}(a_{xy}-b_{xy})\|x-y\|\leq M\rho.
\]
On a Euclidean price realization, the same restriction would lower-bound the
scaled selector exponent $G=cg_\varphi(P_k)$ only by
\[
 G\geq\rho^{-1}\sum_{xy}(a_{xy}-b_{xy})\|x-y\|
 \quad(\rho>0),
\]
whose right-hand side is at most $M$.  Thus no member of this entire class of
wire-product restrictions, including graph-tensor rank upper bounds, can by
itself force $G=\Omega(\sqrt k)$.  If $\rho=0$, its Euclidean weighted length is
nonpositive and yields no positive lower bound on $G$.  This argument does
not say that a cut mixture is a price: products of powers of flattening ranks
need not satisfy the cover inequality.

\paragraph{EX-21: a failed route from cut mixtures to actual selector bounds.}
EX-11 also shows that the endpoint metric, divided by a constant at most
$1+1/\sqrt2$, lies in the convex hull of normalized flattening wire
profiles. Indeed, remove the empty cut and renormalize the remaining
coefficients; their total mass is
\[
 M_k=(1+1/\sqrt2)\bigl(1-(\sqrt2-1)^k\bigr)\leq1+1/\sqrt2.
\]
Every remaining cut splits the endpoint set, so each of its flattening
ranks is an ordinary normalized spectral point on that set. Consequently
no linear inequality in the normalized endpoint wire exponents can force
a selector scale diverging with $k$. This includes arbitrary such valid
linear inequalities, not merely a preselected list of graph-rank bounds.
It says nothing about nonlinear constraints on the joint wire spectrum.

A tempting next step is to take a geometric mean of the cut flattening
ranks and invoke a lower-functional extension theorem. The following
four-party example shows the missing hypothesis explicitly. Let
$A=12|34$, $B=14|23$, and
\[
 f(T)=\sqrt{\operatorname{rank}_A(T)\operatorname{rank}_B(T)}.
\]
For the binary wires $T=\langle2\rangle_{12}$ and
$S=\langle2\rangle_{14}$ one has
\[
 f(T)=f(S)=\sqrt2,\qquad f(T\oplus S)=3.
\]
Thus $f$ already fails the cover inequality. It is nevertheless
multiplicative, homogeneous on natural multiples, and superadditive,
so one might try Alman, Li, and Pratt (2026), Theorem~3.9, p.~21.
That theorem has an additional monomial-dominance requirement, and
this example violates it by a fixed exponential factor:
\[
 \binom nk f(T^kS^{n-k})
 =\binom nk\,2^{n/2}
 \leq(2\sqrt2)^n
 =\left(\frac{2\sqrt2}{3}\right)^n f(T\oplus S)^n
 \qquad(0\leq k\leq n).
\]
The loss cannot be $2^{-o(n)}$. Therefore that theorem does not turn
the constant-mass cut certificate into a dominating spectral point.
The rank computations and this obstruction hold in every
characteristic; the natural multiplicities here are tensor direct-sum
multiplicities, not scalar coefficients reduced modulo the characteristic.

The exact supremum requested in Q3 remains undecided if the hub-displacement
method permits arbitrary selectors, multiple hubs, or a different circuit
decomposition. EX-19 proves that the exponent supremum is exactly $4/3$
across arbitrary graphs with integer bit widths, arbitrary antipodal unit
endpoint configurations, independently and globally optimized two hubs,
and the stated binary-seed selector accounting. Its upper bound has a
$(\log B)^{1/3}$ factor; EX-18 removes that factor for opposite hubs.
The exact $5/4$ barrier in the orthogonal model and the exact $\sqrt k$
binary-seed cost are complementary restricted statements. A matching
unrestricted upper barrier would require a definition of the admissible
method class and a theorem covering constructions outside these models.
No such unrestricted theorem is established here.

\subsection{A multiscale extension to arbitrary radial wire lengths}
\label{sec:ex-multiscale}

\paragraph{EX-23: the $4/3$ supremum for a common-origin two-hub model.}
The equal-radius restriction in EX-19 can be removed. Let a graph have
positive integer bit widths $b_e$, maximum incident load $B$, and $k\geq2$
vertices. For each edge choose a unit vector $x_e$ and place its endpoints
at $r_e^+x_e$ and $-r_e^-x_e$, where $r_e^+,r_e^-\geq0$ and
$r_e^++r_e^->0$. Thus every target wire segment passes through one common
origin; the two end radii need not agree. Its target cost is
$\sum_e b_e(r_e^++r_e^-)$. Suppose the finite endpoint set has a binary
pointwise-averaging certificate of cost $C$, and the method charges at
least $akC$ for every balanced cut, with $a>0$ fixed. Allow the two hubs
to be placed independently anywhere in Euclidean space. If every
balanced cut admits strict improvement after this selector debit and a
bridge of width $h$, then
\[
 h=O_a\bigl(B^{4/3}\log(2+B)\bigr).
\]
The exponent supremum in this larger model is exactly $4/3$, under the
same flat-case reduction used for the attainment in EX-18, since that
equal-radius construction is included. The conclusion still
concerns the specified binary-certificate debit, not a lower bound on
the actual spectral unit-tensor exponent.

\emph{A simultaneous matrix-balancing device.}
We first give the additional ingredient. Let $H_1,\ldots,H_k$ be real
symmetric matrices, regarded as vectors in the real Hilbert space of
symmetric matrices with its Hilbert--Schmidt inner product. Put
\[
 V=\sum_a \operatorname{vec}(H_a)\operatorname{vec}(H_a)^{\mathsf T}.
\]
On the range of $V$, each vector
$V^{-1/2}\operatorname{vec}(H_a)$ has norm at most one, since
$\operatorname{vec}(H_a)\operatorname{vec}(H_a)^{\mathsf T}\preceq V$.
If $g$ is a standard Gaussian on this range, the matrix represented by
$V^{1/2}g$ has the same law as $\sum_a g_aH_a$, where the $g_a$ are
independent standard Gaussians. Consequently the inverse image under
$V^{1/2}$ of an operator-norm ball is a symmetric convex body whose
Gaussian measure is controlled by Tropp's Gaussian matrix-series
bound. This allows a matrix constraint to be one block in a single
Banaszczyk application, together with scalar and Euclidean-norm
constraints. No assertion about a randomly chosen final signing is
needed.

\emph{The simultaneous cut lemma.}
Partition any retained edges into $m\geq1$ bands, with a band parameter
$r_j>0$ such that
\[
 r_j/2<\max(r_e^+,r_e^-)\leq r_j
 \qquad(e\hbox{ in band }j).
\]
Let $D\geq1$ be the dimension spanned by all directions and set
$\ell=\ln(32mD)$. Write, for each band,
\[
 W_j=\sum_{e\in j}b_e,\quad
 A_j=\sum_{e\in j}b_ex_ex_e^{\mathsf T},\quad Q_j=W_jI-A_j,
 \quad
 Q_{aj}=\sum_{e\ni a,\ e\in j}b_e(I-x_ex_e^{\mathsf T}).
\]
Thus $0\preceq Q_{aj}\preceq BI$ and
$\sum_aQ_{aj}=2Q_j$. Put
$v_{aj}=\sum_{e\ni a,\ e\in j}\pm b_ex_e$, with signs fixed by the
edge endpoints, so $\sum_av_{aj}=0$. Choose a top eigenvector $u_j$ of
$A_j$ and write $v_{aj}=\alpha_{aj}u_j+w_{aj}$. Empty bands contribute
zero blocks. There is a balanced cut $S$, provided $k\geq150m$, for
which all bands simultaneously satisfy
\begin{align}
 \|R_j\|&\leq L\,BC/r_j, & L&=50m,\label{eq:ex-band-resultant}\\
 \sup_{x,y}\{\text{target cost in band }j-
                 \text{fan-in cost in band }j\}
       &\leq KBr_j, & K&=60000(m^2+m\ell),
       \label{eq:ex-band-saving}
\end{align}
where $R_j=\sum_{a\in S}v_{aj}$.

Here are full details of the common signing. The binary certificate,
as in EX-16, gives a zonotope $Z=\sum_\omega[-g_\omega,g_\omega]$
containing every half-difference of endpoint positions, and an
ellipsoid matrix $Q_0$ with
\[
 Z\subseteq Q_0^{1/2}\mathbb B_2,
 \qquad \operatorname{tr}Q_0\leq(C/2)^2.
\]
Since $(r_e^++r_e^-)/2>r_j/4$ in band $j$, it follows that
$x_e\in(4/r_j)Z$ and
$\|(r_j/(4B))Q_0^{-1/2}v_{aj}\|\leq1$.
Other norm-at-most-one blocks are
\[
 \alpha_{aj}/B,\qquad
 B^{-1/2}(A_j|_{u_j^\perp})^{-1/2}w_{aj},
\]
where the second claim follows from
$w_{aj}w_{aj}^{\mathsf T}\preceq B A_j|_{u_j^\perp}$.
Use two further matrix-whitening blocks per band, applying the
preceding device to $P_jQ_{aj}P_j$ on $u_j^\perp$ and, when $Q_j$
is invertible, to $Q_j^{-1/2}Q_{aj}Q_j^{-1/2}$. Include the scalar
count block $1$. There are at most $5m+1$ blocks, so division by
$\Lambda=5\sqrt{5m+1}$ makes each whole input vector have norm at
most $1/5$.
All selector and transverse pseudoinverses act only on the ranges of
their positive semidefinite matrices; zero-dimensional blocks are
omitted. The same convention applies to covariance whitening.

For each selector block impose
$\|Q_0^{1/2}z\|\leq4\sqrt m\,(C/2)$; for each principal and the
count block impose $|t|\leq4\sqrt m$; for each transverse block
impose $\|(A_j|_{u_j^\perp})^{1/2}z\|\leq4\sqrt{mW_j}$.
Markov's inequality gives Gaussian failure probability at most
$1/(16m)$ for each such block. For the perpendicular matrix block,
\[
 \sum_a(P_jQ_{aj}P_j)^2\preceq2BW_j I.
\]
For the full whitened block, writing $q_j=\lambda_{\min}(Q_j)>0$,
the analogous sum is at most $2B/q_j$ times the identity.
Tropp's Gaussian bound therefore permits respective body thresholds
$2\sqrt{BW_j\ell}$ and $2\sqrt{(B/q_j)\ell}$, each with failure
probability at most $1/(16m)$. The product body has Gaussian measure
at least $1-(5m+1)/(16m)\geq5/8$. Banaszczyk's theorem now gives a
single signing $\sigma_a\in\{-1,1\}$ obeying every constraint.

Take $S=\{a:\sigma_a=1\}$. The count imbalance is less than $50m$,
so $k\geq150m$ ensures balance. Since
$R_j=\tfrac12\sum_a\sigma_av_{aj}$, the other bounds give
\[
 \|R_j\|\leq50mBC/r_j,\qquad
 |(R_j)_{u_j}|\leq25mB,\qquad
 \|(R_j)_\perp\|^2\leq600m^2BW_j.
\]
When $W_j\geq3000mB\ell$, the perpendicular matrix signing has
operator norm at most $W_j/2$, hence
$Q_{S,j}|_{u_j^\perp},Q_{S^c,j}|_{u_j^\perp}\succeq W_jI/4$.
When $q_j\geq3000mB\ell$, full whitening additionally gives
$Q_{S,j},Q_{S^c,j}\succeq Q_j/2$.

To check the free-saving bound, project either hub to the ball of
radius $r_j$ when considering this band alone; this decreases every
fan-in distance. For $0\leq r\leq r_j$ and $\|z\|\leq r_j$,
\[
 \|rx-z\|\geq r-x\cdot z+
       \frac{\|z\|^2-(x\cdot z)^2}{4r_j}.
\]
This follows from the same denominator identity as EX-17; its
denominator is at most $2(r+\|z\|)\leq4r_j$. The zero-denominator
case follows directly. Scaling the hubs by $r_j$ reduces the saving
calculation to EX-19. If $W_j<3000mB\ell$, the entire target cost
is at most $2W_jr_j$. Otherwise, if
$q_j\geq3000mB\ell$, the full-whitening case bounds saving divided
by $r_j$ by
$2|(R_j)_{u_j}|+8\|(R_j)_\perp\|^2/W_j$.
In the remaining case the mixed-block estimate of EX-19 gives
\[
 2|(R_j)_{u_j}|+16\|(R_j)_\perp\|^2/W_j+16q_j.
\]
All three bounds are at most $KB$. If the directional span of a band
is one-dimensional, the transverse terms vanish and the principal
bound alone suffices. This proves the simultaneous cut lemma.

\emph{Truncation and optimization.}
Let $r_{\max}=\max_e\max(r_e^+,r_e^-)$ and
$\delta=aC/(4B)$. Because $C$ is at least every endpoint distance,
$r_{\max}\leq C$. Edges with maximum radius below $\delta$ have
total target cost at most
\[
 2\delta\sum_e b_e\leq\delta kB=akC/4.
\]
They cannot save more than this amount. If $r_{\max}\leq\delta$,
strict improvement is already impossible. Otherwise place the
remaining edges in dyadic bands with
$r_j=2^{-j}r_{\max}$. The number of bands obeys
\[
 m\leq1+\log_2(4B/a)=O_a(\log(2+B)).
\]
For hub separation $d=\|x-y\|$, the supporting norm inequality
at the common origin also bounds band $j$'s saving by
$R_j\cdot(x-y)\leq\|R_j\|d$. The simultaneous cut therefore has
retained-edge saving at most
\[
 \sum_j\min\{KBr_j,LBCd/r_j\}
 \leq4B\sqrt{KLCd}.
\]
For completeness, split the dyadic sequence at
$r_*=(LCd/K)^{1/2}$. The sum of radii below $r_*$ is at most
$2r_*$, and the sum of inverse radii above $r_*$ is at most
$2/r_*$. This proves the displayed bound, including when one part
is empty. Strict improvement after the bridge and selector thus
requires
\[
 \frac34akC<4B\sqrt{KLCd}-hd
 \leq\frac{4B^2KLC}{h},
 \qquad
 h<\frac{16B^2KL}{3ak}.
\]
Independently, the uniformly random balanced-cut calculation in
EX-16 and the supporting norm inequality at the origin give
$h\leq O(B\sqrt k)$, for arbitrary radii. If $k<150m$, this
already suffices for the claimed final bound. Otherwise combine
the last two bounds to obtain
\[
 h\leq O_a\bigl(B^{4/3}(KL)^{1/3}\bigr).
\]

It remains to control $K$ without a hidden dimension dependence.
The simultaneous free-saving bounds, with the bridge cost omitted,
give $\tfrac34akC<\sum_jKBr_j\leq2KBC$, hence
$k=O_a(BK)$. Integer widths give
$D\leq|E|\leq\sum_e b_e\leq kB/2$. Therefore
\[
 k\leq O_a\bigl(B\log(2+B)
              (\log(2+B)+\log(2+k))\bigr).
\]
For sufficiently large $B$, this implies $k<B^2$ by the increasing
function $x/(\log B+\log x)$, and then
$k=O_a(B(\log B)^2)$. Thus
$\ell=O_a(\log B)$, $K=O_a((\log B)^2)$, and
$KL=O_a((\log B)^3)$. Substitution proves
$h=O_a(B^{4/3}\log(2+B))$. The bounded range of $B$ is absorbed
in the constant. \hfill$\square$

The growing-seed conclusion of EX-20 extends to this radial model,
with the explicit bound
\[
 h=O_a\bigl((q-1)^{1/3}B^{4/3}\log(2+B)\bigr).
\]
To keep the constants uniform in $q$, write $K_0$ for the original
average seed debit, so $K_0\geq r_{\max}$ and the selector debit is
$akK_0$. Binary refinement has cost $C'\leq(q-1)K_0$.
In the proof use $\delta=aK_0/(4B)$, so the number of radial bands
is still $O_a(\log B)$ independently of $q$. The free-saving and
dimension-bootstrap estimates are unchanged, whereas the resultant
bound becomes $LB(q-1)K_0/r_j$. Accordingly the optimized bound is
$h\leq O_a(B^2KL(q-1)/k)$, which gives the displayed conclusion
after combination with $h\leq O(B\sqrt k)$.
Subpolynomially growing seed sizes therefore still have fixed power
supremum $4/3$ in this larger model. To obtain
$h\geq cB^{4/3+\varepsilon}$ by this route requires
$q-1\geq\Omega_{a,c}(B^{3\varepsilon}/(\log B)^3)$.

\paragraph{EX-24: a sharper dependence on the seed size.}
The linear binary-refinement cost used in EX-20 admits a square-root
improvement. For a seed image $\{a_1,\ldots,a_r\}$, set
$\bar a=r^{-1}\sum_i a_i$ and $b_i=a_i-\bar a$. Choose independent
Rademacher signs $\varepsilon_i$, put
$v_\varepsilon=\sum_i\varepsilon_i b_i$, and replace the seed map
by the binary-image map
\[
 Y_\varepsilon(p)=\bar a+
        \varepsilon_{i(p)}v_\varepsilon,
 \qquad X(p)=a_{i(p)}.
\]
Its expectation is $X(p)$. Moreover,
\begin{align*}
 \mathbb E\operatorname{diam}Y_\varepsilon(P)
 &\leq2\mathbb E\|v_\varepsilon\|
 \leq2\sqrt{\sum_i\|b_i\|^2}\\
 &=2\sqrt{\frac1r\sum_{i<j}\|a_i-a_j\|^2}
 \leq\sqrt{2(r-1)}\operatorname{diam}X(P).
\end{align*}
A constant seed requires no cost. Thus for images of size at most
$q\geq2$ one may replace the factor $q-1$ throughout the preceding
refinements by
\[
 \rho_q=\min\{q-1,\sqrt{2(q-1)}\}.
\]
In particular, with constants independent of $q$,
\[
 \begin{array}{ll}
 \text{equal-radius model:}&
 h=O_a\bigl(q^{1/6}B^{4/3}(\log(2+B))^{1/3}\bigr),\\[2pt]
 \text{common-origin radial model:}&
 h=O_a\bigl(q^{1/6}B^{4/3}\log(2+B)\bigr).
 \end{array}
\]
Hence a power improvement by $\varepsilon>0$ in the radial model
requires $q=\Omega_{a,c}(B^{6\varepsilon}/(\log B)^6)$.
The refinement distribution has finite support; its cardinality
does not enter the geometric certificate cost or the barrier proof.

\paragraph{EX-26: an arbitrary number of star-connected hubs.}
Return to unit antipodal endpoints and positive integer cable widths.
Let $r=2^t\ge2$ divide $k$, and put $s=k/r$. For every partition
of the circuit vertices into $r$ groups of size $s$, permit one
independently positioned hub per group and connect those hubs by
a star, with every star edge of width $h$. Assume the selector
debit is at least $ak(\log_2r)C$, with the same binary certificate
cost $C$ and fixed $a>0$ for all partitions. If every such partition
admits a strict improvement over the baseline $2W$, then, uniformly
in $r$,
\[
 h=O_a\!\left(B^{4/3}
       \left(1+\frac{\log(2+B)}{\log r}\right)^{1/3}\right).
 \tag{EX26.1}
\]
In particular the fixed power exponent cannot exceed $4/3$ even
when the number of star hubs grows. This statement specifies an
additional certificate model; it does not assert a new circuit
decomposition theorem or a bound for general bridge trees.

\emph{Proof.}
For a group $S_j$ put $R_j=\sum_{a\in S_j}v_a$.
Write $z_0$ for the center hub. Convexity at the origin gives
\[
 \text{total cost}-2W
 \ge -\sum_jR_j\cdot z_j
          +h\sum_{j\ne0}\|z_j-z_0\|
 =\sum_{j\ne0}\bigl(h\|z_j-z_0\|
                    -R_j\cdot(z_j-z_0)\bigr),
 \tag{EX26.2}
\]
before the positive selector debit, since $\sum_jR_j=0$.
Consequently a strict improvement requires
$h<\max_{j\ne0}\|R_j\|$.

There are two useful equal-size partitions. First minimize
$\sum_j\|R_j\|^2$ among all such partitions. Swapping
$a\in S_i$ and $b\in S_j$ cannot improve this objective, so
\[
 (R_i-R_j)\cdot(v_a-v_b)\le\|v_a-v_b\|^2.
\]
Summing over the $s^2$ choices gives
$s\|R_i-R_j\|^2\le4s^2B^2$.
Averaging $R_i-R_j$ over $j$ then yields
$\|R_i\|\le2B\sqrt s$ for every group.

For the second partition use the certificate ellipsoid $Q_0$
from EX-16, with $\operatorname{tr}Q_0=(C/2)^2$ and
$\|Q_0^{-1/2}v_a\|\le B$. All inverses are on its range.
At each node of a recursive bisection, pair its vertices
arbitrarily and apply Banaszczyk's theorem to the difference
vectors
$Q_0^{-1/2}(v_a-v_b)/(10B)$.
Their norms are at most $1/5$. The body
$\{y:\|Q_0^{1/2}y\|\le C\}$ has Gaussian measure at least
$3/4$, by its second moment. Assign opposite signs within each
pair according to the resulting signing. The signed resultant
error has norm at most $10BC$. Each child has exactly half the
vertices and resultant $(R_{\rm parent}\pm E)/2$, with
$\|E\|\le10BC$. Starting with resultant zero, every descendant
therefore has resultant norm less than $10BC$. Because $r$ is
a power of two dividing $k$, this produces exactly $r$ groups
of size $s$. Equation (EX26.2), applied to the two possibly
different partitions, proves
\[
 h\le10B\min\{C,\sqrt s\}.
 \tag{EX26.3}
\]

We next construct a third partition whose total free hub saving
is at most $500rB\ell$, where $D$ is the span dimension of the
wire directions and $\ell=\log(32rD)>1$. Put
\[
 A=\sum_e b_ex_ex_e^{\mathsf T},\quad Q=WI-A,
 \quad Q_a=\sum_{e\ni a}b_e(I-x_ex_e^{\mathsf T}).
\]
Then $0\preceq Q_a\preceq BI$, $\sum_aQ_a=2Q$.
Let $u$ be a top eigenvector of $A$, let $P$ project onto
$u^\perp$, and put $q=u^{\mathsf T}Qu$.
As before $Q_\perp\succeq(W/2)I$ and $Q_\perp\preceq WI$.
Sort the numbers $\alpha_a=u\cdot v_a\in[-B,B]$ and form
consecutive blocks of $r$ vertices. Independently within each
block, assign its vertices to the $r$ groups by a uniform
permutation. Every group has size $s$. Since the sum of the
ranges of these consecutive blocks is at most $2B$, while
$\sum_a\alpha_a=0$, every outcome satisfies
\[
 |u\cdot R_j|\le2B\qquad\text{for every }j.
 \tag{EX26.4}
\]
Independence between blocks and the weighted Cauchy--Schwarz
bound $\sum_a\|Pv_a\|^2\le2BW$ give
\[
 \mathbb E\sum_j\|PR_j\|^2\le2BW.
\]
Thus with probability at least $3/4$ their total is at most $8BW$.

Write $H_j=\sum_{a\in S_j}Q_a$, so $\mathbb EH_j=2Q/r$.
For a fixed group the centered matrix contributions of the
blocks are independent, have norm at most $B$, and their
total variance on $u^\perp$ is at most $2BW/r$.
The two-sided matrix Bernstein inequality of Tropp (2012),
Theorem~6.1(ii), pp.~23--24 (also Theorem~1.4, p.~4), gives
\[
 \Pr\!\left\{\|PH_jP-2Q_\perp/r\|>W/(2r)\right\}
 \le2D\exp\{-W/(18Br)\}.
 \tag{EX26.5}
\]
Indeed the variance assertion follows from
$\mathbb E(X-\mathbb EX)^2\preceq\mathbb EX^2
\preceq B\mathbb EX$ for each positive matrix contribution.
If $q>0$, applying the same argument after conjugation by
$Q^{-1/2}$ gives
\[
 \Pr\!\left\{\|Q^{-1/2}H_jQ^{-1/2}-2I/r\|>1/r\right\}
 \le2D\exp\{-q/(5Br)\}.
 \tag{EX26.6}
\]
These bounds use only independence between blocks; assignments
to distinct groups need not be independent.

If $W<64rB\ell$, the entire baseline $2W$ is already less than
$128rB\ell$. Otherwise (EX26.5), a union bound over the $r$
groups, and the preceding Markov bound have a common outcome
for which
\[
 \frac{W}{2r}I\preceq PH_jP\preceq\frac{3W}{r}I,
 \qquad \sum_j\|PR_j\|^2\le8BW.
 \tag{EX26.7}
\]
When $q\ge64rB\ell$, use (EX26.6) in place of (EX26.5)
to get additionally, in a common outcome, $H_j\succeq Q/r$.
The failure union in either use is less than $1/4$.

For a free hub one may project its position into the unit ball.
The exact distance inequality in EX-17 bounds the free saving
of group $j$ by
\[
 \sup_{\|z\|\le1}
       \{R_j\cdot z-\tfrac14z^{\mathsf T}H_jz\}.
 \tag{EX26.8}
\]
In the regime $q\ge64rB\ell$, the unrestricted quadratic
maximum and (EX26.4) give a total saving at most
\[
 r\sum_jR_j^{\mathsf T}Q^{-1}R_j
 \le\frac{4r^2B^2}{q}
          +\frac{2r}{W}\sum_j\|PR_j\|^2
 <17rB.
\]
In the remaining regime $q<64rB\ell$, put
$a_j=u^{\mathsf T}H_ju$ and $b_j=PH_ju$.
Positive semidefiniteness and (EX26.7) imply
$\|b_j\|^2\le(3W/r)a_j$. Writing $z=tu+w$, with
$w\perp u$ and $|t|\le1$, and completing a square in the
mixed term gives
\[
 z^{\mathsf T}H_jz
 \ge\frac{W}{4r}\|w\|^2-12a_j.
\]
Therefore (EX26.8) is at most
$|u\cdot R_j|+4r\|PR_j\|^2/W+3a_j$.
Since $\sum_ja_j=2q$, the total is at most
\[
 2rB+32rB+6q<418rB\ell.
\]
If $D=1$, the perpendicular spaces are absent and the direct
bound $\sum_j|R_j|\le2rB$ suffices. This proves the claimed
$500rB\ell$ free-saving bound in every case.

The assumed feasibility for this third partition now gives
\[
 sC\log_2r<\frac{500}{a}B\ell.
 \tag{EX26.9}
\]
For $z=\min\{C,\sqrt s\}$, one has $z^3\le sC$;
(EX26.3) and (EX26.9) imply
$h=O_a(B^{4/3}(\ell/\log r)^{1/3})$.
Finally integer widths give $D\le W\le rsB/2$, while
$C\ge2$ because the certificate preserves a unit antipodal
pair. Hence (EX26.9) implies
\[
 s\le O_a\bigl(B(1+\log s+\log(2+B))\bigr),
\]
so $s=O_a(B\log(2+B))$ by the same elementary bootstrap as
EX-19. In particular
$\ell/\log r=O_a(1+\log(2+B)/\log r)$, proving (EX26.1).
\hfill$\square$

For seeds of size at most $q$, charged by their original
average debit $K_0$, the same argument uses $K_0$ in (EX26.9)
and $\rho_qK_0$ in the certificate resultant bound.
It adds a factor $\rho_q^{1/3}=O(q^{1/6})$ to (EX26.1).
The logarithmic bootstrap remains independent of $q$ since
$K_0\ge2$. Arbitrary endpoint radii and general bridge trees
are not included in EX-26.

\paragraph{EX-27: general bridge trees with uniform angular curvature.}
Retain the equal-size partition, unit-antipodal-endpoint, and selector
assumptions of EX-26, but let the bridges be any fixed tree on the
$r=2^t$ labelled hubs. Put $D=\dim\operatorname{span}\{x_e\}$ and
\[
 \ell=\log\bigl(64r^2(D+1)\bigr),\qquad Q=WI-A.
\]
Assume, for a fixed $0<\eta\le1$, that
\[
 Q\succeq\eta WI,\qquad \frac Wr\ge64\eta^{-2}B\ell.
 \tag{EX27.1}
\]
The tree and its hub labels are fixed before the partition is chosen:
group $j$ is assigned to hub $j$. If every labelled equal-size
partition admits a strict improvement including
the debit $ak(\log_2r)C$, then
\[
 h=O_{a,\eta}\bigl(B^{4/3}\ell^{2/3}\bigr)
   =O_{a,\eta}\bigl(B^{4/3}(\log(2+rB))^{2/3}\bigr).
 \tag{EX27.2}
\]
Thus for any fixed $M$, allowing arbitrary tree shapes and at most
$B^M$ hubs cannot improve the fixed power exponent beyond $4/3$
under (EX27.1). The additional angular-curvature and average-width
conditions are substantive assumptions, not conclusions of the
general two-hub or star-hub theorems.

\emph{Proof.}
Root the hub tree and order its vertices by depth-first preorder.
Every rooted subtree is an interval in this ordering. The paired
Banaszczyk bisections from EX-26 assign the circuit vertices to
consecutive dyadic blocks of hub labels. Every such block has
resultant norm at most $10BC$. Each interval is a disjoint union
of at most $2\log_2r$ dyadic blocks. Hence the resultant on either
side of every tree edge has norm at most $20BC\log_2r$.
Summation by parts on the tree, followed by convexity of each
fan-in at zero, shows that this partition cannot improve cost
when $h\ge20BC\log_2r$. Therefore
\[
 h\le20BC\log_2r.
 \tag{EX27.3}
\]

We construct a second equal-size partition with simultaneous
curvature and interval bounds. Divide the $k$ circuit vertices
arbitrarily into blocks of $r$, and independently assign each
block to the hub labels by a uniform permutation. Every group
has exactly $s=k/r$ vertices. As in EX-26, with
$H_j=\sum_{a\in S_j}Q_a$, matrix Bernstein and (EX27.1) give,
with failure probability less than $1/16$ after the union over
all hubs,
\[
 H_j\succeq\frac{\eta W}{r}I\qquad(j=1,\ldots,r).
 \tag{EX27.4}
\]
Indeed the mean is $2Q/r$, the centered block norm is at most
$B$, the variance is at most $2BW/r$, and the threshold
$\eta W/r$ has failure bound
$2D\exp\{-\eta^2W/(5Br)\}$ for a single hub.

Here is a careful sampling-without-replacement argument for the
needed interval bounds. Fix a label subset $I$ of size $u\le r/2$.
Within one circuit-vertex block, reveal the $u$ vectors assigned
to these labels. The Doob martingale for their final sum has
increment at draw $j$ equal to
\[
 \frac{r-u}{r-j}
       (v_{\rm draw}-\overline v_{\rm remaining}),
 \qquad 1\le j\le u,
 \tag{EX27.5}
\]
where the remaining mean is taken before that draw. Its norm
is at most $2B$. Its conditional squared norm is at most
$(2/r)\sum_{a\text{ in block}}\|v_a\|^2$, since at least
$r/2$ vectors remain and the displayed prefactor is at most
one. Concatenating the block martingales gives final value
$R_I=\sum_{j\in I}R_j$, because $\sum_av_a=0$.
The predictable quadratic variation after self-adjoint vector
dilation has operator norm at most
\[
 \frac{2u}{r}\sum_a\|v_a\|^2\le\frac{4BWu}{r}.
 \tag{EX27.6}
\]
Both diagonal blocks of the dilation variance are bounded by
the scalar sum of the conditional squared vector norms, so
no dimension factor enters this variance estimate.

Tropp (2011), \emph{Freedman's inequality for matrix martingales},
\emph{Electronic Communications in Probability}, 16, 262--270,
Theorem~1.2 and Corollary~1.3, printed p.~264,
\href{https://doi.org/10.1214/ECP.v16-1624}
{doi:10.1214/ECP.v16-1624}, now imply
\[
 \Pr\{\|R_I\|>8\sqrt{BWu\ell/r}+4B\ell\}
 \le(D+1)e^{-\ell}.
 \tag{EX27.7}
\]
The theorem is applied with increment bound $2B$ and the
deterministic variance bound (EX27.6). For $u>r/2$, apply
the same argument to the complementary labels and use
$R_I=-R_{I^c}$; the weaker displayed bound with $u$ remains
valid. The full-label resultant is zero. There are at most
$r^2$ preorder intervals, so (EX27.4) and the bounds
\[
 \|R_I\|\le M(|I|),\qquad
 M(u)=8\sqrt{BWu\ell/r}+4B\ell
 \tag{EX27.8}
\]
hold simultaneously for some equal-size partition. This
argument uses exact equal sizes throughout, not independent
unconditioned labels.

Fix that partition. Projecting all hubs onto the unit ball
does not increase fan-in or bridge lengths, because Euclidean
metric projection is nonexpansive. The exact distance bound
from EX-17 and (EX27.4) therefore bound its maximum saving,
before the selector debit, by
\[
 F_h=\sup_{z_1,\ldots,z_r}
 \left\{\sum_jR_j\cdot z_j
       -\frac{\eta W}{4r}\sum_j\|z_j\|^2
       -h\sum_{ij\in E(T)}\|z_i-z_j\|\right\}.
 \tag{EX27.9}
\]
The unrestricted supremum only enlarges the upper bound.

Choose a dyadic length $L$ dividing $r$ and partition the
preorder list into consecutive blocks of $L$ hubs. Replace
each $R_j$ by the mean $\overline R_j$ of its block. The
residuals sum to zero on whole blocks. For any tree-edge
subtree interval, only its two partial end blocks can
contribute to its residual resultant. For a partial block
$J\subseteq K$, that resultant is
$R_J-(|J|/L)R_K$ and has norm at most $2M(L)$.
Consequently every residual tree-edge flow has norm at most
$4M(L)$. If $4M(L)\le h$, tree summation by parts absorbs
the residual linear term into the bridge term in (EX27.9),
giving
\[
 F_h\le\frac r{\eta W}\sum_j\|\overline R_j\|^2.
 \tag{EX27.10}
\]
For $L=1$ there is no residual, so this inequality holds
without a condition on $h$.

Put $v_0=BW\ell/r$. Assumption (EX27.1) gives
$B\ell\le\sqrt{v_0}/8$. Thus
\[
 4M(L)\le34\sqrt{v_0L},\qquad
 M(L)^2\le129v_0L\quad(L\ge1).
\]
The second inequality follows from
$(a+b)^2\le2a^2+2b^2$ and
$32B^2\ell^2\le v_0/2$.
Since each block contributes $\|R_K\|^2/L$ to the sum
in (EX27.10), this sum is at most $129BW\ell/L$.
If $h^2/(34^2v_0)<2$, use $L=1$.
If $2\le h^2/(34^2v_0)<r$, take the largest dyadic
$L\le h^2/(34^2v_0)$, which is at least half that number.
If $h^2/(34^2v_0)\ge r$, take the single block $L=r$;
its mean is zero, so $F_h=0$.
In every case, for $h>0$,
\[
 F_h\le\frac{2^{20}}{\eta}
             \frac{B^2W\ell^2}{h^2}.
 \tag{EX27.11}
\]
The case $h=0$ already satisfies the claimed conclusion.

Feasibility for this partition and $2W\le kB$ imply
\[
 h^2 C\log_2r
 <\frac{2^{19}}{a\eta}B^3\ell^2.
\]
Together with (EX27.3), this gives
$h^3\le(20\cdot2^{19}/(a\eta))B^4\ell^2$,
which proves the first bound in (EX27.2).
For the second, the free-saving partition of EX-26 does not
depend on the bridge topology when the bridges are ignored.
It still gives $sC\log_2r=O_a(B\ell)$, and $C\ge2$.
Integer widths imply $D\le W\le rsB/2$; the same bootstrap
as in EX-26 yields $s=O_a(B\log(2+B))$ and hence
$\ell=O_a(\log(2+rB))$.\hfill$\square$

Replacing binary seeds by seeds of size at most $q$ adds the
factor $\rho_q^{1/3}=O(q^{1/6})$ to (EX27.2), by using
the original debit in the feasibility inequality and the
refined binary debit only in (EX27.3). The proof above does
not treat vanishing angular curvature or groups whose
average width is below the threshold in (EX27.1).

\paragraph{EX-28: arbitrary bridge trees, without angular or group-width assumptions.}
Under the unit-antipodal-endpoint, positive-integer-width, and
all-equal-partition certificate model of EX-26, allow any fixed tree
on $r=2^t$ labelled hubs, where $t\ge1$ and $r\mid k$.
The tree and its labels are fixed before the partition is chosen,
and feasibility is required for every labelled equal partition,
with group $j$ assigned to hub $j$. Then
\[
 h=O_a\!\left(
 B^{4/3}\bigl(\log(2+rB)\bigr)^{2/3}
       +B\log(2+rB)\log r\right).
 \tag{EX28.1}
\]
In particular, for every fixed $M>0$, allowing any tree shape and
$r\le B^M$ hubs has fixed-power exponent supremum $4/3$ in this
model, using the two-hub construction for attainment under the
same supplied flat-case reduction assumed in EX-18. This removes
both extra hypotheses in EX-27. The model still charges the stated
binary-certificate selector debit; it does not identify that debit
with a lower bound for an arbitrary spectral selector.

\emph{A centered matrix estimate for pipage rounding.}
We first record explicitly the consequence of Harvey and Olver
(2014) that is needed. Let $A_i\succeq0$, $\|A_i\|\le R$,
and let $\xi$ be the integral output of randomized pipage rounding
from initial coordinates $p_i$. Put
$\mu=\sum_i p_iA_i$ and
$\sigma^2=\|\sum_i p_iA_i^2\|$.
Then
\[
 \Pr\left\{\left\|\sum_i\xi_iA_i-\mu\right\|\ge x\right\}
 \le2d\exp\left\{-\frac{x^2}{2(\sigma^2+Rx/3)}\right\}.
 \tag{EX28.2}
\]
Here $d$ is the matrix dimension. For completeness, the extension
from their stated uncentered Chernoff bound is as follows. For
fixed real $\theta\ne0$, consider
\[
 g_\theta(y)=\operatorname{tr}\exp\left(
       -\theta\mu+
       \sum_i\log\bigl(I+y_i(e^{\theta A_i}-I)\bigr)\right).
\]
Their Theorem~3.11 permits an arbitrary fixed real symmetric
matrix outside the two logarithms. It makes $g_\theta$ concave
under every pipage swap when $\theta>0$. For $\theta<0$ both
swap-direction matrices are negative semidefinite; reversing
the swap parameter reduces to the same theorem with positive
semidefinite direction matrices. The matrices inside logarithms
remain positive definite, including at the endpoints.
Thus the conditional expectation of $g_\theta$ cannot increase
during the rounding, even for adaptively selected pairs.
At an integral output it equals
$\operatorname{tr}\exp(\theta(\sum_i\xi_iA_i-\mu))$.
At the initial point, $\log(I+X)\preceq X$ gives the bound
\[
 g_\theta(p)\le
 d\exp\left\{
 \frac{e^{|\theta|R}-1-|\theta|R}{R^2}\,\sigma^2
 \right\}.
\]
The scalar exponential remainder bound is applied by functional
calculus to each $A_i$. Exponential Markov bounds for both signs,
followed by the usual Bernstein optimization, prove (EX28.2).
The precise primary source is Harvey and Olver,
\emph{Pipage rounding, pessimistic estimators and matrix
concentration}, SODA 2014, pp.~926--945,
\href{https://doi.org/10.1137/1.9781611973402.69}
{doi:10.1137/1.9781611973402.69}; the inspected author version
has Theorem~3.1 on PDF p.~4 and Theorem~3.11 on PDF p.~6.

\emph{One partition controlling all dyadic blocks.}
Let $u$ be a top eigenvector of $A$, and sort the circuit vertices
by $\alpha_a=u\cdot v_a$. Form consecutive blocks of $r$ vertices.
Within each such block, recursively split by sorting the surviving
vertices by $\alpha_a$, pairing adjacent entries, and assigning
opposite members of every pair independently and uniformly to
the two children. Leaves are the $r$ hub labels in depth-first
preorder of the bridge tree. Every hub receives one vertex from
each initial block, so the group sizes are exactly $k/r$.

Let $I$ be a dyadic interval of $L$ labels, and put $p=L/r$.
If an initial circuit-vertex block has principal-coordinate range
$\Delta$, every recursive half has principal sum equal to half
the parent's sum plus an error of absolute value at most
$\Delta/2$. Iteration gives error at most $\Delta$ from $p$
times the initial block sum. The ranges of the consecutive
initial blocks sum to at most $2B$. Since $\sum_a\alpha_a=0$,
\[
 |u\cdot R_I|\le2B
 \quad\text{for every dyadic }I,
 \qquad R_I=\sum_{j\in I}R_j.
 \tag{EX28.3}
\]

For each fixed $I$, its selected vertices are also a valid pipage
rounding output from all inclusion probabilities $p$. Along its
ancestral path, a pair of equal surviving probabilities
$(p',p')$ becomes $(2p',0)$ or $(0,2p')$ with probability
$1/2$. Pairing may depend on the earlier outcomes; the preceding
concavity argument permits this. The base polytope is the direct
sum, over the initial blocks, of uniform matroids of rank $L$.
Thus (EX28.2) applies separately to every fixed dyadic interval.

In addition to $Q_a$, apply it to the positive semidefinite lifts
\[
 V_a=\begin{pmatrix}d_a&v_a^{\mathsf T}\\v_a&d_aI\end{pmatrix}.
\]
They obey $\|V_a\|\le2B$ and $\sum_aV_a=2WI$, because
$\|v_a\|\le d_a$ and $\sum_av_a=0$.
Consequently $\sum_aV_a^2\preceq4BW I$, while
$\sum_aQ_a^2\preceq2BQ$.
For $q=u^{\mathsf T}Qu>0$, also apply the estimate to
$Q^{-1/2}Q_aQ^{-1/2}$, whose norm is at most $B/q$ and
whose sum is $2I$.
Set $\ell=\log(64r^2(D+1))$. A union bound over fewer than
$2r$ dyadic intervals and these three matrix families gives
a common outcome for which, for every such $I$,
\begin{align*}
 \|R_I\|&\le8\sqrt{BWp\ell}+8B\ell,\\
 \|H_I-2pQ\|&\le8\sqrt{BWp\ell}+8B\ell,\\
 \|Q^{-1/2}H_IQ^{-1/2}-2pI\|
   &\le8\sqrt{Bp\ell/q}+8B\ell/q\quad(q>0),
 \tag{EX28.4}
\end{align*}
where $H_I=\sum_{j\in I}H_j$.
The off-diagonal block of the centered lift sum is $R_I$, so
its norm is bounded by the norm of that whole matrix.
The total failure probability is below $1/4$.

Every preorder interval is a disjoint union of at most two
dyadic blocks at each scale. Their square-root sizes sum to at
most $8\sqrt{|I|}$. Therefore the chosen outcome also satisfies
\[
 \|R_I\|,\ \|H_I-2|I|Q/r\|
 \le M(|I|),\qquad
 M(v)=64\sqrt{BWv\ell/r}+16B\ell t,
 \tag{EX28.5}
\]
for every preorder interval. This passage uses no assertion
that an arbitrary interval itself has a pipage distribution.

\emph{Free saving after grouping consecutive hubs.}
Take any dyadic length $L\mid r$, write $p=L/r$, and group the
hub labels into consecutive blocks of that length. For a block
$I$, define
\[
 f_I=\sup_{\|z\|\le1}
       \{R_I\cdot z-\tfrac14z^{\mathsf T}H_Iz\}.
\]
We claim
\[
 \sum_I f_I\le7000(r/L)B\ell.
 \tag{EX28.6}
\]
If $pW<1024B\ell$, the trivial bound
$f_I\le\|R_I\|\le\sum_{a\text{ assigned to }I}d_a$
gives $\sum_If_I\le2W<2048(r/L)B\ell$.
Otherwise (EX28.4) gives
\[
 \tfrac12pWI\preceq PH_IP\preceq3pWI,
 \qquad \|PR_I\|^2\le129BpW\ell,
\]
where $P$ projects onto $u^\perp$.
If $pq\ge1024B\ell$, it also gives $H_I\succeq pQ$.
The unrestricted quadratic maximum and (EX28.3) then yield
\[
 f_I\le\frac1pR_I^{\mathsf T}Q^{-1}R_I
 \le\frac{4B^2}{pq}+258B\ell\le259B\ell.
\]
If $pq<1024B\ell$, use the mixed-block estimate in EX-26,
with $W/r$ there replaced by $pW$. Writing
$a_I=u^{\mathsf T}H_Iu$, it gives
\[
 f_I\le2B+\frac4{pW}\|PR_I\|^2+3a_I.
\]
Summing and using $\sum_Ia_I=2q$ proves (EX28.6), since
$2+516+6\cdot1024<7000$.
When $D=1$ the perpendicular part is absent and
$f_I\le|R_I|\le2B$, so the same claim holds.

\emph{Absorbing both vector and matrix residuals on the bridge tree.}
Project hubs into the unit ball. For the chosen length $L$,
replace $R_j$ and $H_j$ by the means $\overline R_j$ and
$\overline H_j$ of their consecutive block. Each residual
sums to zero on whole blocks. A tree-edge subtree is a preorder
interval, so its residual flow has contributions from at most
two partial blocks. By (EX28.5), both its vector residual norm
and its matrix residual operator norm are at most $4M(L)$.
For hub positions in the unit ball,
\[
 \bigl\|zz^{\mathsf T}-ww^{\mathsf T}\bigr\|_*
 \le(\|z\|+\|w\|)\|z-w\|\le2\|z-w\|.
\]
Tree summation by parts therefore bounds the residual linear
term by $4M(L)$ times total bridge length and the absolute
residual quadratic term, including its factor $1/4$, by
$2M(L)$ times that length. If $6M(L)\le h$, both are
absorbed by the width-$h$ bridge cost. The remaining saving
is at most $\sum_I f_I$: each of the $L$ hubs in one block
has the same averaged objective, scaled by $1/L$.
For $L=1$ the residuals vanish and no absorption condition is
needed.

If $h<256B\ell t$, the second term of (EX28.1) already
bounds it. Otherwise set
\[
 X=\frac{h^2r}{2^{24}BW\ell}.
\]
For $X<1$ take $L=1$. For $1\le X<r$, take the largest
dyadic $L\le X$, so $L\ge X/2$.
For $X\ge r$, take $L=r$. In the latter two cases
\[
 6M(L)\le\frac{384}{2^{12}}h+96B\ell t<h.
\]
The single-block case has $R_I=0$ and $H_I\succeq0$,
so it permits no positive saving. In the other two cases,
(EX28.6) gives the uniform bound
\[
 \text{saving before selector}
 \le2^{38}\frac{B^2W\ell^2}{h^2}.
 \tag{EX28.7}
\]
For $X<1$, this follows by multiplying the $L=1$ bound by
$h^2<2^{24}BW\ell/r$; otherwise use $L\ge X/2$.

The certificate-based tree-flow partition in (EX27.3)
still gives $h\le20BCt$. Feasibility for the new partition
and $2W\le kB$ give
$h^2Ct\le(2^{37}/a)B^3\ell^2$.
Combining them proves
$h=O_a(B^{4/3}\ell^{2/3})$ in the large-$h$ case.
The preceding small-$h$ case adds $O(B\ell t)$.
Finally (EX28.6) with $L=1$ yields
$sCt=O_a(B\ell)$; $C\ge2$ and
$D\le W\le rsB/2$ give the same bootstrap
$s=O_a(B\log(2+B))$, hence
$\ell=O_a(\log(2+rB))$. This proves (EX28.1).
\hfill$\square$

For $q$-point seeds the first term in (EX28.1) is multiplied
by $\rho_q^{1/3}=O(q^{1/6})$; the second is unchanged.
For polynomially many hubs and subpolynomially growing seed
sizes, the fixed-power supremum in the stated model remains
$4/3$.

The same conclusion holds for any fixed connected labelled hub
graph whose bridge widths are all at least $h$. Delete edges to
obtain a fixed spanning tree and lower its retained widths to
$h$. These operations can only increase the available saving,
so feasibility of the original graph implies the feasibility
needed above. This bounds the minimum bridge width; it does
not bound the maximum width of one edge in a nonuniform graph.
The graph and the group-to-hub assignment must still be fixed
before the partition is chosen. Choosing a new graph or
permuting the groups after seeing the partition is not covered.
Arbitrary non-antipodal layouts, arbitrary radii with many hubs,
and a better actual spectral selector estimate are also not
covered by EX-28.

\paragraph{EX-29: arbitrary common-origin radii and polynomially many hubs.}
The common-origin radial model of EX-23 can also be combined with
the fixed labelled graph model of EX-28. Assume positive integer
edge widths, maximum incident load $B$, $k$ circuit vertices,
and endpoints $r_e^+x_e,-r_e^-x_e$ with $r_e^\pm\ge0$ and
$r_e^++r_e^->0$. Let $r=2^t\mid k$, $t\ge1$.
Fix a connected labelled graph on these $r$ hubs before choosing
the partition; group $j$ is assigned to hub $j$. Suppose every
labelled equal partition permits strict improvement with all
bridge widths at least $h$ and selector debit at least $aktC$,
where $C$ is the binary pointwise-averaging certificate cost.
Then
\[
 h=O_a\!\left((\log(2+B))^2\log(2+rB)
                   \bigl(B^{4/3}+B\log r\bigr)\right).
 \tag{EX29.1}
\]
In particular, for each fixed $M>0$, if $r\le B^M$, then
\[
 h=O_{a,M}\bigl(B^{4/3}(\log(2+B))^3\bigr).
 \tag{EX29.2}
\]
The fixed-power exponent supremum is therefore $4/3$ throughout
this stated model, with attainment conditional on the same
supplied flat-case reduction assumed in EX-18. The conclusion bounds the minimum bridge
width and retains the prescribed certificate debit and the
fixed labelled assignment quantifiers.

\emph{A simultaneous recursive halving lemma.}
We give the needed deterministic replacement for the pipage
construction. Use $m\ge1$ radial bands and the notation
$W_j,Q_j,Q_{aj},v_{aj},u_j$ of EX-23. Throughout this proof,
$D$ is the span dimension of the directions retained after the
radial truncation below. Project hubs and the selector ellipsoid
onto that span: retained fan and bridge lengths and the
ellipsoid trace can only decrease. This makes the later bound
$D\le W$ refer to the same retained edges. Let
\[
 \ell=\ln(64m(D+1)),\qquad K=5\sqrt{5m\ell},
 \qquad \Theta=10K^2=1250m\ell.
\]
For each parent set of an even number of circuit vertices,
pair its vertices arbitrarily and assign opposite members of
each pair to its two equal children. One Banaszczyk application
chooses the pair signs simultaneously for all bands.

Here are the blocks and their bounds. For a pair $(a,b)$ use
the principal block $(\alpha_{aj}-\alpha_{bj})/(2B)$ and
the selector block
\[
 \frac{r_j}{4B}Q_0^{-1/2}\frac{v_{aj}-v_{bj}}2.
\]
Both have norm at most one, by EX-23. Use three additional
matrix covariance-whitening blocks per band, from the pair
differences divided by two of
\[
 V_{aj}=\begin{pmatrix}d_{aj}&v_{aj}^{\mathsf T}\\
                      v_{aj}&d_{aj}I\end{pmatrix},
 \qquad Q_{aj},\qquad Q_j^{-1/2}Q_{aj}Q_j^{-1/2},
 \tag{EX29.3}
\]
where $d_{aj}=\sum_{e\ni a,\ e\in j}b_e$ and the last
family is omitted when $q_j=\lambda_{\min}(Q_j)=0$.
All covariance inverses and selector pseudoinverses are restricted
to their ranges, and zero blocks are omitted. Every whitened
pair vector has norm at most one, by the covariance argument
in EX-23. There are at most $5m$ blocks, so dividing the whole
input vector by $\Lambda=5\sqrt{5m}$ gives norm at most $1/5$.

For clarity, if $0\preceq M_a\preceq RI$ is any of the three
positive semidefinite families and $M_S=\sum_{a\in S}M_a$
is its parent sum, the pair matrices $H_{ab}=(M_a-M_b)/2$
satisfy
\[
 \sum_{(a,b)}H_{ab}^2\preceq\frac R2 M_S.
\]
Tropp's Gaussian matrix-series bound permits the image norm
threshold $\sqrt{R\|M_S\|\ell}$ with failure probability at
most $1/(32m)$. For the principal block use threshold
$4\sqrt m$, and for the selector block use
$\|Q_0^{1/2}z\|\le4\sqrt m(C/2)$.
Their Gaussian failure probabilities are at most $1/(16m)$
each. The product body's Gaussian measure is at least
$1-2m/(16m)-3m/(32m)=25/32>1/2$.
Thus Banaszczyk gives equal-size children for which
\[
 M_{S'}=\tfrac12M_S+E,\qquad
 \|E\|\le K\sqrt{R\|M_S\|}.
 \tag{EX29.4}
\]
At the same split the principal error is at most $45mB$,
and the selector resultant error is at most $90mBC/r_j$.
No count correction is used: the pair assignments give exact
halves.

Iterate this construction to depth $t$. The following elementary
recursion makes its matrix consequences explicit. If a node has
relative size $p=2^{-d}$, (EX29.4) implies
\[
 \|M_{\rm node}-pM_{\rm root}\|
 \le5K\sqrt{Rp\|M_{\rm root}\|}+5K^2R.
 \tag{EX29.5}
\]
Indeed, writing $b_S=\sqrt{\|M_S\|}$ gives
$b_{S'}\le b_S/\sqrt2+K\sqrt{R/2}$ and hence
$b_{\rm node}\le\sqrt p\,b_{\rm root}
 +K\sqrt R/(\sqrt2-1)$.
Summing the deviation recursion with weights
$2^{-(d-i)}$ gives coefficients
$2/(\sqrt2-1)<5$ for both terms in (EX29.5).

The three root sums in (EX29.3) are $2W_jI$, $2Q_j$,
and $2I$; their individual norm bounds are $2B$, $B$,
and $B/q_j$, respectively. Consequently every dyadic
hub-label block $I$, with $p=|I|/r$, simultaneously obeys
\begin{align*}
 \|R_{I,j}\|,\ \|H_{I,j}-2pQ_j\|
     &\le4\sqrt{\Theta BW_jp}+\Theta B,\\
 \|Q_j^{-1/2}H_{I,j}Q_j^{-1/2}-2pI\|
     &\le4\sqrt{\Theta Bp/q_j}+\Theta B/q_j
                              \quad(q_j>0),\\
 |u_j\cdot R_{I,j}|&\le\Theta B,
       \qquad \|R_{I,j}\|\le\Theta BC/r_j.
 \tag{EX29.6}
\end{align*}
The last two bounds follow from geometric summation of the
constant split errors and $\sum_av_{aj}=0$. Leaves are
assigned the fixed spanning tree's depth-first preorder labels,
so each group has exactly $k/r$ vertices. The construction
chooses one partition satisfying all band constraints together.

\emph{Coarse fan bounds and tree flows.}
For a dyadic length $L\mid r$, put $p=L/r$ and group the
hub labels into consecutive blocks of length $L$. For one band,
define
\[
 f_{I,j}=\sup_{\|z\|\le1}
       \{R_{I,j}\cdot z-\tfrac14z^{\mathsf T}H_{I,j}z\}.
\]
The calculation of EX-28, with (EX29.6), gives
\[
 \sum_I f_{I,j}\le2000(r/L)\Theta B.
 \tag{EX29.7}
\]
Here are constants sufficient to verify all regimes. If
$pW_j<256\Theta B$, the sum is at most
$2W_j<512(r/L)\Theta B$. Otherwise
$pW_jI/2\preceq P_jH_{I,j}P_j\preceq3pW_jI$ and
$\|P_jR_{I,j}\|^2\le33\Theta BpW_j$.
If $pq_j\ge256\Theta B$, full whitening gives
$H_{I,j}\succeq pQ_j$ and $f_{I,j}\le67\Theta B$.
In the remaining case the mixed-block estimate gives
$f_{I,j}\le\Theta B+132\Theta B+3a_{I,j}$,
where $a_{I,j}=u_j^{\mathsf T}H_{I,j}u_j$ and
$\sum_Ia_{I,j}=2q_j<512(r/L)\Theta B$.
The total is below $1669(r/L)\Theta B$.
The one-dimensional case follows directly from the principal
bound.

Decomposing intervals into dyadic blocks also gives
\[
 \|R_{I,j}\|,\ \|H_{I,j}-2|I|Q_j/r\|
 \le M_j(|I|),\qquad
 M_j(v)=32\sqrt{\Theta BW_jv/r}+2\Theta Bt.
 \tag{EX29.8}
\]
Average both the vectors and matrices within consecutive
blocks of length $L$. As in EX-28, each residual tree-edge
flow has norm at most $4M_j(L)$, and its linear and quadratic
terms together cost at most $6M_j(L)$ per unit bridge length.

For unequal radii this statement is applied after projecting
each hub separately to the ball of radius $r_j$ for band $j$
and dividing its position by $r_j$. This can only decrease
that band's fan costs. The exact distance inequality in
EX-23 gives its normalized quadratic upper bound, multiplied
by $r_j$. Projection is nonexpansive, so this prefactor cancels
the factor $1/r_j$ in the variation of the normalized hub
positions. The residual bound is therefore still
$6M_j(L)$ times the original tree variation
$\mathrm{TV}(z)=\sum_{ab\in E(T)}\|z_a-z_b\|$.

The averaged vector flow admits an additional bound. Every
tree-edge interval consists of whole averaging blocks and at
most two partial blocks. Its whole-block middle is an interval,
whose resultant is at most $2t\Theta BC/r_j$ by (EX29.6).
The two fractional end-block resultants contribute at most
$2\Theta BC/r_j$. Thus the averaged flow is at most
$4\Theta BCt/r_j$. Dropping the negative quadratic term
and summing by parts shows that the averaged band objective
is at most
\[
 \frac{4\Theta BCt}{r_j}\,\mathrm{TV}(z).
 \tag{EX29.9}
\]

\emph{Truncation and optimization.}
Put $r_{\max}=\max_e\max(r_e^+,r_e^-)$ and
$\delta=atC/(4B)$. Edges with maximum radius below
$\delta$ have total target cost at most $aktC/4$.
Since $r_{\max}\le C$, feasibility is impossible when
$r_{\max}\le\delta$; otherwise the retained edges lie in
$m=O_a(\log(2+B))$ dyadic bands with parameters
$r_j=2^{-j}r_{\max}$. Write $W=\sum_jW_j\le kB/2$.

If $h<96m\Theta Bt$, the second term in (EX29.1)
will suffice. Otherwise, independently for each band set
\[
 X_j=\frac{h^2r}{2^{22}m^2\Theta BW_j}.
\]
Take $L_j=1$ for $X_j<1$, the largest dyadic
$L_j\le X_j$ for $1\le X_j<r$, and $L_j=r$
for $X_j\ge r$. In the latter two cases
\[
 6M_j(L_j)\le\frac{192}{2^{11}}\frac hm
                 +12\Theta Bt<\frac h{2m}.
\]
For $L_j=1$ residuals vanish. Hence the total residual
debit is at most $(h/2)\mathrm{TV}(z)$.
By (EX29.7), the remaining averaged band objective is at most
\[
 a_j=2^{35}\frac{m^2\Theta^2B^2W_jr_j}{h^2}.
\]
This includes $X_j<1$ by the same multiplication argument
as EX-28; when $X_j\ge r$, the averaged objective is zero.
By (EX29.9) it is also at most $b_j\mathrm{TV}(z)$,
where $b_j=4\Theta BCt/r_j$.
Writing $v=\mathrm{TV}(z)$, retained-edge saving less
bridge cost is therefore bounded by
\begin{align*}
 \sum_j\min\{a_j,b_jv\}-\tfrac h2v
 &\le \frac{\sqrt{2^{37}m^3\Theta^3B^3CtW}}h\sqrt v
                   -\tfrac h2v\\
 &\le 2^{36}\frac{m^3\Theta^3B^3CtW}{h^3}.
 \tag{EX29.10}
\end{align*}
The first inequality uses $\min(a,b)\le\sqrt{ab}$ and
$\sum_j\sqrt{W_j}\le\sqrt{mW}$. Strict improvement
after the remaining selector debit $3aktC/4$, followed by
$W\le kB/2$, now gives
$h=O_a(m\Theta B^{4/3})$.

Finally, (EX29.7) with $L=1$, ignoring bridges, bounds
retained free saving by
$2000r\Theta B\sum_jr_j\le4000r\Theta BC$.
Thus feasibility forces $st=O_a(\Theta B)$, $s=k/r$.
Integer widths give $D\le W\le rsB/2$.
Using $m=O_a(\log(2+B))$ and $t=\log_2r\ge1$,
the resulting inequality implies
$s=O_a(B(\log(2+B))^2)$ by the same logarithmic
bootstrap as EX-23. Consequently
$\ell=O_a(\log(2+rB))$ and
$m\Theta=O_a((\log(2+B))^2\log(2+rB))$.
This proves (EX29.1); polynomially many hubs give
(EX29.2). Deleting extra hub-graph edges and lowering
retained bridge widths to $h$ proves the connected-graph
statement.\hfill$\square$

For seeds of size at most $q$, use their original debit $K_0$
in the selector and radial cutoff, and the refined binary cost
$C'\le\rho_qK_0$ only in the selector blocks and (EX29.9).
The first term in (EX29.1) gains the factor
$\rho_q^{1/3}=O(q^{1/6})$; the second and the dimension
bootstrap are unchanged. Subpolynomial seed sizes retain
the $4/3$ fixed-power supremum. This theorem still does not
cover non-common-origin target segments, graphs or label
assignments selected after seeing the partition, or a selector
whose actual spectral cost is smaller than the charged
pointwise-averaging certificate.

\paragraph{EX-30: graphs chosen after the partition, with subpolynomially many hubs.}
Retain the radial, integer-width and selector assumptions of
EX-29, but now allow the connected hub graph and the assignment
of the $r$ equal groups to its hubs to be chosen after seeing
each partition. If every equal partition is feasible, then
\[
 h=O_a\!\left(
 rB^{4/3}
 \frac{\bigl(\log(2+B)\log(2+rB)\bigr)^{2/3}}
      {(\log_2r)^{1/3}}\right).
 \tag{EX30.1}
\]
Thus if $r=B^{o(1)}$, the fixed-power exponent supremum is
still $4/3$, even with these adaptive choices, with attainment
conditional on the supplied flat-case reduction in EX-18. This does not
give the same conclusion for adaptively chosen graphs with
a polynomially growing number of hubs.

\emph{Proof.}
Use the simultaneous leaf partition from EX-29, with its
$m$ radial bands and $\Theta=1250m\ell$.
The free-saving bound (EX29.7), at $L=1$, is
$2000r\Theta Br_j$ for band $j$.
Each individual group's band resultant has norm at most
$\Theta BC/r_j$. For any spanning tree of any subsequently
chosen connected hub graph, every edge flow is a sum over
some subset of the groups, so its norm is at most
$r\Theta BC/r_j$. The supporting fan inequality at the
common origin therefore bounds the same band's saving by
$(r\Theta BC/r_j)v$, where $v$ is the tree's total length.
The dyadic radial sum gives
\[
 \sum_j\min\{2000r\Theta Br_j,
                    r\Theta BCv/r_j\}
 \le4\sqrt{2000}\,r\Theta B\sqrt{Cv}.
\]
After subtracting the width-$h$ bridge cost and maximizing
over $v\ge0$, the retained saving is at most
$8000r^2\Theta^2B^2C/h$.
The same radial truncation loses at most $aktC/4$, so
\[
 h=O_a\!\left(\frac{r\Theta^2B^2}{st}\right),
 \qquad s=k/r,\quad t=\log_2r.
 \tag{EX30.2}
\]
This argument is uniform over the graph and label assignment.

For a second partition, apply the equal-group local-swap
argument of EX-26 to the full directional resultants
$v_a=\sum_{e\ni a}\pm b_ex_e$. It gives
$\|R_i\|\le2B\sqrt s$ for every group, independently
of the radii. Since all group resultants sum to zero,
every subset resultant has norm at most $rB\sqrt s$.
Tree summation by parts and the supporting fan inequality
therefore forbid strict improvement for any connected hub
graph when $h\ge rB\sqrt s$. Consequently
\[
 h\le rB\sqrt s.
\]
Combining this with (EX30.2) gives
$h=O_a(rB^{4/3}\Theta^{2/3}t^{-1/3})$.
The free-saving bootstrap in EX-29 is independent of the
bridge graph and still gives
$\Theta=O_a(\log(2+B)\log(2+rB))$, proving (EX30.1).
As before, the two partitions need not coincide, because
feasibility is assumed for every equal partition.
\hfill$\square$

With $q$-point seeds the right side gains the factor
$\rho_q^{1/3}=O(q^{1/6})$. If both the number of hubs
and the seed sizes are subpolynomial in $B$, these factors
do not change the fixed-power exponent. The theorem still
concerns common-origin segments and the prescribed
pointwise-averaging selector cost.

\paragraph{EX-31: a sparsity obstruction to the balanced matching-cover shortcut.}
Let $G_r(q)$ be the graph tensor of $K_{r,r}$ and let
$H_r(q)=K_r(q)\otimes K_r(q)$ be the tensor of the two
disjoint cliques on its two parts. Use the canonical tensor-product
edge bases, with $r\ge2$ and $q\ge2$. Suppose a cover of
$H_r(q)$ uses $A$ restrictions of $G_r(q)$, such that in every
restriction each left local map has at most $u$ nonzero entries
in any row and each right local map has at most $v$ nonzero
entries in any column. Then, over every field,
\[
 A(uv)^r\ge q^{r(r-1)/2}.
 \tag{EX31.1}
\]
The same necessary inequality holds for a border limit of such
$A$-term covers with the same sparsity bounds.

\emph{Proof.}
Fix one valid left-clique output basis tuple, for example the
tuple with all its edge indices zero. The corresponding row
of the left-versus-right flattening of $H_r(q)$ is the full
right-clique tensor: it has exactly $q^{r(r-1)/2}$ nonzero
coordinates, all with coefficient one.

For one restricted source term, each fixed left output basis
vector is compatible with at most $u$ input basis vectors.
Thus there are at most $u^r$ compatible tuples of source
inputs at the left parties. Such a tuple specifies every
source edge index, since every source edge has exactly one
left endpoint. For each of these assignments, the right
local maps have at most $v$ possible nonzero output coordinates
each, hence at most $v^r$ right output tuples. The resulting
row has support size at most $(uv)^r$. Cancellation can only
decrease that support. The support of a sum of $A$ rows is
contained in their union, proving (EX31.1).

For the border assertion, the set of row vectors with at most
$A(uv)^r$ nonzero coordinates is a finite union of coordinate
linear subspaces, hence Zariski closed. The same row-support
bound therefore holds on the border closure. This argument
does not require a positivity assumption.\hfill$\square$

In particular, local partial-permutation maps require
$A\ge q^{r(r-1)/2}$. More generally, at $q=2$, polynomial
row and column sparsities in $r$ force
$\log_2 A\ge r(r-1)/2-O(r\log r)$.
They cannot establish the $2^{o(r^2)}$ balanced-cover gate in
SH-33. Counting only $r!$ matching terms is therefore
insufficient if each term merely selects or relabels basis
coordinates. A successful determinant or Pfaffian identity
could evade this obstruction through dense local maps; no
bound for arbitrary restrictions, or for arbitrary local
changes of basis, follows from this proposition.

\subsection{A geometric model with both flatness and contraction}

\paragraph{EX-32: CAT(0), all flats and all trees do not imply mutual exclusion.}
For every finite label set $S$, let $\mathcal M(S)$ be the
pointwise closure of the pseudometrics obtained by mapping $S$
into a complete CAT(0) cube complex and multiplying distances
by an arbitrary nonnegative constant. Then this family is
closed under pullback, nonnegative scaling and pointwise limits;
each of its members has a CAT(0) realization; and it contains
every finite Euclidean configuration and every finite weighted
tree metric. Nevertheless every $d\in\mathcal M(
\operatorname{Mat}_8(\mathbb F_2))$ satisfies
\[
 D_d(8)\le\frac{255\sqrt2}{128}D_d(1),
 \qquad
 \left(\frac{255\sqrt2}{128}\right)^2
      =\frac{65025}{8192}<8.
 \tag{EX32.1}
\]
Thus this family has both the flat property and a uniform
strict rank-step contraction. It is a model for the listed
geometric properties, not a claimed class of actual prices
over any field.

\emph{Proof.}
Pullbacks, scaling and closure follow directly from the
definition. Euclidean space $\mathbb R^D$ has its usual
complete unit cubulation, so every finite Euclidean metric
already occurs before taking closure. Approximate the edge
lengths of any finite weighted tree by nonnegative rational
lengths. After a common rescaling and subdivision, each
positive rational-edge tree is a unit-edge CAT(0) cube complex.
Limits also permit zero edge lengths. Hence all the weighted
tree metrics belong to $\mathcal M$.

For a convergent sequence of such finite distance arrays,
choose one labelled point as basepoint in each realizing
space. All labelled points stay a bounded distance from their
basepoints. Their metric ultralimit therefore realizes the
limit array, after identifying zero-distance labels. Pointwise
ultralimits of constant-speed geodesics are geodesics, and
the CAT(0) squared-distance comparison passes to the limit.
Taking the metric completion if necessary preserves that
comparison. This proves the asserted CAT(0) realization.

Fujiwara and Toyoda (2012), Theorem~1.4, arXiv version~1,
p.~2, prove that every complete CAT(0) cube complex has
Izeki--Nayatani invariant at most $1/2$. The independently
checked Poincar\'e argument PG-IN-GATE consequently gives
\[
 D_d(n)^2\le8(1-2^{-n})^2D_d(1)^2
\]
for every finite image in such a space. This inequality is
unchanged by scaling and passes to pointwise limits because
both sides are finite sums of squared distances. At $n=8$
it is exactly (EX32.1).\hfill$\square$

Even the smaller family of pointwise limits of finite images
in finite products of complete real trees contains all flats
and all finite weighted trees. The invariant-zero product
result of Izeki and Nayatani, cited and checked in PG-IN-GATE,
gives the sharper $n=4$ bound $D_d(4)\le(15/8)D_d(1)$
for that smaller family. Neither construction verifies the
tensor-restriction inequalities required of actual prices.
Accordingly, it identifies a limitation of a purely geometric
exclusivity argument and does not refute Q1c over either
$\mathbb C$ or $\mathbb F_2$.

\paragraph{EX-33: a sharp obstruction to the finite-characteristic Weyl shortcut.}
Let $F$ have characteristic $p>0$. If matrices
$A_1,\ldots,A_n,B_1,\ldots,B_n\in\operatorname{Mat}_d(F)$
satisfy
\[
 [A_i,B_j]=\delta_{ij}I_d,
\]
without any assumptions on $[A_i,A_j]$ or $[B_i,B_j]$, then
\[
 p^n\le d^2,\qquad n\le2\log_p d.
 \tag{EX33.1}
\]
For $d=p^s$ this bound is attained with $n=2s$.
Thus omitting same-side commutation does not permit a family
with $n=d^c$ for fixed $c>0$ and unbounded $d$.

\emph{Proof.}
Let $\mathcal B$ be the unital subalgebra generated by the
$B_j$, and let $D_i=\operatorname{ad}(A_i)$ on
$\mathcal B$. It is a derivation with
$D_i(B_j)=\delta_{ij}I$. On every ordered monomial,
\[
 D_i(B_1^{e_1}\cdots B_n^{e_n})
   =e_i B_1^{e_1}\cdots B_i^{e_i-1}\cdots B_n^{e_n}.
\]
This formula uses the Leibniz rule and does not commute
any two $B_j$. The $D_i$ commute on $\mathcal B$ because
their commutator is a derivation vanishing on every generator.

Suppose the $p^n$ ordered monomials with $0\le e_i<p$
had a nontrivial linear relation. Choose an exponent vector
$e$ of maximum total degree among its nonzero coefficients,
and apply $D_1^{e_1}\cdots D_n^{e_n}$.
Every other term vanishes: survival would require its
exponent vector to be coordinatewise at least $e$, which,
by maximal total degree, would make it equal to $e$.
The surviving term is its nonzero coefficient multiplied
by $\prod_i e_i!$ and $I_d$. All these factorials are
nonzero in characteristic $p$, a contradiction.
Thus $\dim_F\mathcal B\ge p^n$, while
$\mathcal B\subseteq\operatorname{Mat}_d(F)$ has dimension
at most $d^2$.

For sharpness, on
$V=(F[x]/(x^p))^{\otimes s}$ let $M_j$ multiply by
$x_j$ and let $\partial_j$ differentiate with respect to
$x_j$. Differentiation is well-defined on the quotient,
and $[\partial_j,M_j]=I$; operators on different tensor
factors commute. Set
\[
 A_{2j-1}=\partial_j,\quad B_{2j-1}=M_j,
 \qquad A_{2j}=M_j,\quad B_{2j}=-\partial_j.
\]
These $2s$ pairs satisfy all the required cross commutators
on the $p^s$-dimensional space $V$, including in
characteristic two.\hfill$\square$

\paragraph{EX-34: every folding into at most three superparties permits the balanced conversion.}
Fix $r\ge2$, $q\ge2$ and an arbitrary partition of the $2r$
labelled vertices into at most three groups. After grouping the
tensor factors according to this partition, the tensor
$H_r(q)=K_r(q)\otimes K_r(q)$ is an exact restriction of
$G_r(q)=K_{r,r}(q)$ over every field.

\emph{Proof.}
Add empty groups if needed. Write $a_i$ and $b_i$ for the
numbers of left and right vertices in group $i$, and put
$c_i=a_i-b_i$. Then $\sum_i c_i=0$. Between two distinct
superparties the source and target respectively have
\[
 S_{ij}=a_i b_j+a_j b_i,\qquad
 T_{ij}=a_i a_j+b_i b_j
\]
parallel $q$-dimensional wires. In particular,
\[
 T_{ij}-S_{ij}=c_i c_j.
 \tag{EX34.1}
\]
If every $c_i$ is zero, all these demands agree. Otherwise,
after relabelling and possibly replacing every $c_i$ by
$-c_i$, assume $c_1,c_2\ge0$ and
$c_3=-c_1-c_2$. The only possible positive excess is the
$c_1c_2$ wires demanded between groups $1$ and $2$.
The unused source capacities on the other two pairs are
\[
 S_{13}-T_{13}=c_1(c_1+c_2)\ge c_1c_2,
 \qquad
 S_{23}-T_{23}=c_2(c_1+c_2)\ge c_1c_2.
\]
Route those extra $1$--$2$ wires through group $3$, using
one unused wire from each of these pairs. The local map at
group $3$ contracts their two coordinates with
$\sum_{j=1}^q e_j^*\otimes e_j^*$, producing exactly one
$q$-dimensional wire between groups $1$ and $2$.
Use the remaining demanded wires directly and discard
surplus wires by coordinate projections. This construction
uses only zero and one coefficients.

Edges internal to a superparty contribute fixed nonzero
local vectors. They impose no additional obstruction:
a local linear map can send the source's internal vector
to the target's internal vector while carrying out the
specified maps on the boundary wires. Thus the complete
grouped target is obtained by local linear maps.
\hfill$\square$

Consequently no restriction-monotone invariant applied only
after such a folding can obstruct the balanced conversion;
this includes ordinary ranks and their border versions, as
well as three-party Koszul flattenings or substitution bounds.
The assertion is uniform in the chosen partition and in the
field. A slightly weaker, but still asymptotically exact,
statement holds for four groups, as follows.

\emph{Four-group extension.}
For every grouping into at most four superparties,
\[
 H_r(q)^{\otimes2}\le G_r(q)^{\otimes2}
 \tag{EX34.2}
\]
over every field. Continue to use the surpluses in
(EX34.1). If at most one $c_i$ has either sign, route every
positive excess through that opposite-sign group. The total
load required on the wire bundle incident to a group of
absolute imbalance $a$ is $a(P-a)\le aP$, where $P$
is the total positive imbalance; its available surplus is
$aP$. This gives an integral routing already for one copy.

The remaining case has two positive and two negative
imbalances. Denote them by $a,b,-c,-d$, with
$a+b=c+d$. Relabelling within the pairs and interchanging
the two pairs allows us to assume
$a\le c\le d\le b$. The two excess demands are $ab$
between the positive groups and $cd$ between the negative
groups; the four cross capacities are $ac,ad,bc,bd$.
Set
\[
 x=ac-\frac{a^2}{2},\qquad y=\frac{a^2}{2}.
\]
Route the $ab$ demand through the $c$-group in amount $x$
and through the $d$-group in amount $ab-x$. Route the $cd$
demand through the $a$-group in amount $y$ and through the
$b$-group in amount $cd-y$. All four amounts are
nonnegative. The first two cross capacities $ac,ad$ are
saturated, while the unused capacities on $bc,bd$ are
respectively $(c-a)^2$ and $(d-a)^2$. Thus this is a feasible
half-integral wire routing. Taking two tensor copies makes
every wire count integral, proving (EX34.2). Internal edges
are handled as in the three-group argument.

It follows that an asymptotic restriction monotone applied
only after any fixed folding into at most four groups cannot
obstruct the balanced conversion. One-copy border invariants
can still distinguish some four-group instances; EX-35 gives
an explicit example. No assertion is made for five or more
superparties, or for invariants retaining all $2r$ parties.

\paragraph{EX-35: the binary four-cycle has a one-copy border obstruction that disappears at the tensor square.}
Let $C_4(2)$ have cyclically labelled parties $1,2,3,4$, and
let $P=\mathrm{EPR}_2^{13}\otimes\mathrm{EPR}_2^{24}$.
Over any field of characteristic different from two,
$P$ is not a border restriction of $C_4(2)$. Nevertheless
$P^{\otimes2}$ is an exact restriction of
$C_4(2)^{\otimes2}$ over every field, by (EX34.2).

Here is a direct polynomial proof of the border assertion.
For a four-qubit tensor $T$, contract parties $1,3$ with
coordinates $x=(x_0,x_1)$ and $z=(z_0,z_1)$. Regard the
result as a $2\times2$ matrix $M_T(x,z)$ on parties $2,4$,
and form its determinant. This is a polynomial of bidegree
$(2,2)$, whose coefficient matrix $K_T$ in the monomial
bases $(x_0^2,x_0x_1,x_1^2)$ and
$(z_0^2,z_0z_1,z_1^2)$ is $3\times3$.
Then $I_6(T)=\det K_T$ is a degree-six polynomial in $T$.

A restricted cycle tensor has coefficients
\[
 T_{abcd}=\operatorname{tr}(A_aB_bC_cD_d),
 \qquad A_a,B_b,C_c,D_d\in\operatorname{Mat}_2(F).
\]
Work first at generic local maps over an algebraic closure.
Each of the matrix pencils $A(x)$ and $C(z)$ has two
distinct rank-one members; choose those as its physical
basis. For every corner pair $(x_i,z_j)$, both selected
matrices have rank one. Writing them as $uv^{\mathsf T}$
and $st^{\mathsf T}$ gives
\[
 (M_T)_{bd}
   =(v^{\mathsf T}B_b s)(t^{\mathsf T}D_d u),
\]
so this output matrix has rank at most one. Consequently
the four corner entries of $K_T$ vanish, and it has the form
\[
 K_T=\begin{pmatrix}0&*&0\\ *&*&*\\0&*&0\end{pmatrix}.
\]
Its determinant is zero. Invertible changes of the physical
bases multiply this determinant by nonzero characters, so
its vanishing is basis-independent. Polynomiality and the
density of the generic set extend the identity to all local
maps and to their border closure.

For the crossed Bell-pair tensor,
$M_P(x,z)=(x_0z_0+x_1z_1)I_2$, whence
\[
 K_P=\operatorname{diag}(1,2,1),\qquad I_6(P)=2.
\]
This is nonzero in the stated characteristics, proving the
obstruction. The displayed invariant gives no obstruction
in characteristic two; EX-36 settles that case by a separate
explicit construction.
An exact symbolic expansion has $144$ monomials, all with
coefficients $+1$ or $-1$. In particular the coefficient
gcd is one, so dividing the invariant by two cannot repair
this failure. The program
\path{programs/ex_cycle_sextic.py}, parameters
\path{inputs/ex_cycle_sextic_parameters.json}, and full
coefficient record \path{outputs/ex_cycle_sextic_results.json}
preserve this computation and the target evaluation.

The invariant and the complex cycle variety are identified
in Bernardi, De Lazzari, and Gesmundo (2023),
\emph{Dimension of tensor network varieties},
\emph{Communications in Contemporary Mathematics},
25(10), 2250059,
\href{https://doi.org/10.1142/S0219199722500596}
{doi:10.1142/S0219199722500596}, Theorem~5.3,
printed pp.~23--24 of the MEGA preprint. Their primary
\href{https://fulges.github.io/code/BDG-DimensionTNS.html}
{supplementary page} provides the degree-six calculation.
The rank-one-pencil argument above verifies the needed
vanishing independently and supplies the stated
characteristic qualification. This example warns against
inferring an asymptotic obstruction even from a proved
one-copy border obstruction.

\paragraph{EX-36: an explicit characteristic-two border conversion, with an all-field exact obstruction.}
For every field $F$ of characteristic two, there are polynomial
local maps over $F[t]$ taking $C_4(2)$ to
\[
 tP+t^2U,\qquad
 P=\mathrm{EPR}_2^{13}\otimes\mathrm{EPR}_2^{24},\qquad
 U=e_1\otimes e_0\otimes e_0\otimes e_0.
 \tag{EX36.1}
\]
However, over every field, $P$ is not an exact restriction of
$C_4(2)$. Thus this is a same-field border conversion in
characteristic two, in contrast to the characteristic-not-two
border obstruction in EX-35.

\emph{Explicit degeneration.}
Put
\[
 E=\begin{pmatrix}0&0\\0&1\end{pmatrix},\qquad
 X=\begin{pmatrix}0&1\\1&0\end{pmatrix}.
\]
Represent each local map by its two matrix slices and choose
\[
 \begin{array}{c|cc}
     &0&1\\\hline
 A&I&tE\\
 B&I&X\\
 C&tE&I\\
 D&I&X
 \end{array}
\]
With the usual cyclic contraction, direct multiplication
over the integer polynomial ring gives
\[
 \operatorname{tr}(A_aB_bC_cD_d)
 =2\,\mathbf1_{a=0,c=1}\delta_{bd}
   +t\delta_{ac}\delta_{bd}
   +t^2\mathbf1_{a=1,b=0,c=0,d=0}.
 \tag{EX36.2}
\]
Indeed $\operatorname{tr}I=2$, $\operatorname{tr}E=1$,
$\operatorname{tr}X=0$, and $EXE=0$ verify all sixteen
entries. In characteristic two the first term vanishes,
which proves (EX36.1) without any extension of the ground
field or interpolation assumption.

\emph{Exact obstruction.}
Suppose an exact representation existed, and extend the
field to an algebraic closure. Conciseness of $P$ forces
$A_0,A_1$ to be linearly independent. Their matrix pencil
therefore contains a nonzero rank-one member
$A(\alpha)=uv^{\mathsf T}$: a nonzero homogeneous binary
determinant has a projective root, and if the determinant
vanishes identically every nonzero member has rank one.
Contracting party $1$ with $\alpha$ gives, for every $z$,
\[
 R\,C(z)\,S=(\alpha_0z_0+\alpha_1z_1)I_2,
 \quad
 R_{b,:}=v^{\mathsf T}B_b,\quad
 S_{:,d}=D_d u.
\]
Since $\alpha\ne0$, one choice of $z$ makes the right
side invertible. Hence $R,S$ are invertible, and the displayed
identity forces
$C(z)=(\alpha\cdot z)R^{-1}S^{-1}$ for all $z$.
Thus $C_0,C_1$ are proportional, contradicting conciseness
at party $3$. This proves the exact obstruction in every
characteristic.\hfill$\square$

The example also rules out a general relative-closedness
principle based only on invertibility of all adjacent
two-site blocked maps. The normalized family $P+tU$ tends
to $P$, whose four adjacent-pair flattenings all have rank
four. Their determinants remain nonzero over $F((t))$.
Each associated blocked source-to-output map is square of
size four and must consequently be invertible. Nevertheless
the limit is not an exact cycle restriction. This statement
does not rule out a stronger principle in characteristic
zero or for the specific balanced $K_{3,3}$ target.

\emph{Independent algebraic-geometry certificate in characteristic two.}
The geometric closure of the binary cycle image over
$\overline{\mathbb F}_2$ is exactly $V(I_6)$. Here are
complete finite certificates and the structural argument.
The polynomial $I_6$ is invariant under $\mathrm{SL}_2^4$.
An irreducible factor of an invariant is fixed up to scalar
by this connected group; that scalar is an algebraic
character. Such a character is trivial, since each
$\mathrm{SL}_2$ is generated by its additive unipotent
subgroups and an additive group has no nontrivial
algebraic character into the multiplicative group.
Every factor is therefore invariant. Diagonal torus weights
exclude invariants of odd degree. A nontrivial factorization
of a sextic would consequently contain an invariant
quadratic factor.

In degree two the torus weights allow only the eight
complementary-coordinate monomials. Polarization is
injective on their span: in characteristic two a quadratic
form with zero polarization is a sum of squares, and these
squares have incompatible torus weights. The space of
invariant bilinear forms is one-dimensional, the tensor
product of the four one-dimensional alternating-form
spaces for the standard $\mathrm{SL}_2$ modules. Hence
the only invariant quadratic, up to scalar, is
\[
 H=\sum_{i=0}^7 t_i t_{15-i},
\]
where $t_i$ are the tensor coordinates in binary
lexicographic order. Direct unipotent substitution verifies
that $H$ is invariant; the preserved eight-variable
coefficient system has rank seven over $\mathbb F_2$.
The explicit binary point
\[
 (t_0,\ldots,t_{15})
 =(0,1,1,1,1,1,1,1,0,1,0,0,1,0,0,1)
\]
has $H=0$ and $I_6=1$ in $\mathbb F_2$. Thus $H$ does
not divide $I_6$, proving that the sextic is irreducible
over $\overline{\mathbb F}_2$.

Finally, the preserved certificate supplies a point of the
32-parameter cycle parametrization and a $15\times15$
Jacobian minor whose determinant over the integers is
$-1$. The image closure therefore has dimension at least
15 in characteristic two. It is irreducible and contained
in the irreducible 15-dimensional hypersurface $V(I_6)$,
so equality follows. This independently confirms the
geometric-border assertion; the explicit polynomial maps
above provide the stronger same-field degeneration.

The coordinate point, selected Jacobian rows and columns,
full minor, determinant, unipotent quadratic-invariance
linear system, and all sixteen coefficients of (EX36.2)
are recorded by \path{programs/ex_cycle_sextic.py} and its
input/output JSON files. The explicit degeneration and
exact obstruction were independently re-proved by both
the parent and the shortening agent. Neither this example
nor its characteristic distinction improves the asymptotic
balanced-cover gate: its tensor-square routing conversion
was already available over every field in EX-34.

The parent's independent reconstruction of EX-36 is preserved as
\path{programs/rt_verify_cycle_characteristic.py} and
\path{outputs/rt_cycle_characteristic_verification.json}.
It recomputes the sextic, unipotent constraints, witnesses,
Jacobian minor and all sixteen integer coefficients.

\paragraph{EX-37: balanced graph conversions are excluded in border by overlapping orbit tangents.}
Let $G_r(q)$ denote the graph tensor of $K_{r,r}$ and
$H_r(q)=K_r(q)\otimes K_r(q)$ the two disjoint cliques on
the same labelled parts. For every $r\ge4$, every $q\ge2$
and every field, $H_r(q)$ is not a border restriction of
$G_r(q)$, even after extending the field to its algebraic
closure. The same conclusion holds for $r=3,q=2^N$, $N\ge1$, in
characteristic zero and characteristic two.

\emph{An elementary orbit dimension and separability lemma.}
Work over an algebraically closed field. For any simple
graph $J$ with $v$ vertices, $e$ edges of dimension $q$, and
vertex dimensions $n_i=q^{\deg(i)}$, let $T_J$ be its graph
tensor. Reorder the ambient tensor factors by edges, and
write each edge space of dimension $q^2$ as
$\langle I_e\rangle\oplus W_e$, with
$\dim W_e=q^2-1$. One can take $W_e$ to be the kernel
of the coordinate-$(0,0)$ functional; no trace normalization
is used. The image of the infinitesimal action of
$\operatorname{End}(F^{n_i})$ is the full tensor product
on edges incident to $i$, with fixed $I_e$ on all other
edges. Its dimension is $n_i^2$.

In the sum of these coordinate subspaces, the empty edge
subset is counted $v$ times, every singleton edge subset
twice, and every subset of at least two edges at most once.
Consequently the differential of the physical group action
has rank
\[
 d_J=\sum_i n_i^2-e(q^2-1)-v+1.
 \tag{EX37.1}
\]
There is an explicit stabilizer family of complementary
dimension $e(q^2-1)+v-1$: assign an invertible matrix $g_e$
to each edge, acting by $g_e$ and $g_e^{-\mathsf T}$ at
its endpoints, and multiply the vertex operators by
scalars $\lambda_i$ with $\prod_i\lambda_i=1$.
Normalize one nonzero entry of every $g_e$ to one.
The resulting parameter family is injective, since equality
of tensor products of nonzero operator factors determines
each factor up to scalar, and the normalizations fix those
scalars. Thus the orbit has dimension at most $d_J$.
Its dimension is at least the differential rank $d_J$,
so equality holds. Algebraic group orbits over an
algebraically closed field are smooth; hence the
differential image is their entire tangent space.
This proves the needed orbit dimension and separability
in every characteristic, also for disconnected graphs.

\emph{The common tangent space.}
Let $X$ be the closure, in the target format, of the
restriction image of $G_r(q)$. It is irreducible.
For a permutation $\sigma\in S_r$, group $L_i$ with
$R_{\sigma(i)}$. Contracting the internal wire in each
group leaves the complete graph on $r$ grouped parties,
with edge dimension $q^2$. Its boundary dimension at a
group and the grouped output dimension are both
$q^{2(r-1)}$. The resulting blocked local maps are
generically invertible: projecting each matching wire to
one fixed basis label gives an explicit invertible choice.
It follows that $X$ is contained in the closure of this
grouped complete-graph orbit, for every $\sigma$.

The target $H_r(q)$ belongs to each of these grouped
orbits: its left and right clique edges combine into
$q^2$-dimensional wires between the grouped parties.
Write $T_i$ for the infinitesimal physical-action image
at vertex $i$ of one clique. By the preceding lemma,
the tangent space of the grouped orbit at $H_r(q)$ is
\[
 \mathcal T_\sigma=\sum_{i=1}^r T_i\otimes T_{\sigma(i)}.
 \tag{EX37.2}
\]
If $H_r(q)\in X$, then
$T_{H_r(q)}X\subseteq\bigcap_\sigma\mathcal T_\sigma$.

Use the common edge-space direct-sum decomposition on
each clique. For an edge subset $S$, let $A(S)$ be the
vertices incident to every edge in $S$. Its size is $r$
for the empty subset, two for a singleton, one for a
nonempty star containing at least two edges, and zero
otherwise. The component indexed by a left subset $S$
and right subset $U$ belongs to $\mathcal T_\sigma$
exactly when $\sigma(A(S))\cap A(U)\ne\varnothing$.
This holds for every permutation if and only if
\[
 |A(S)|+|A(U)|>r.
 \tag{EX37.3}
\]
The equivalence follows by matching $A(S)$ into the
complement of $A(U)$ whenever their sizes sum to at
most $r$.

For $r\ge4$, two nonempty subsets have availability
sizes at most two, so cannot satisfy (EX37.3).
The surviving components are precisely those of the
ordinary physical orbit of $H_r(q)$, of dimension
\[
 d=2r q^{2(r-1)}-r(r-1)(q^2-1)-2r+1.
 \tag{EX37.4}
\]
Thus membership of $H_r(q)$ in $X$ would imply
$\dim X\le\dim T_{H_r(q)}X\le d$.

\emph{A dimension contradiction without a generic-rank assumption.}
For each perfect matching $M$ of $K_{r,r}$, projecting
its wires to fixed labels produces the graph tensor of
$J_M=K_{r,r}\setminus M$, in the same output format.
The physical orbit of $J_M$ lies in $X$ and has dimension
$d$ by (EX37.1), because $J_M$ and $H_r(q)$ have the same
degrees and edge count. Two distinct matchings give
distinct physical orbits: the flattening at a pair
$\{L_i,R_j\}$ has rank $q^{2(r-1)}$ if that edge was
deleted and $q^{2(r-2)}$ otherwise. An irreducible variety
of dimension $d$ cannot contain two disjoint locally
closed orbits of dimension $d$, since each would contain
a nonempty open subset. Hence $\dim X\ge d+1$, a
contradiction. This proves the assertion for $r\ge4$.

\emph{The binary three-by-three case.}
For $r=3,q=2$, condition (EX37.3) permits one additional
kind of component: a singleton edge on each side.
The physical orbit dimension is $73$, and the nine
edge pairs contribute $9(2^2-1)^2=81$. Therefore
\[
 \dim\bigcap_\sigma\mathcal T_\sigma=154.
 \tag{EX37.5}
\]
The exact certificate
\path{programs/ex_balanced_tangent.py}, with fixed
binary parameters \path{inputs/ex_balanced_tangent_parameters.json},
computes the $4096\times192$ Jacobian of the restriction
parametrization at the displayed point. It preserves a
$155\times155$ minor that is nonsingular modulo both
two and 101. Thus $\dim X\ge155$ in characteristic two
and characteristic zero, contradicting (EX37.5).
The full Jacobian ranks modulo both primes are 160,
but the retained minor already suffices. Entries, selected
rows and columns, determinant residues, and a separate
combinatorial tangent count are saved in
\path{outputs/ex_balanced_tangent_results.json}.
An independent reconstruction by the root agent, using
\path{programs/rt_verify_balanced_tangent.py}, reproduces
the complete integer Jacobian hash, both full ranks, the
minor determinants $1$ modulo two and $95$ modulo 101,
and every site and five-site rank used below; see
\path{outputs/rt_balanced_tangent_verification.json}.

The same finite certificate extends to every binary tensor
power. Let $S_i$ be the image of the derivative with respect
to the $i$th local map at the certified binary point.
Over each of $\mathbb F_2$ and $\mathbb F_{101}$, the
program verifies
\[
 \dim S_i=32,\qquad
 \dim\sum_iS_i=160,\qquad
 \dim\sum_{j\ne i}S_j=138
\]
for all six sites. The quotient $P_i$ by the sum of the
other five subspaces therefore has rank 22 on $S_i$.
At the $N$th tensor-product point, the derivative image
at site $i$ is exactly $S_i^{\otimes N}$: tensor products
of variations of the original local map span all variations
of the enlarged local map. The direct sum of the maps
$P_i\otimes\operatorname{id}$ annihilates every other
site's component and has image dimension
\[
 6\cdot22\cdot32^{N-1}=132\cdot32^{N-1}.
 \tag{EX37.6}
\]
Thus the enlarged restriction image has at least that
dimension in characteristic two. Reduction of its integer
Jacobian modulo 101 proves the same lower bound in
characteristic zero. For $q=2^N$ and $N\ge2$,
\[
 132\cdot32^{N-1}=\frac{33}{8}q^5
 >15q^4-24q^2+10,
\]
which is the common tangent dimension obtained from
(EX37.3) for $r=3$. Indeed $q\ge4$ gives
$(33/8)q\ge33/2>15$. This completes the tensor-power
extension without constructing a large new Jacobian.
\hfill$\square$

For $r\ge4$ this excludes a one-term border conversion
at every equal-size tensor power, since edge dimensions
become $q^N$. It does not exclude covers with
subexponentially many terms or extra source width.
Accordingly, the asymptotic balanced-cover gate in SH-33
remains open; this theorem closes the proposed shortcut
through a single equal-size degeneration.

There is an exact limitation to a first-order extension of
this proof to covers. Each grouped orbit closure $Y_\sigma$
is a cone containing $H_r(q)$ and spans the entire ambient
tensor space. For every $T\in Y_\sigma$, the curve
$H_r(q)+tT$ lies in its two-term secant cone. Therefore
the tangent space of that secant cone at $H_r(q)$ contains
the full linear span of $Y_\sigma$, and is the entire
ambient space. Intersecting these first-order secant
tangents is already vacuous at two terms. A quantitative
selector or multi-term obstruction requires additional,
higher-order information.

\subsection{What classical network-coding identities can transfer}
\label{sec:ex-coding-transfer}

\paragraph{EX38: exact transfer and a rank-metric limitation.}
The recent preprint of Zhang, Li, and Li (2026),
\emph{Network Coding Can Beat Routing in Undirected Multiple-Unicast
Networks}, gives a particularly concrete lead: its scalar-instance
definition (Definition III.2, printed p.~4) requires a deterministic
causal code that uses each physical edge exactly once. The reversible
construction on pp.~5--6 is described by integer coefficient identities.
The argument below states precisely what follows from such a code;
it does not assume that the entire new preprint or its Lean development
has been independently verified.

\begin{proposition}[EX38: a deterministic code gives an all-field restriction]
Let a finite deterministic causal protocol over an alphabet of size $q$
deliver $B_i$ independent symbols from $s_i$ to $t_i$ exactly.
Let $D_e$ be the total number of scheduled uses of undirected edge $e$,
counting both directions. Over every scalar field $F$, the product of
source wires of dimensions $q^{D_e}$ restricts to the product of demand
wires of dimensions $q^{B_i}$, with unit physical factors at internal
vertices.
\end{proposition}
\begin{proof}
Give an edge one virtual index for its entire scheduled transcript.
At each vertex use the zero--one table that enforces all its local
encoding rules and decoding rules, with its external source and sink
indices exposed. For a fixed tuple of source inputs, causal induction
determines exactly one global transcript. Exact decoding makes the
exposed sink tuple equal to the corresponding source tuple. Thus the
contracted coefficient is exactly one on the demand-wire support and
zero elsewhere. This is an integer identity with no cancellations,
so it specializes to every $F$. In particular, the alphabet field of
a linear code need not equal the scalar field of the resulting tensors.
\end{proof}

This transfer alone cannot violate Euclidean wire length. Every
one-term restriction between wire products obeys every cut-rank
inequality. Nonnegative combinations and limits give the same inequality
for every $L_1$ pseudometric, including Euclidean distance by Crofton's
representation. A coding advantage for a shortest-path metric is
therefore not itself a flatness counterexample or a strict shortening.

There is a stronger obstruction to using linear codes to violate the
specific square-root rank metrics in the contraction gate.

\begin{proposition}[EX38-RANK: concave rank distances respect linear codes]
Suppose the preceding protocol is linear over a commutative field $K$.
For matrices $P_v\in\operatorname{Mat}_n(K')$ assigned to its vertices,
where $K'/K$ is any field extension, and every nondecreasing concave
function $\varphi:\{0,1,\ldots,n\}\to\mathbb R_{\ge0}$ with
$\varphi(0)=0$, one has
\[
 \sum_i B_i\varphi\!\left(\operatorname{rank}(P_{s_i}-P_{t_i})\right)
 \le
 \sum_e D_e\varphi\!\left(\operatorname{rank}(P_u-P_v)\right),
 \qquad e=\{u,v\}.
 \tag{EX38.1}
\]
In particular, this holds for $\varphi(j)=\sqrt j$.
\end{proposition}
\begin{proof}
First take $\varphi(j)=j$. Replicate each scalar symbol of the code
as a vector in $(K')^n$, with every scalar coefficient acting on all
coordinates. Let $x$ be the tuple of original source vectors and
$x_a$ the vector transmitted at event $a:u\to v$ in the original
execution. Start a new execution with source vector $P_{s_i}x_i$
for each coordinate of message $i$, and add the error
\[
 (P_v-P_u)x_a
\]
when event $a$ arrives at $v$. By induction through the causal event
order, every local stored vector at vertex $v$ is $P_v$ times its
original value: the claim holds for its own inputs, scalar local
operations commute with $P_v$, and the injected error converts
$P_ux_a$ to $P_vx_a$. The decoded vector is consequently $P_{t_i}x_i$.
Without the errors, linearity and exact decoding would instead give
$P_{s_i}x_i$.

Let $E$ send $x$ to the list of these injected errors, and let $D$ be
the linear transfer from independent event errors to decoded outputs
when source inputs are zero. We have the exact factorization
\[
 \bigoplus_i I_{B_i}\otimes(P_{t_i}-P_{s_i})=DE.
 \tag{EX38.2}
\]
The map $E$ factors through
$\bigoplus_{a:u\to v}\operatorname{im}(P_v-P_u)$.
Taking ranks in (EX38.2) proves (EX38.1) for $\varphi(j)=j$,
with one summand per scalar event and hence $D_e$ summands per
physical edge. No sign, characteristic, or positivity assumption
is used.

For $1\le k\le n$, pass if necessary to an infinite extension and
choose common matrices $L\in\operatorname{Mat}_{k\times n}$ and
$R\in\operatorname{Mat}_{n\times k}$ generically. For every one of
the finitely many edge and demand differences $Q$, simultaneously,
\[
 \operatorname{rank}(LQR)=\min\{k,\operatorname{rank}Q\}.
\]
For each fixed $Q$ the required nonzero minor exists for some $L,R$;
the product of the finitely many corresponding nonzero polynomials
is still nonzero. Apply the proved rank inequality to the vertex
matrices $LP_vR$. This proves (EX38.1) for every truncated rank
$\min\{k,j\}$.

Write $\Delta_k=\varphi(k)-\varphi(k-1)$ and
$\Delta_{n+1}=0$. Monotonicity and concavity say
$\Delta_k-\Delta_{k+1}\ge0$, and telescoping gives
\[
 \varphi(j)=\sum_{k=1}^n
       (\Delta_k-\Delta_{k+1})\min\{k,j\}.
\]
The nonnegative linear combination of the truncated-rank
inequalities proves the assertion.
\end{proof}

The root and shortening agents independently re-proved the error
factorization, simultaneous generic compression, and discrete concavity
step. Thus any integer-coefficient code, interpreted over
$\mathbb F_2$, already respects $\sqrt{\operatorname{rank}(A-B)}$
on every finite set of binary matrices. Its graph-tensor identity,
although valid over both $\mathbb C$ and $\mathbb F_2$, cannot provide
a violated linear inequality on that rank-building configuration.
For a code available only in another characteristic, the conclusion
applies to matrices over that coding field. In the particular
eight-point configuration of SH-20, Proposition
\ref{rt:universal-rank-witness} supplies an explicit integral
realization with exactly the same rank distances over every field.
Consequently even a code defined only over $\mathbb F_9$, or any
other coding field, cannot make its direct wire inequality exclude
that configuration. No cross-characteristic realization of the full
metric on $\operatorname{Mat}_n(\mathbb F_2)$ is asserted. The
proposition is a limitation of causal linear coding, not of arbitrary
tensor restrictions, coherent covers, or nonlinear codes.

There is also a classical obstruction to a direct complete-bipartite
specialization. Jain, Vazirani, Yeung, and Yuval (2005), Theorem~4
(author PDF pp.~4--5), imply that when every demand has both endpoints
in one side of a bipartite network, the total entropy of the independent
messages is at most the total capacity in either direction. Adding
the two directional inequalities gives
$2\sum_iB_i\le\sum_eD_e$. On $K_{r,r}$ with unit shared edge
capacity and all $r(r-1)$ within-side demands, the common coding rate
is therefore at most $r/[2(r-1)]$. The hypotheses include classical
causality and independent messages; this bound is not asserted for
an arbitrary GHZ-assisted tensor restriction.

\subsection{An adaptive-topology countermodel to a broader exponent barrier}
\label{sec:ex-adaptive-countermodel}

\paragraph{EX39: singleton groups with an existential hub tree.}
The fixed-topology and fixed-label quantifiers in EX-29 cannot simply
be dropped. In the broader geometric model that allows a connected hub
graph, or the assignment of groups to its labels, to be chosen after
the labelled equal partition is known, polynomially many hubs permit
arbitrarily large fixed powers of $B$. This is a countermodel to that
broader geometric upper-bound statement, not a circuit construction.

Fix the positive selector-debit constant $a$ and integers
$k=2^t\ge4$, $1\le B\le k/2$. Put $n=k/2$. On circuit vertices
$L_0,\ldots,L_{n-1},R_0,\ldots,R_{n-1}$ take the simple
$B$-regular bipartite graph
\[
 L_iR_{i+s\bmod n},\qquad 0\le i<n,\quad0\le s<B,
\]
with unit integer widths. Thus its maximum incident load is $B$ and
its total wire width is
\[
 W=nB=kB/2.
\]
Place every terminal belonging to an $L_i$ at $+e_1$ and every
terminal belonging to an $R_i$ at $-e_1$. All target segments pass
through the origin and have length two, so the target cost is $2W$.
The endpoint set has exact binary-certificate cost $C=2$: the
identity map on its two points is a certificate, and Jensen's
inequality forces every certificate to cost at least their distance.

Take $r=k$ hubs, so each labelled equal group is a singleton.
For any such labelled partition, put the hub assigned to $L_i$ at
$+e_1$ and the hub assigned to $R_i$ at $-e_1$. Choose one left
hub and one right hub as roots; connect every other left hub to
the left root, every other right hub to the right root, and connect
the two roots. This is a tree on all $k$ labelled hubs. Give every
tree edge width
\[
 h=W/2=kB/4.
\]
For the sequence below this is an integer. Every fan has zero
length. The $k-2$ same-side tree edges also have zero length;
the remaining edge has length two. With selector debit exactly
$aktC=2akt$, the total cost is
\[
 2h+2akt=W+2akt<2W
 \quad\Longleftrightarrow\quad B>4at.
 \tag{EX39.1}
\]
Thus this broader model genuinely admits strict improvement for
every labelled singleton partition, with all bridge widths equal
to $h$.

For any fixed real $M\ge1$, choose
$B=2^j$ and $k=2^{\lceil Mj\rceil+1}$. Then $B\le k/2$,
$k=\Theta(B^M)$, $t=O_M(\log B)$, and (EX39.1) holds for all
sufficiently large $j$. Yet
\[
             h=kB/4=\Theta(B^{M+1}).
 \tag{EX39.2}
\]
Even $M=1$ gives $h=\Theta(B^2)$. Letting the allowed polynomial
bound on the number of hubs vary gives no universal finite power
supremum in this broad abstract geometry.

The tree was chosen using the left/right types of the singleton
groups. Equivalently, one could fix the displayed two-star tree
and permit a subsequent type-dependent permutation of the groups
onto its labels. Neither operation is part of EX-29's fixed labelled
model. More substantially, neither supplied manuscript has been
shown to produce this existential choice of hub topology with the
stated bridge width from an arbitrary circuit. The construction
therefore does not establish a stronger direct-case circuit lower
bound, contradict EX-29, or determine the unrestricted supremum in
Q3. It identifies a necessary quantifier in any proposed universal
geometric limitation.

\emph{The countermodel also has an exact actual selector price.}
This limitation is not caused by an unknown normalization of its
two-point geometry. For an explicit family, take $B$ through powers
of two, fix an integer $M\ge1$, and set $k=2B^M$. A simple
$B$-regular bipartite graph exists with $k/2$ vertices on each side,
for example by taking $B$ cyclic matchings between the two sides.
Its total width is $W=kB/2$, with unit graph-edge widths. Place the
terminal at each left vertex at $+1$ and each right terminal at
$-1$. The target has cost $2W=kB$.

Use singleton groups, so $r=k$ and $t=\log_2k$. After each labelled
singleton partition, choose or relabel the two-star hub tree by the
left/right types as above: every left hub is at $+1$, every right
hub at $-1$, same-type hubs form their respective stars, and one
edge joins the two star centers. All fan lengths and all same-type
bridge lengths vanish. If every tree edge has width
$h=W/2=kB/4$, the total bridge cost is $2h=W$.

On all physical endpoint and hub labels, take the ordinary spectral
point given by flattening rank across the division between positions
$+1$ and $-1$. Its binary wire metric is zero within a position and
one between the two positions. Scale by $c=2$. This realizes every
Euclidean distance in the construction exactly. The partial binary
GHZ on all endpoint labels has flattening rank two, so its actual
scaled exponent is $G=2$. Charging exactly the postulated debit
$aktG=2ak\log_2k$ gives total abstract cost
\[
                      W+2ak\log_2k<2W
             \quad\text{whenever}\quad B>4a\log_2k.
\]
For every fixed $M$ this holds for all sufficiently large $B$, while
$h=\Theta(B^{M+1})$. Thus an exact actual spectral price does not
repair the geometric barrier under this adaptive-topology model.

There is, however, an elementary obstruction to interpreting the
postulated debit as the bit count of a genuine cover. The target
tensor has flattening rank $2^W$ across all left physical parties
versus all right physical parties. Every network just described
has at most rank $2^h$ across that output cut: cutting its unique
positive-length bridge separates the two types, and local maps and
internal contractions cannot increase this rank. Therefore any
$A$-term cover using these layouts must satisfy
\[
                         A\,2^h\ge2^W,
 \qquad \log_2A\ge W-h=W/2.
\]
In particular, the proposed budget $\log_2A\le ak\log_2k$ is
impossible in the range giving strict abstract improvement. With
the actual required bit count, bridge cost plus selector cost is
at least $2h+G\log_2A\ge2W$. The actual spectral price exists;
the proposed cheap cover does not. No circuit decomposition is
therefore obtained from this adaptive geometric example.


\subsection{A two-hub bound without a common segment origin}
\label{sec:ex-nonconcurrent-twohub}

\paragraph{EX40: equal-length segments with arbitrary centers.}
Let a loopless graph on $k\ge2$ circuit vertices have positive integer
widths $b_e$, total width $W=\sum_e b_e$, and maximum incident load
$B>0$. Each oriented edge $e=ab$ has terminals $p_e$ at $a$ and
$q_e$ at $b$ in a Euclidean space. Assume that all target segments
have one common positive length; after a uniform rescaling write
\[
 x_e=p_e-q_e,\qquad \|x_e\|=1.
\]
Their centers are arbitrary. The target cost is $W$.

Use precisely the binary pointwise-averaging certificate defined
before EX-16, now on the whole endpoint set
$P=\{p_e,q_e:e\in E\}$, with cost $C$. For each balanced cut
$k/3\le |S|\le2k/3$, connect all terminals belonging to $S$ to
one hub $z$, and all remaining terminals to a second hub $w$.
Allow both hubs to depend on the cut and to lie anywhere.
The two hubs are joined by one bridge of width at least $h$,
and the method charges selector debit at least $akC$, with
fixed $a>0$. If every balanced cut admits total cost strictly
less than $W$, then
\[
 h=O_a\!\left(B^{4/3}(\log(2+B))^{1/3}\right).
 \tag{EX40.1}
\]
Thus the fixed-power supremum is again $4/3$ in this specified
equal-length, two-hub, charged-certificate model. Attainment
remains conditional on the supplied flat-case reduction used
in EX-18, whose concurrent equal-length construction is included.
No claim about the actual spectral selector exponent is made.

\emph{The reduction that removes segment centers.}
Put $d=z-w$. If an edge crosses the cut, the triangle inequality
gives
\[
 \|p_e-z\|+\|q_e-w\|\ge\|x_e-d\|;
\]
with the opposite endpoint assignment it gives $\|x_e+d\|$.
If both endpoints belong to the same group, their fan cost is
at least $\|p_e-q_e\|=1$. Consequently the total fan saving
is at most
\[
 S_S(d)=\sum_{e\text{ crosses }S}
        b_e\{1-\|\sigma_e x_e-d\|\},
 \tag{EX40.2}
\]
where $\sigma_e=1$ if the oriented initial vertex belongs to
$S$, and $\sigma_e=-1$ otherwise. This is an upper bound on
the saving of the original layout, not an asserted realization
of the relaxed expression. Since all points $\sigma_e x_e$
lie in the unit ball, radial projection of $d$ to that ball
increases (EX40.2) and decreases the bridge cost. Thus its
supremum may be taken over $\|d\|\le1$. Projection onto the
span of the $x_e$ has the same property.

Write
\[
 v_a=\sum_{e\ni a}\pm b_ex_e,\quad
 A=\sum_e b_ex_ex_e^{\mathsf T},\quad Q=WI-A.
\]
The signs in $v_a$ are induced by the oriented segment
differences, so $\sum_av_a=0$ and
$R_S=\sum_{a\in S}v_a=\sum_{e\text{ crosses }S}\sigma_eb_ex_e$.
Let $D\ge1$ be the dimension of the direction span and put
$\ell=\ln(64(D+1))$. We prove that some balanced cut satisfies
\[
                  \sup_d S_S(d)\le20000B\ell.
 \tag{EX40.3}
\]

\emph{A paired random cut and its crossed-edge curvature.}
Let $u$ be a unit top eigenvector of $A$ and write
$v_a=\alpha_au+w_a$. If $k$ is odd add one zero-load dummy
vertex for this pairing argument. Sort the $\alpha_a$,
pair consecutive vertices, and give the two vertices in pair
$i$ opposite signs $\epsilon_i,-\epsilon_i$, independently
over the pairs. The positive-sign set on the original vertices
is always balanced. As in EX-19,
\[
 |R_u|\le B,\qquad
 \mathbb E\|R_\perp\|^2\le BW.
 \tag{EX40.4}
\]
Indeed the sum of adjacent gaps is at most $2B$, while
weighted Cauchy--Schwarz at the vertices gives
$\sum_a\|w_a\|^2\le2BW$. Hence
$\|R_\perp\|^2\le4BW$ with probability at least $3/4$.

Set $M_e=b_e(I-x_ex_e^{\mathsf T})$ and
\[
 Q_{\rm cut}=\sum_{e\text{ crosses }S}M_e.
\]
An edge whose endpoints form one paired pair always crosses;
every other edge crosses with probability one half. Therefore
\[
 \mathbb E Q_{\rm cut}
 =\tfrac12 Q+\tfrac12\sum_{e\text{ inside a paired pair}}M_e
 \succeq Q/2.
 \tag{EX40.5}
\]
This is a quadratic, rather than linear, function of the
independent pair signs. Nevertheless it has a useful matrix
martingale bound.

Give each vertex a fixed internal pair sign $\tau_a\in\{1,-1\}$,
so its cut sign is $\tau_a\epsilon_i$ if it belongs to pair $i$.
Reveal pair signs in order. The Doob martingale for
$Q_{\rm cut}-\mathbb E Q_{\rm cut}$ has increment
\[
 D_i=-\frac{\epsilon_i}{2}
       \sum_{\substack{e=ab:\\a\in i,\ b\in j<i}}
          \epsilon_j\tau_a\tau_b M_e.
 \tag{EX40.6}
\]
The total width in the sum is at most $2B$. For any vector $y$,
weighted Cauchy--Schwarz and
$(I-x_ex_e^{\mathsf T})^2=I-x_ex_e^{\mathsf T}$ give
\[
 \|D_i y\|^2
 \le \frac14\left(\sum_{e\text{ in }(i,<i)}b_e\right)
             \sum_{e\text{ in }(i,<i)}y^{\mathsf T}M_ey
 \le\frac B2\sum_{e\text{ in }(i,<i)}y^{\mathsf T}M_ey.
\]
Each edge between different pairs occurs in exactly one such
sum. Since the square of the increment is measurable before
$\epsilon_i$ is revealed, this proves
\[
 \|D_i\|\le B,\qquad
 \sum_i\mathbb E[D_i^2\mid\epsilon_1,\ldots,\epsilon_{i-1}]
 \preceq (B/2)Q.
 \tag{EX40.7}
\]
The same proof applies after compression to $u^\perp$.
If $q=\lambda_{\min}(Q)>0$, apply it instead to
$M'_e=Q^{-1/2}M_eQ^{-1/2}$. Now
$\|M'_e\|\le b_e/q$, so
\[
 \|D'_i\|\le B/q,\qquad
 \sum_i\mathbb E[(D'_i)^2\mid\epsilon_1,\ldots,\epsilon_{i-1}]
 \preceq B I/(2q).
 \tag{EX40.8}
\]
For example, the second inequality follows from the weighted
Cauchy--Schwarz calculation with
$N_e=M'_e/b_e$, using $N_e^2\preceq N_e/q$ and
$\sum_eM'_e=I$.

Tropp's matrix Freedman inequality (2011, Theorem~1.2,
published p.~264), applied also to the negatives, gives the
two-sided tail
$2D\exp[-t^2/(2(V+Rt/3))]$ from variance bound $V$
and increment bound $R$. One may use the weaker parameters
$V=2BW,R=2B$ in (EX40.7), and
$V=2B/q,R=2B/q$ in (EX40.8).
For $W\ge4096B\ell$, deviation threshold $W/8$ on
$u^\perp$, together with (EX40.5), gives
\[
 Q_{\rm cut}|_{u^\perp}\succeq WI/8
 \tag{EX40.9}
\]
except on an event of probability at most $1/16$.
Here every eigenvalue of $A$ other than its largest is at most
$W/2$, because $\operatorname{tr}A=W$.
If also $q\ge4096B\ell$, whitened deviation threshold $1/4$
gives
\[
                         Q_{\rm cut}\succeq Q/4
 \tag{EX40.10}
\]
with the same failure bound. These events intersect the
event in (EX40.4) of probability at least $3/4$.
When a condition is not met, its corresponding concentration
event is simply not required. The one-dimensional case has
zero transverse space and can be handled directly.

\emph{Uniform control of the relaxed saving.}
For $\|d\|\le1$, the exact norm identity used in EX-17 gives
\[
 S_S(d)\le R_S\cdot d-\tfrac14d^{\mathsf T}Q_{\rm cut}d.
 \tag{EX40.11}
\]
If $W<4096B\ell$, the trivial bound $S_S(d)\le W$
already proves (EX40.3). Otherwise choose a cut satisfying
(EX40.4) and (EX40.9), and also (EX40.10) when applicable.

If $q\ge4096B\ell$, maximize the quadratic bound without its
unit-ball constraint to obtain
\[
 \sup_d S_S(d)
 \le4R_S^{\mathsf T}Q^{-1}R_S
 \le4\left(\frac{B^2}{q}
             +\frac{2\|R_\perp\|^2}{W}\right)<33B.
\]
In the remaining case $q<4096B\ell$, write
$Q_{\rm cut}$ in the decomposition $\mathbb Ru\oplus u^\perp$.
Its top entry $a$ is at most $q$, since $0\preceq Q_{\rm cut}\preceq Q$;
the squared norm of its off-diagonal block is at most $aW$,
by positivity and $\|Q_{\rm cut}\|\le W$.
Using (EX40.9) and Young's inequality gives
\[
 d^{\mathsf T}Q_{\rm cut}d
 \ge\frac W{16}\|d_\perp\|^2-15q\,d_u^2.
\]
Since $\|d\|\le1$, (EX40.11) implies
\[
 \sup_d S_S(d)
 \le B+\frac{16\|R_\perp\|^2}{W}+\frac{15q}{4}
 \le65B+15360B\ell<20000B\ell.
\]
If $D=1$, (EX40.4) alone gives $S_S(d)\le B$.
This proves (EX40.3) for arbitrary segment centers.

\emph{Selector accounting and the exponent.}
Strict improvement for the cut in (EX40.3) requires
$akC<20000B\ell$. For any cut,
$\|\sigma_ex_e-d\|\ge1-\sigma_ex_e\cdot d$ in (EX40.2),
so a bridge of width $h$ permits strict improvement only if
$h<\|R_S\|$. The binary certificate contains each
half-difference $x_e/2$ in its zonotope. Applying EX-16 with
the vertex-vector bound $v_a/(2B)\in Z$ therefore gives,
for $k\ge43$, a possibly different balanced cut with
$\|R_S\|\le5\sqrt2 BC$. Its random-pair bound also gives
a balanced cut with $\|R_S\|\le B\sqrt{k/2}$.
It is legitimate to use different cuts, because the hypothesis
requires improvement for every balanced cut with the same
certificate charge. Consequently
\[
 h\le B\min\{\sqrt{k/2},\,5\sqrt2 C\},
 \qquad
 kC<20000B\ell/a.
 \tag{EX40.12}
\]
For bounded $k$ the first inequality already gives $h=O(B)$.

Jensen's inequality gives $C\ge1$. Integer widths give
$D\le |E|\le W\le kB/2$. Thus (EX40.12) first implies
$k=O_a(B\log(2+B))$ by the same logarithmic bootstrap as
EX-19, and hence $\ell=O_a(\log(2+B))$. Substituting the
second inequality into the first and optimizing over $k$
gives (EX40.1), with the balancing scale
$k\asymp(B\log(2+B))^{2/3}$.
The proof does not require a common midpoint, common supporting
line intersection, or bounds on the original segment centers.
It does require equal target lengths and exactly the stated
two-hub and prescribed selector-debit model.


\paragraph{EX41: unequal lengths, an aspect-ratio bound, and a $3/2$ ceiling.}
Keep the exactly-two-hub and charged-certificate assumptions of EX40,
but allow arbitrary positive target lengths $l_e=\|p_e-q_e\|$ and
arbitrary segment centers. Put
\[
 L=\max_e l_e,\qquad \delta=\min_e l_e,\qquad \kappa=L/\delta.
\]
Then
\[
 h=O_a\!\left(
 B\min\left\{
       (\kappa B\log(2+B))^{1/3},
       (B\log(2+B))^{1/2}
       \right\}\right).
 \tag{EX41.1}
\]
In particular, arbitrary length spread has a fixed-power upper ceiling
$3/2$ in this model; bounded spread has supremum $4/3$ with the same
conditional attainment as EX40. If $\kappa=B^{o(1)}$, no fixed power
strictly larger than $4/3$ is possible. A $4/3$ bound for completely
unrestricted length spread is not proved by this argument.

\emph{Proof.}
Use the unit directions $x_e=(p_e-q_e)/l_e$ in the definitions
of $v_a,R_S,A,Q$ and in the paired random cut of EX40.
The spring relaxation is now
\[
 S_S(d)=\sum_{e\text{ crosses }S}
 b_e\{l_e-\|\sigma_e l_ex_e-d\|\}.
\]
Its supremum occurs with $\|d\|\le L$, because its spring points
are in that ball. For every $0<l_e\le L$ and $\|d\|\le L$,
the same denominator identity gives
\[
 \|l_ex_e-d\|
 \ge l_e-x_e\cdot d+
      \frac{\|d\|^2-(x_e\cdot d)^2}{4L}.
\]
The zero-denominator case follows directly, as in EX17.
After writing $d=Lz$, the resulting saving bound is
\[
 S_S(Lz)\le L\{R_S\cdot z-z^{\mathsf T}Q_{\rm cut}z/4\},
 \qquad \|z\|\le1.
\]
Thus the entire EX40 random-cut and curvature proof applies
without change and gives some balanced cut with free saving at
most $20000BL\ell$, where $\ell=\ln(64(D+1))$.
Selector payment consequently requires
\[
                  k(C/L)<20000B\ell/a.
 \tag{EX41.2}
\]
Since $C\ge L$, positive integer widths and
$D\le W\le kB/2$ again yield
$k=O_a(B\log(2+B))$ and $\ell=O_a(\log(2+B))$.
The balanced random-resultant bound therefore gives the second
term of (EX41.1), with no dependence on $\delta$.

For the first term, the certificate zonotope contains
$l_ex_e/2$ for every edge. Hence it contains $\delta x_e/2$,
and EX16 with this rescaling gives
$h\le5\sqrt2 BC/\delta$ for $k\ge43$.
Combine this with $h\le B\sqrt{k/2}$ and (EX41.2), and optimize
over $k$ exactly as in EX40. The result is
$h=O_a(B(\kappa B\log(2+B))^{1/3})$.
Bounded $k$ is handled by the random-resultant bound alone.
Zero-length target segments, if allowed, may be omitted in this
upper-bound relaxation because their direct cost is zero and
their fan cost is nonnegative; $\delta$ then denotes the minimum
remaining positive length.

The multiscale step not resolved by the preceding argument is specific.
EX23 balanced
several principal components together with \emph{linear} vertex
curvature matrices. Here the crossed-edge curvature is a
quadratic function of the cut signs. One cannot replace the
independent paired signs by a Banaszczyk signing and retain the
matrix martingale conclusion without an additional argument.
EX42 supplies such an argument by decoupling the cut into two signings.


\paragraph{EX42: arbitrary lengths and centers, by two-stage cut decoupling.}
The unrestricted length-spread gap in EX41 can be removed while keeping
exactly two hubs. In the model of EX40, allow every target segment
$p_e-q_e=l_ex_e$ to have an arbitrary positive length $l_e$ and arbitrary
center, with $\|x_e\|=1$. Keep positive integer widths, maximum incident
load $B$, the same binary pointwise-averaging certificate of cost $C$,
and selector debit at least $akC$ for every balanced cut. The graph,
endpoint positions, and certificate are fixed before the cut is chosen;
both hub positions may depend arbitrarily on the cut. If every balanced cut
admits strict improvement after paying a bridge of width at least $h$,
then
\[
             h=O_a\bigl(B^{4/3}\log(2+B)\bigr).
 \tag{EX42.1}
\]
Thus the fixed-power supremum is $4/3$ throughout this arbitrary-length,
arbitrary-center, exactly-two-hub model, with attainment still
conditional on the supplied flat-case reduction used in EX18. The
selector debit is an assumption of the model, not a proved lower bound
on an actual spectral GHZ exponent. No additional hubs or topology
choice is covered by this theorem.

The new ingredient is a simultaneous cut lemma. Retain any set of
edges, partitioned into $m\ge1$ bands satisfying
$r_j/2<l_e\le r_j$ in band $j$. Let $D\ge1$ be the dimension of the
span of the retained directions, and define
\[
 \ell=\ln(64m(D+1)),\qquad T=2^{18},\qquad
 L_0=100m,\qquad K_0=2^{23}(m^2+m\ell).
\]
As in EX23, put
\[
 W_j=\sum_{e\in j}b_e,\quad A_j=\sum_{e\in j}b_ex_ex_e^{\mathsf T},
 \quad Q_j=W_jI-A_j,\quad M_e=b_e(I-x_ex_e^{\mathsf T}),
\]
and
$v_{aj}=\sum_{e\ni a,\ e\in j}\pm b_ex_e$ with the fixed endpoint
orientation. Choose a unit top eigenvector $u_j$ of $A_j$, write
$v_{aj}=\alpha_{aj}u_j+w_{aj}$, and set
$q_j=\lambda_{\min}(Q_j)$. There is one balanced cut $S$, provided
$k\ge300m$, for which every band simultaneously satisfies
\begin{align}
 \|R_j\|&\le L_0BC/r_j,\qquad R_j=\sum_{a\in S}v_{aj},
 \tag{EX42.2}\\
 \sup_d\sum_{\substack{e\in j\\ e\text{ crosses }S}}
 b_e\{l_e-\|\sigma_el_ex_e-d\|\}&\le K_0Br_j.
 \tag{EX42.3}
\end{align}
All matrices below act on the retained direction span. Pseudoinverses
and covariance whitenings act only on their ranges; zero blocks are
omitted.

\emph{First choose an auxiliary partition with curvature.}
Choose independent signs $\tau_a$ at the circuit vertices and let
$U,V$ be their two sign classes. This partition need not be balanced
and is not the final cut. Put
\[
 Q_{UV,j}=\sum_{\substack{e\in j\\e\text{ joins }U\text{ to }V}}M_e.
\]
Its expectation is $Q_j/2$. Reveal the vertex signs in a fixed order.
The centered Doob increment is
\[
 D_{a,j}=-\frac{\tau_a}{2}
       \sum_{\substack{e=ab\in j\\b<a}}\tau_bM_e.
\]
Its incident width is at most $B$. The weighted vector
Cauchy--Schwarz argument in EX40 therefore gives
\[
 \|D_{a,j}\|\le B/2,\qquad
 \sum_a\mathbb E[D_{a,j}^2\mid\tau_1,\ldots,\tau_{a-1}]
       \preceq BQ_j/4.
 \tag{EX42.4}
\]
Each edge occurs in exactly one increment. After full $Q_j$ whitening
when $q_j>0$, the corresponding bounds are $B/(2q_j)$ and
$BI/(4q_j)$. The compressed bounds on $u_j^\perp$ hold as well.
Use the matrix Freedman inequality cited in EX40, with its deliberately
weaker parameters there. A union bound over at most $2m$ events shows
that an auxiliary partition exists with
\begin{align}
 Q_{UV,j}|_{u_j^\perp}&\succeq W_jI/8
       &&\text{if }W_j\ge TmB\ell,\tag{EX42.5}\\
 Q_{UV,j}&\succeq Q_j/4
       &&\text{if }q_j\ge TmB\ell.\tag{EX42.6}
\end{align}
Indeed $T m\ge4096$, so the thresholds used in EX40 apply with
$\ell=\ln(64m(D+1))$; each required event has failure probability
at most $1/(16m)$. Fix one such $U,V$.

\emph{Sign the two parts successively.}
The certificate gives the zonotope $Z$ and positive semidefinite
matrix $Q_0$ from EX23, with
$Z\subseteq Q_0^{1/2}\mathbb B_2$ and
$\operatorname{tr}Q_0\le(C/2)^2$.
Since $l_ex_e/2\in Z$ and $l_e>r_j/2$, the following blocks have
norm at most one for every vertex:
\[
 \frac{r_j}{4B}Q_0^{-1/2}v_{aj},\qquad
 \frac{\alpha_{aj}}B,\qquad
 B^{-1/2}(A_j|_{u_j^\perp})^{-1/2}w_{aj},\qquad 1.
 \tag{EX42.7}
\]
The last block is included just once, for the total count. The first
three are included for each band. These norm bounds use only the
original vertex vectors, so they remain valid after restriction to
either $U$ or $V$; neither part's unsigned vector sum must vanish.

First apply Banaszczyk's theorem to the vertices in $U$ with the
blocks (EX42.7), divided by
$\Lambda=5\sqrt{5m+1}$. Use the same scalar, selector, and transverse
Gaussian-body thresholds as EX23: respectively $4\sqrt m$,
$4\sqrt m(C/2)$ after applying $Q_0^{1/2}$, and
$4\sqrt{mW_j}$ after applying $(A_j|_{u_j^\perp})^{1/2}$.
The number of blocks is at most $3m+1$, each contributes failure
probability at most $1/(16m)$, and the Gaussian measure of the
product body exceeds one half. Fix the resulting signs
$\sigma_u$ for $u\in U$.

For $v\in V$ define the now fixed symmetric matrices
\[
 H_{v,j}=\sum_{\substack{u\in U\\e=uv\in j}}\sigma_uM_e.
 \tag{EX42.8}
\]
They need not be positive semidefinite. Nevertheless, weighted
Cauchy--Schwarz, applied to $H_{v,j}y$, gives
\[
 H_{v,j}^2\preceq
 B\sum_{\substack{u\in U\\e=uv\in j}}M_e,
 \qquad
 \sum_{v\in V}H_{v,j}^2\preceq BQ_{UV,j}\preceq BQ_j.
 \tag{EX42.9}
\]
For the full whitening, put
$H'_{v,j}=Q_j^{-1/2}H_{v,j}Q_j^{-1/2}$ when $q_j>0$.
Writing $M'_e=b_eN_e$ gives $N_e^2\preceq N_e/q_j$ and hence
\[
             \sum_{v\in V}(H'_{v,j})^2\preceq (B/q_j)I.
 \tag{EX42.10}
\]
These bounds hold uniformly in the already chosen signs on $U$.

Apply Banaszczyk's theorem a second time, now to the vertices in $V$.
Use the same vector/count blocks, together with the two Hilbert--Schmidt
covariance-whitening blocks from EX23 for
$P_jH_{v,j}P_j$ on $u_j^\perp$ and $H'_{v,j}$, respectively.
Here $P_j$ is orthogonal projection onto $u_j^\perp$. The covariance
device is valid for arbitrary symmetric matrices, so the possible
indefiniteness in (EX42.8) is harmless. The two Gaussian matrix-series
thresholds may be taken as
$2\sqrt{BW_j\ell}$ and $2\sqrt{(B/q_j)\ell}$, by
(EX42.9)--(EX42.10) and Tropp's Gaussian matrix-series bound cited in
EX23. There are at most $5m+1$ blocks. Their product body has Gaussian
measure at least $1-(5m+1)/(16m)\ge5/8$, so the same scaling by
$\Lambda$ is valid.

Combine the two signings into $\sigma_a$ on the entire vertex set.
Each signing has count discrepancy less than $50m$, so
$|\sum_a\sigma_a|<100m$ and $k\ge300m$ ensures that the positive
set $S$ is balanced. Since the global sum of each $v_{aj}$ is zero,
$R_j=\tfrac12\sum_a\sigma_av_{aj}$. The vector blocks and
$\Lambda^2\le150m$ give
\[
 \|R_j\|\le100mBC/r_j,\quad
 |(R_j)_{u_j}|\le50mB,\quad
 \|(R_j)_\perp\|^2\le2500m^2BW_j.
 \tag{EX42.11}
\]
This proves (EX42.2).

\emph{Why the second signing controls the final curvature.}
Let $Q_{{\rm cut},j}$ denote the sum of $M_e$ over edges crossing
the final cut. The edges between $U$ and $V$ contribute exactly
\[
 \frac12\left(Q_{UV,j}-\sum_{v\in V}\sigma_vH_{v,j}\right).
 \tag{EX42.12}
\]
The crossed edges lying within either auxiliary part add positive
semidefinite matrices, so (EX42.12) is a lower bound for
$Q_{{\rm cut},j}$. The second signing makes the perpendicular
matrix error at most $2\Lambda\sqrt{BW_j\ell}$, and the full
whitened error at most $2\Lambda\sqrt{(B/q_j)\ell}$. For
$W_j\ge TmB\ell$, the former is at most $W_j/16$, since
$4\Lambda^2\le600m$ and $T>256\cdot600$. For
$q_j\ge TmB\ell$, the latter is at most $1/8$.
Together with (EX42.5)--(EX42.6), this yields
\begin{align}
 Q_{{\rm cut},j}|_{u_j^\perp}&\succeq W_jI/32
       &&\text{if }W_j\ge TmB\ell,\tag{EX42.13}\\
 Q_{{\rm cut},j}&\succeq Q_j/16
       &&\text{if }q_j\ge TmB\ell.\tag{EX42.14}
\end{align}
The auxiliary random signs and the final Banaszczyk signs play
different roles. No martingale concentration statement is asserted
for the final signing, and no change of probability measure is used.

\emph{Saving in each band.}
The spring reduction in EX41 is valid for arbitrary segment centers.
For this band alone, project the difference of the hubs to the ball
of radius $r_j$, which can only increase its relaxed saving. Write
$d=r_jz$ with $\|z\|\le1$. The denominator bound from EX41 gives
\[
 S_j(r_jz)\le r_j\left(
 R_j\cdot z-\tfrac14z^{\mathsf T}Q_{{\rm cut},j}z\right).
 \tag{EX42.15}
\]
If $W_j<TmB\ell$, the entire band cost is at most $W_jr_j$,
which proves (EX42.3). Otherwise use (EX42.13).
When $q_j\ge TmB\ell$, (EX42.14) and unrestricted quadratic
maximization give
\[
 \sup_z\left(R_j\cdot z-
       \tfrac14z^{\mathsf T}Q_{{\rm cut},j}z\right)
 \le16R_j^{\mathsf T}Q_j^{-1}R_j
 \le120000m^2B.
\]
Here $Q_j|_{u_j^\perp}\succeq W_jI/2$, $q_j\ge B$, and
(EX42.11) were used.
In the remaining case $q_j<TmB\ell$, the top entry of
$Q_{{\rm cut},j}$ in the decomposition
$\mathbb Ru_j\oplus u_j^\perp$ is at most $q_j$; positivity
and $\|Q_{{\rm cut},j}\|\le W_j$ bound the squared norm of its
off-diagonal block by $q_jW_j$. Young's inequality and
(EX42.13) therefore give
\[
 z^{\mathsf T}Q_{{\rm cut},j}z
 \ge\frac{W_j}{64}\|z_\perp\|^2-63q_jz_u^2.
\]
Consequently the expression in parentheses in (EX42.15) is at most
\[
 50mB+64\frac{\|(R_j)_\perp\|^2}{W_j}+16q_j
 \le50mB+160000m^2B+16TmB\ell
 \le K_0B.
\]
For $D=1$, all transverse blocks vanish and the principal bound
alone gives the same conclusion. This proves (EX42.3).

\emph{Truncation, dyadic summation, and the final exponent.}
Let $l_{\max}=\max_e l_e\le C$, and discard edges of length below
$\delta=aC/(4B)$. Their entire target cost is at most
$\delta\sum_e b_e\le akC/8$, so their possible saving is at most
$akC/4$. If no edge remains, strict improvement is impossible.
Otherwise the retained dyadic bands have
$r_j=2^{-j}l_{\max}$ and
$m\le1+\log_2(4B/a)=O_a(\log(2+B))$.
Project to the span of retained directions when applying the cut
lemma; this does not increase the certificate ellipsoid trace.

For the actual hub difference $d$, the supporting norm inequality
also gives $S_j(d)\le R_j\cdot d\le\|R_j\|\|d\|$.
Thus the one cut in the lemma has retained saving at most
\[
 \sum_j\min\{K_0Br_j,L_0BC\|d\|/r_j\}
 \le4B\sqrt{K_0L_0C\|d\|},
\]
by the same dyadic split as EX23. Paying the bridge and selector
requires
\[
 \frac34akC
 <4B\sqrt{K_0L_0C\|d\|}-h\|d\|
 \le\frac{4B^2K_0L_0C}{h},
 \qquad
 h<\frac{16B^2K_0L_0}{3ak}.
 \tag{EX42.16}
\]
Independently, apply EX16's uniformly balanced random-cut calculation
to the unit direction resultants of all original edges. For any cut,
the supporting spring inequality permits improvement only if
$h<\|R_S\|$. Since every balanced cut must improve, a possibly
different cut gives $h\le B\sqrt{k/2}$. This handles $k<300m$;
otherwise its combination with (EX42.16) gives
$h=O_a(B^{4/3}(K_0L_0)^{1/3})$.

Finally, omit the bridge charge and use the same simultaneous free
saving bound to obtain
$\tfrac34akC<\sum_jK_0Br_j\le2K_0BC$.
Positive integer widths imply
$D\le\#\{\text{retained edges}\}\le kB/2$. Exactly the logarithmic
bootstrap in EX23 now yields
$k=O_a(B(\log(2+B))^2)$ and $\ell=O_a(\log(2+B))$.
Hence $K_0=O_a((\log(2+B))^2)$ and
$L_0=O_a(\log(2+B))$, proving (EX42.1).
The use of a different auxiliary or final cut for a separate upper
bound is legitimate because the hypothesis covers every balanced
cut with the same selector charge. \hfill$\square$



\subsection{Reducing the actual selector question to endpoint normalization}
\label{sec:ex-endpoint-normalization}

\paragraph{EX43: a fixed-scale endpoint-to-hub extension.}
Fix a field and an integer $k\ge2$, and put
$P_k=\{\pm e_1/2,\ldots,\pm e_k/2\}$.
Let $\Delta(P_k)$ denote the ordinary normalized tensor spectrum on these
parties. Define
\[
 \Gamma_k=\inf\{c>0:\ \exists\psi\in\Delta(P_k),\quad
           c\,d_\psi(p,q)=\|p-q\|\text{ for all }p,q\in P_k\},
 \tag{EX43.1}
\]
with value $+\infty$ when the set is empty. When finite, this infimum is
attained. The substantive reduction is the following: for every feasible
$c,\psi$ in (EX43.1), and every finite set
\[
 H\subseteq \{x\in\mathbb R^k:\|x\|\le\rho_k\},\qquad
 \rho_k=\tfrac12\,2^{-\lceil\log_2k\rceil/2}
             \ge\frac1{2\sqrt{2k}},
 \tag{EX43.2}
\]
there is a spectral point $\varphi$ on $P_k\cup H$ satisfying
\[
 c\,d_\varphi(x,y)=\|x-y\|\quad(x,y\in P_k\cup H),
 \qquad g_\varphi(P_k)=1.
 \tag{EX43.3}
\]
The enlarged point can itself be taken ordinary normalized. In particular,
the scale $c$ and the partial GHZ exponent are preserved simultaneously.
Preservation of every old tensor value is not needed or asserted here.

\emph{Compactness and the meaning of $\Gamma_k$.}
The ordinary spectrum on a fixed finite party set is compact in the
pointwise topology: for each tensor $T$, its value lies between zero and
its finite tensor rank, and normalization, multiplicativity, and the
cover inequalities are closed conditions. The normalization lemma used
in EX13--15 identifies every nontrivial enlarged spectral point with a
positive power of an ordinary one. If scaled wire metrics on $P_k$
converge to its Euclidean metric, normalize their full unit exponents
to one, absorbing the exponent into the scale. Call the resulting
scales $c_n$. A wire is a restriction of the full unit tensor, whereas
the unit tensor is a restriction of a product of wires along any
spanning tree. Thus
\[
 \max_{p,q}c_n d_{\psi_n}(p,q)\le c_n
 \le\sum_{e\in T}c_n d_{\psi_n}(e).
\]
The right side stays bounded and the left side tends to one. A compact
subnet therefore gives an exact feasible pair in (EX43.1). Conversely,
an exact pair is a price by definition. The same compactness argument
applied to a minimizing sequence proves attainment. In particular,
$\Gamma_k<\infty$ if and only if the endpoint metric is a price, and
$\Gamma_k\ge1$ whenever it is finite. This argument does not deduce a
dimension-independent upper bound on $\Gamma_k$.

\emph{Midpoint extension preserving the original GHZ exponent.}
A tensor-preserving completion is also stated in Liu's unrefereed manuscript (2026a),
\emph{A uniform exponent bound for asymptotic tensor rank},
Lemma~2.1, printed p.~2. Here the finite-label spectral compactness
step additionally records the exact partial GHZ value needed in
(EX43.3). We use the exact form of Anthropic Proposition~A.7, printed pp.~21--22,
not just its metric consequence A.8. If $I$ is the current finite party
set containing the original $P_k$, choose its anchor tensor to be
\[
 C_I=\langle2\rangle_{P_k}
           \prod_{\{x,y\}\subseteq I}\langle2\rangle_{xy}.
 \tag{EX43.4}
\]
Apply A.7 to a finite random retraction that fixes $I$ and sends each
new midpoint label to either endpoint of its pair with probability one
half. For every tensor embedded from $I$, the extension upper bound in
A.7 is its old value, since every retraction is the identity on $I$.
The anchor product has exactly its old value. Each of its nonzero
factors can only decrease; multiplicativity therefore makes every
factor retain its old value. Consequently all old wire distances and
the value of $\langle2\rangle_{P_k}$ are unchanged.
This factor-preservation mechanism is also used explicitly in the
proof of Anthropic Lemma~E.8, printed pp.~35--36.

Initially $g_\psi(P_k)=1$. It remains one at every extension step.
Monotonicity under deleting parties from a unit tensor gives
$g(P_k)\le g(I)$, while the cover inequality gives $g(I)\le1$.
Thus the enlarged full unit exponent is also one, and the earlier
normalization lemma makes the enlarged point ordinary normalized.
The new midpoint distances are at most half the old endpoint distance;
the triangle inequality forces equality. Their CN inequality follows
from the checked price-curvature theorem, Anthropic Theorem~C.1,
printed p.~25.

Starting from $P_k$, repeatedly add midpoints for every pair in the
current finite set. Identify points at distance zero and complete the
resulting metric space. The proof of Anthropic Proposition~C.3,
printed p.~28, then gives a complete CAT(0) space $X$. Work with its
metric multiplied by the original scale $c$, so $P_k$ embeds with its
Euclidean distances. This use of the price-space construction has the
following additional normalization record. For any finite set
$F\subseteq X$ containing $P_k$, approximate its new points by labels
from finite midpoint stages, while keeping all original labels fixed.
Pull back those stages' ordinary spectral points to the fixed label
set $F$. They still have $g(P_k)=1$, and their scaled wire metrics
converge to the metric of $F$. Compactness of the spectrum on $F$
therefore supplies a point with the exact metric and $g(P_k)=1$.
There is no rescaling at this limiting step. The same argument covers
coincident labels by pullback. Thus every finite part of this completion
has a realizing spectral point with the original scale and partial GHZ
normalization intact.

\emph{A balanced-block rigidity lemma.}
Let $f:\mathbb R^d\to X$ be a $1$-Lipschitz map into a complete CAT(0)
space that preserves the distances of all vertices $\{\pm e_i\}$.
Then $f$ is isometric on the ball of radius
\[
                     R_d=2^{-\lceil\log_2d\rceil/2}.
 \tag{EX43.5}
\]
It is also isometric on the union of that ball with the original
vertices. This improves the sequential-coordinate radius in the
earlier rigid-ball lemma; only flat squares are required.

Here is a proof with the dimension dependence explicit. For $d=1$,
distance preservation at the two ends and the triangle inequality
force the whole intervening segment to be mapped isometrically.
For $d\ge2$, put $s=\lceil\log_2d\rceil$ and split the coordinates
into two nonempty orthogonal blocks, each with at most $2^{s-1}$
coordinates. Inductively the restriction of $f$ to each block is
isometric on its centered ball of radius
$R=2^{-(s-1)/2}$. Take any $z=x+y$ in the full ball of radius
$R/\sqrt2$, with $x,y$ in the two blocks. Choose vectors $u,v$ of
norm $R$ in those respective blocks, in the directions of $x,y$
when these are nonzero. Then
\[
 \|x\|+\|y\|\le\sqrt2\|z\|\le R,
 \qquad z\in\operatorname{conv}\{u,v,-u,-v\}.
\]
The two diameters $[u,-u]$ and $[v,-v]$ are preserved. CAT(0)
roundness, combined with the four nonexpansive side bounds, gives
\[
 8R^2\le d(fu,fv)^2+d(fv,f(-u))^2
            +d(f(-u),f(-v))^2+d(f(-v),fu)^2\le8R^2.
\]
Every side is therefore exact. The square equality case of the flat
rectangle lemma (Anthropic Lemma~C.4(f), printed p.~29, proved on
p.~65) realizes this square as a Euclidean square in $X$.
Moreover $f$ agrees with that isometry on the entire Euclidean
diamond: it agrees on each boundary segment by equality of its end
distance, and then on each chord joining opposite boundary segments
by the same equality and uniqueness of geodesics. Consequently
$d(fz,f(-z))=2\|z\|$.

Now take any two $z,w$ in the ball. Their two antipodal diameters
are preserved, and the Euclidean parallelogram identity gives
\[
 4(\|z\|^2+\|w\|^2)
 \le\sum_{\text{four cyclic sides}}d(fz,fw)^2
 \le2\|z-w\|^2+2\|z+w\|^2
 =4(\|z\|^2+\|w\|^2),
\]
where the cyclic vertices are $z,w,-z,-w$. All four side bounds
are equalities, proving isometry of the ball. The same argument
with $z$ an original vertex and $w$ in the ball proves the extra
assertion about their union. This completes the induction. It also
shows explicitly why preserving only the ball's internal distances
would have been insufficient for the endpoint-to-hub application.

\emph{Application to the fixed-scale completion.}
Use the Kirszbraun theorem for a complete CBB(0) domain and complete
CAT(0) target: Alexander, Kapovitch, and Petrunin (2011),
\emph{Alexandrov meets Kirszbraun}, \S5, Theorem~5.4, author-PDF p.~11.
It extends the isometric vertex map $P_k\to X$ to a $1$-Lipschitz
map $f:\mathbb R^k\to X$. Its hypotheses hold for Euclidean space
and the complete CAT(0) space just constructed; no properness of
$X$ is assumed. Scale the balanced-block lemma by one half. Then
$f$ is isometric on $P_k$ together with the ball in (EX43.2).
Apply the preserved-normalization finite-subset property of $X$
to $f(P_k\cup H)$. This proves (EX43.3).

\emph{The strictly smaller gate for the complete-network argument.}
In Anthropic Proposition~E.10, printed pp.~36--37, every hub has
norm $\eta=2/(9\sqrt k)$. Since
$2/(9\sqrt k)<1/(2\sqrt{2k})\le\rho_k$, all hubs for all balanced
cuts belong to the ball in (EX43.2). There are finitely many of them.
Thus an endpoint minimizer for $\Gamma_k$ extends to one spectral
point realizing all required endpoint and hub distances exactly,
with scaled endpoint GHZ exponent equal to $\Gamma_k$.
Pullback to the physical end labels preserves the distances; if
some of the $2k$ endpoint positions are unused, their actual partial
unit exponent is at most one, which only improves this selector
upper bound.

Accordingly EX10 no longer needs to assume simultaneous control on
the enlarged endpoint-and-hub configuration. Under the flat hypothesis,
it is enough to prove, for all sufficiently large $k$,
\[
                         \Gamma_k=O(k^\beta),\qquad\beta\ge0.
 \tag{EX43.6}
\]
If this holds, the complete-network comparison, conditional on the
claimed circuit reduction being correct, yields
\[
            2^{\Omega((\log m)^\alpha)},\qquad
            \alpha=1+\frac1{2(1+\beta)}.
 \tag{EX43.7}
\]
The parameter calculation is EX10: choose cable width
$\ell$ a sufficiently large multiple of $k^\beta$, put $B=k\ell$,
and obtain bridge scale $B\sqrt k$.
The established binary certificate gives $\Gamma_k\le\sqrt k$
under the flat hypothesis, by Anthropic Proposition~E.9, printed
pp.~35--36, and compactness. Explicitly, restrict each approximate
realization to $P_k$, write $\alpha=g(P_k)>0$, power-normalize that
restricted point by $1/\alpha$, and absorb $\alpha$ into the old
metric scale. The new scale is the bounded quantity $c\alpha$
from E.9; its compact limit is at most $\sqrt k$. This recovers $4/3$. A bound
$\Gamma_k=O(1)$ would give $3/2$ by this same direct argument.
No such improved bound is proved here. The unresolved object is now
the ordinary spectrum on just the $2k$ endpoints, rather than a
simultaneous spectral realization on all exponentially many hubs.



\paragraph{EX44: a finite-cover test for the endpoint normalization.}
The endpoint quantity in EX43 has the following exact operational
formulation. This is a dual reformulation, not an improved estimate of
$\Gamma_k$. Write $d_{pq}=\|p-q\|$ on $P_k$, and set
\[
 W_k=\prod_{\{p,q\}\subseteq P_k}\langle2\rangle_{pq},
 \qquad L_k=\sum_{\{p,q\}\subseteq P_k}d_{pq}
          =k+\sqrt2\,k(k-1).
 \tag{EX44.1}
\]
Fix $t>0$. A resource is any tensor
\[
 S(m,\mathbf r)=\langle m\rangle_{P_k}
                  \prod_{\{p,q\}\subseteq P_k}\langle r_{pq}\rangle_{pq},
 \qquad m,r_{pq}\in\mathbb Z_{\ge1},
\]
and its assigned cost is
\[
                     w_t(m,\mathbf r)
                 =m\prod_{\{p,q\}}r_{pq}^{d_{pq}/t}.
 \tag{EX44.2}
\]
For an arbitrary tensor $T$ on these parties, let $p_t(T)$ be the
infimum of the sums of the assigned costs over all finite covers of
$T$ by such resources. Covers mean sums of party-local restrictions,
with the same convention as Anthropic Definition~A.1. Let
\[
                       R_t(T)=\inf_{N\ge1}p_t(T^N)^{1/N}.
 \tag{EX44.3}
\]
Then, over the chosen field,
\[
 \boxed{\quad \Gamma_k\le t
       \quad\Longleftrightarrow\quad
       R_t(W_k)=2^{L_k/t}.\quad}
 \tag{EX44.4}
\]
Equivalently, $\Gamma_k>t$, including $\Gamma_k=+\infty$, if and
only if for some integer $N\ge1$ there is a finite cover
\[
 W_k^N\preceq\sum_{j=1}^A S(m_j,\mathbf r_j),
 \qquad
 \sum_{j=1}^A m_j\prod_{\{p,q\}}r_{pq,j}^{d_{pq}/t}
                           <2^{NL_k/t}.
 \tag{EX44.5}
\]
No border closure is part of this finite certificate.

To prove this, first observe that $p_t$ is finite: a tensor of finite
rank is a restriction of a full unit tensor, which is an allowed
resource. It is an upper functional. Indeed, covers compose and
concatenate, and resource costs multiply under tensor products.
Also $p_t(1)=1$: every nonzero resource has assigned cost at least one,
every cover of $1$ has a nonzero term, and the resource with all
dimensions one attains one. This proves normalization even when
$t<1$; no assumption that the assigned wire costs are below tensor
rank is made. The regularization argument in Anthropic
Proposition~A.5, printed pp.~20--21, shows that $R_t$ is again an
upper functional and satisfies $R_t(T^a)=R_t(T)^a$ for every
integer $a\ge1$. The direct resource for $W_k$ gives
$R_t(W_k)\le2^{L_k/t}$.

If $\Gamma_k\le t$, take the attained ordinary spectral point
$\psi$ with scale $c=\Gamma_k$. The enlarged spectral point
$\varphi=\psi^{c/t}$ is valid because $0<c/t\le1$ and
$(\sum_j a_j)^{c/t}\le\sum_j a_j^{c/t}$ for nonnegative $a_j$.
It has
\[
 \varphi(\langle r\rangle_{pq})=r^{d_{pq}/t},\qquad
 \varphi(\langle m\rangle_{P_k})=m^{c/t}\le m.
\]
Consequently $\varphi\le p_t$ and, by multiplicativity,
$\varphi\le R_t$. Its value on $W_k$ is the stated upper bound,
which proves equality.

Conversely, suppose equality holds. Apply the saturation lemma
PG-SAT to $R_t$ and $W_k$. It gives an enlarged spectral point
$\varphi\le R_t\le p_t$ with
$\varphi(W_k)=2^{L_k/t}$. Each binary wire has value at most
$2^{d_{pq}/t}$, and their product attains the product of these
positive upper bounds. Thus every wire attains its bound. Let
$\alpha=\log_2\varphi(\langle2\rangle_{P_k})$. The resource with
$m=2$ and all $r_{pq}=1$ gives $\alpha\le1$. Every wire is a
restriction of this full unit, so $\alpha>0$. By the normalization
lemma of EX13--15, $\psi=\varphi^{1/\alpha}$ is ordinary normalized.
It realizes the endpoint metric with scale $t\alpha\le t$.
This proves (EX44.4). Finally, failure of equality is strict because
of the direct upper bound; the infimum in (EX44.3), followed by the
definition of $p_t$, turns strict inequality into one finite cover.
The reverse implication is immediate. This proves (EX44.5).

The certificate keeps the GHZ dimension separate from the number of
summands. Both contribute to the cost, and treating a coherent
selector as free would change the problem. The result supplies a
finite search target on $2k$ parties but establishes no new growth
bound for $\Gamma_k$ and no new circuit exponent by itself.


\paragraph{EX45: a $3/2$ ceiling with an actual spectral selector.}
The second bound of EX41 has a wider selector scope than its statement
there. It does not require a binary pointwise-averaging certificate.
Keep the fixed loopless graph, positive integer edge widths, maximum
incident width $B$, arbitrary positive target lengths $l_e$, arbitrary
segment centers, exactly two hubs, and balanced-cut conventions of
EX40. Put $L=\max_e l_e>0$. Suppose that there is a number $G\ge L$,
fixed before the cut is selected, such that for every balanced cut
one can place the two hubs and a bridge of width at least $h$ with
\[
 \text{fan cost}+h\,\|\text{hub}_1-\text{hub}_2\|+akG
                         <\sum_e b_e l_e,
 \tag{EX45.1}
\]
where $a>0$ is fixed. Then
\[
                 h=O_a\bigl(B^{3/2}\sqrt{\log(2+B)}\bigr).
 \tag{EX45.2}
\]
The quantity $G$ may be the actual scaled partial-GHZ exponent of a
spectral point; it need not be the cost of an averaging construction.

Indeed, the free-saving argument of EX41 depends only on the unit
directions and widths. Its paired-sign matrix-martingale construction
gives a balanced cut for which every two-hub placement has free saving
at most
\[
                  20000BL\ell,\qquad
                  \ell=\ln(64(D+1)),
 \tag{EX45.3}
\]
where $D$ is the span dimension of the target directions. The arbitrary
segment centers were removed by the spring relaxation, and the
unequal lengths were handled by the common denominator $4L$; neither
step uses any property of a selector. Assumption (EX45.1) therefore
implies
\[
                      akG<20000BL\ell,
 \qquad k<20000B\ell/a.
 \tag{EX45.4}
\]
The second inequality uses only $G\ge L$.

For every cut, convexity of the norm in the same spring relaxation
gives free saving at most $R_S\cdot d$, where $d$ is the hub
difference and $R_S$ is the signed sum of the crossed unit directions
with their widths. Thus strict improvement requires
$h<\|R_S\|$. The independent balanced random-resultant argument in
EX40 gives a balanced cut with
$\|R_S\|\le B\sqrt{k/2}$, with the same harmless changes of constants
for the inherited odd-$k$ convention. Different cuts may establish
(EX45.4) and this resultant bound, because (EX45.1) is required for
every balanced cut with the same $G$.

Positive integer widths give
$D\le |E|\le\sum_e b_e\le kB/2$. Consequently (EX45.4) has the
form $k\le C_a B\log(2+kB)$. The elementary logarithmic bootstrap
used in EX19 gives $k=O_a(B\log(2+B))$. Substituting this in
$h\le B\sqrt{k/2}$ proves (EX45.2). Bounded $k$ is included in this
bound. In particular, no selector zonotope estimate and no comparison
$G\le C$ with a binary certificate have entered the proof.

Here is why an actual spectral selector satisfies the needed lower
bound. Let $P$ contain all physical endpoint parties, let $\varphi$
be one spectral point on the endpoints and all hub labels used, and
suppose $c\,d_\varphi(p,q)$ equals the Euclidean distance of those
labels. Put
\[
                         G=c\,g_\varphi(P).
\]
For each target edge, its binary wire is a restriction of
$\langle2\rangle_P$: apply the all-ones linear functional at the
other parties. Spectral monotonicity therefore gives
\[
                   G\ge c\,d_\varphi(p_e,q_e)=l_e.
\]
Taking the maximum proves $G\ge L$. The same lower bound holds in
the limit for exact price realizations obtained by compactness.

The debit $akG$ remains an explicit hypothesis. This corollary does
not prove that every possible circuit decomposition must incur that
many selector bits, nor does it allow additional hubs, a changing
graph, or a cut-dependent spectral point. Under the same size encoding as EX18, the exactly-two-hub direct
method with this debit has fixed-power exponent at most $3/2$.
For the circuit consequence, \emph{if the claimed theorem is correct},
the checked $4/3$ construction supplies the lower endpoint,
so this broader actual-selector model has the rigorous interval
$[4/3,3/2]$, rather than an asserted exact supremum. EX43 identifies
an endpoint-normalization bound that would attain the upper endpoint
within the complete-network construction.


\paragraph{EX46: arbitrary endpoint geometry and polynomially many fixed hubs.}
The actual-selector ceiling extends to fixed labelled hub graphs with
many hubs, without a common origin or equal target lengths. Let a
loopless graph on $k$ circuit vertices have positive integer widths
$b_e$, total width $W=\sum_e b_e>0$, and maximum incident width $B$.
Its two terminal positions $p_e,q_e$ are arbitrary Euclidean points.
Let $P$ contain every endpoint party, and suppose its diameter is
positive. Fix $r=2^t\mid k$, $t\ge1$, and a connected labelled graph
on the $r$ hubs. This graph and the assignment of group $j$ to hub
$j$ are fixed before choosing a labelled partition into $r$ groups
of equal size $s=k/r$.

Suppose $G\ge\operatorname{diam}P$ is fixed, and every such partition
admits positions for its hubs giving a strict improvement over the
target cost $\sum_e b_e\|p_e-q_e\|$, when the source pays the fan cost,
bridge width at least $h$ on every hub edge, and selector debit at
least $aktG$, with fixed $a>0$. Then
\[
 h=O_a\!\left(
       \frac{B^{3/2}\log(2+rB)}{\sqrt t}
                +Bt\log(2+rB)\right).
 \tag{EX46.1}
\]
In particular, for every fixed $M>0$, if $r\le B^M$, then
\[
                     h=O_{a,M}\bigl(B^{3/2}\log(2+B)\bigr).
 \tag{EX46.2}
\]
The bound concerns the minimum bridge width. The actual scaled
partial-GHZ exponent is an admissible $G$, as verified below.

\emph{A median supplies the equilibrium needed for the partition proof.}
Count endpoint incidences separately, with the width of their edge
as their weight. Choose a weighted geometric median $o$ of these
incidences, minimizing
\[
                         F_0=\sum_{\text{ends }x}b_x\|x-o\|.
\]
A minimizer exists by coercivity and can be chosen in the convex hull
of the endpoints by Euclidean projection. Translate $o$ to zero and
put $R=\max_x\|x\|$. Then
$0<R\le\operatorname{diam}P\le G$. The first-order subgradient
condition gives vectors $u_x$ with
\[
 u_x=x/\|x\|\quad(x\ne0),\qquad
 \|u_x\|\le1\quad(x=0),\qquad
                         \sum_xb_xu_x=0.
 \tag{EX46.3}
\]
At zero endpoints choose the subgradients in the linear span of the
translated endpoints: orthogonal projection into that span preserves
the equilibrium and cannot increase their norms. Call its dimension
$D$. Orthogonally projecting all hubs into this same span can only
decrease fan and bridge lengths, so the ambient space may be taken
to have dimension $D$. Positive integer widths give
$D\le2|E|\le2W\le kB$.

For each endpoint and every $\|z\|\le R$,
\[
 \|x-z\|\ge\|x\|-u_x\cdot z
       +\frac{\|z\|^2-(u_x\cdot z)^2}{4R}.
 \tag{EX46.4}
\]
For $x\ne0$, rationalize the difference between the norm and
$\|x\|-u_x\cdot z$. Its denominator is at most $4R$; the
zero-denominator case follows directly. For $x=0$, write
$\rho=\|z\|$ and $v=u_x\cdot z$, so $|v|\le\rho\le R$.
The assertion is
$\rho+v\ge(\rho^2-v^2)/(4R)$, which follows from
$(\rho+v)(\rho-v)\le4R(\rho+v)$.

Define, for each circuit vertex,
\[
 d_a=\sum_{x\text{ at }a}b_x,\qquad
 v_a=\sum_{x\text{ at }a}b_xu_x,\qquad
 H_a=\sum_{x\text{ at }a}b_x(I-u_xu_x^{\mathsf T}),
\]
and put
\[
 A=\tfrac12\sum_xb_xu_xu_x^{\mathsf T},\qquad Q=WI-A.
\]
These data have exactly the properties used by the non-selector
part of EX28:
\[
 \|v_a\|\le d_a\le B,\quad \sum_av_a=0,\quad
 0\preceq H_a\preceq BI,\quad \sum_aH_a=2Q,
 \quad A\succeq0,\quad\operatorname{tr}A\le W.
 \tag{EX46.5}
\]
In particular, $Q\succeq0$, and on the orthogonal complement of
a top eigenvector of $A$ one has $Q\succeq(W/2)I$.
The positive lifts
$V_a=\left(\begin{smallmatrix}d_a&v_a^{\mathsf T}\\v_a&d_aI\end{smallmatrix}\right)$
satisfy
\[
 \|V_a\|\le2B,\quad \sum_aV_a=2WI,
 \quad \sum_aV_a^2\preceq4BW I,
 \quad \sum_aH_a^2\preceq2BQ.
 \tag{EX46.6}
\]
The whitened bound used there also follows from
$H_a\preceq BI$ and $\sum_aH_a=2Q$. Thus replacing the earlier
$Q_a$ by $H_a$ changes no concentration hypothesis. The weaker
inequality $\operatorname{tr}A\le W$ suffices everywhere that the
earlier unit directions gave equality. No pairwise antipodality
is used after these properties and the fan inequality are available.

\emph{The bridge-dependent saving estimate.}
Delete extra hub edges to obtain a fixed spanning tree, and lower
its retained bridge widths to $h$; this only enlarges the permitted
saving. Project every hub into the ball of radius $R$ about the
median. Projection fixes the endpoints and is nonexpansive, so
neither fan lengths nor bridge lengths increase. Write $z_j=Ry_j$.
For the group assigned to hub $j$, let
$R_j=\sum_{a\text{ in group }j}v_a$ and
$H_j=\sum_{a\text{ in group }j}H_a$.
By the triangle inequality, the target cost is at most $F_0$.
Using (EX46.4), its saving before the selector is therefore at most
\[
 R\left\{\sum_jR_j\cdot y_j
       -\tfrac14\sum_jy_j^{\mathsf T}H_jy_j
       -h\sum_{ij\text{ in tree}}\|y_i-y_j\|\right\},
 \qquad\|y_j\|\le1.
 \tag{EX46.7}
\]
The nonnegative difference $F_0-\text{target cost}$ was discarded
only to obtain an upper bound.

Apply the sorted hierarchical pipage construction of EX28 to
(EX46.5)--(EX46.6), with
$\ell=\ln(64r^2(D+1))$. Its principal-coordinate discrepancy,
dyadic concentration, and preorder interval estimates are unchanged.
The coarse-block quadratic maxima satisfy
$\sum_I f_I\le7000(r/L)B\ell$ as in (EX28.6): the full-curvature
case uses $Q$ on the perpendicular space, and the near-line case
uses $\sum_Iu^{\mathsf T}H_Iu=2u^{\mathsf T}Qu$.
The vector and matrix residuals are absorbed along the fixed tree
exactly as in (EX28.7), since
$\|yy^{\mathsf T}-ww^{\mathsf T}\|_*\le2\|y-w\|$
on the unit ball. Consequently, when $h\ge256B\ell t$, there is
one labelled equal partition whose maximum saving before the selector
is at most
\[
                        R\,2^{38}\frac{B^2W\ell^2}{h^2}.
 \tag{EX46.8}
\]
The $L=1$ coarse estimate also gives a possibly different partition
with free saving at most $7000RrB\ell$, irrespective of bridges.
These are uses of the already proved quantitative partition argument;
its certificate-based tree-flow inequality is not used.

Feasibility for the partition in (EX46.8), and $2W\le kB$, give
\[
                h^2(G/R)t<(2^{37}/a)B^3\ell^2.
 \tag{EX46.9}
\]
The coarse free-saving partition gives
$st(G/R)<7000B\ell/a$. Since $G/R\ge1$ and $D\le rsB$,
the latter inequality implies
$s=O_a(B\log(2+B))$ by the same logarithmic bootstrap as EX28.
Hence $\ell=O_a(\log(2+rB))$. In the large-$h$ case (EX46.9)
proves the first term of (EX46.1); otherwise the threshold
$h<256B\ell t$ proves its second term. This establishes
(EX46.1)--(EX46.2).

For an actual spectral selector $G=c\,g_\varphi(P)$ in one
Euclidean realization of all involved labels, the restriction
$\langle2\rangle_P\succeq\langle2\rangle_{xy}$ is valid for
\emph{every} pair of endpoint parties, not just the target edges.
Thus $G\ge c\,d_\varphi(x,y)=\|x-y\|$, and so
$G\ge\operatorname{diam}P$. This is the stronger diameter bound
needed for the median argument. If all endpoints coincide, the
target cost is zero and strict improvement by nonnegative costs
is impossible, so the excluded case is vacuous.

The common selector debit $aktG$, the fixed labelled hub graph,
and the requirement for every equal labelled partition remain
essential model assumptions. A graph or label assignment chosen
after seeing the partition is not covered. With the same size
encoding as EX18, this entire model has fixed-power ceiling $3/2$.
For the lower endpoint, \emph{if the claimed theorem is correct},
the two-hub complete-network construction supplies $4/3$.
No exact supremum is asserted for arbitrary actual selectors;
in particular, this argument supplies no improved upper bound
for $\Gamma_k$ and no obstruction to its being bounded.


\paragraph{EX47: subpolynomially many adaptive hubs with actual selectors.}
Keep the endpoint geometry, widths, equal group sizes, and actual
selector debit $aktG$ of EX46, with $G\ge\operatorname{diam}P$.
Now allow the connected hub graph and the assignment of the $r=2^t$
groups to its labels to be chosen after each partition is revealed.
Suppose every equal partition is still strictly feasible and every
bridge has width at least $h$. Then
\[
 h=O_a\!\left(rB^{3/2}
          \sqrt{\frac{\log(2+rB)}{t}}\right).
 \tag{EX47.1}
\]
In particular, if $r=B^{o(1)}$, no fixed power strictly larger than
$3/2$ can arise in this model. The conclusion permits arbitrary
segment centers and lengths and arbitrary topology choices; the
subpolynomial bound on the number of hubs is substantive.

\emph{Proof.}
Use the weighted median, radius $R$, and vertex resultants $v_a$ of
EX46. The free-saving partition there, obtained from its coarse
$L=1$ estimate, ignores all bridges. It therefore applies even if
the graph and label assignment are selected later. Feasibility gives
\[
                       st(G/R)=O_a(B\ell),\qquad
               \ell=\ln(64r^2(D+1)),\quad s=k/r.
\]
As before, $G/R\ge1$ and $D\le rsB$ imply
$\ell=O_a(\log(2+rB))$ and
$s=O_a(B\log(2+rB)/t)$.

For a separate partition, minimize $\sum_{j=1}^r\|R_j\|^2$
over all equal-size partitions, where $R_j=\sum_{a\in S_j}v_a$.
Such a minimum exists because there are finitely many partitions.
For $a\in S_i,b\in S_j$, optimality under swapping $a,b$ gives
\[
 (R_i-R_j)\cdot(v_a-v_b)\le\|v_a-v_b\|^2.
\]
Sum over the $s^2$ choices of $a,b$. Since $\|v_a\|\le B$,
\[
             s\|R_i-R_j\|^2\le4s^2B^2.
\]
The identity $\sum_jR_j=0$ then implies
$\|R_j\|\le2B\sqrt s$ for every group. For every subset
$J$ of the groups, use this bound on $J$ or on its complement to get
\[
                         \left\|\sum_{j\in J}R_j\right\|
                              \le rB\sqrt s.
 \tag{EX47.2}
\]
This controls every subset, so it is unaffected by subsequent
relabelling or graph selection.

For any subsequently selected connected graph, choose a spanning
tree and root it. The convex supporting inequality for each endpoint
at the median gives
\[
 \text{target cost}-\text{fan cost}
       \le F_0-\text{fan cost}\le\sum_jR_j\cdot z_j.
\]
Tree summation by parts expresses the last sum using the resultant
of the group subset below each edge. By (EX47.2), it is at most
$rB\sqrt s$ times the total length of this tree. The source bridge
cost is at least $h$ times that length, including when other graph
edges are present. Thus strict improvement is impossible if
$h\ge rB\sqrt s$. Combining $h<rB\sqrt s$ with the earlier bound
on $s$ proves (EX47.1). The two partitions may differ because
feasibility is assumed for every equal partition.

For $r=B^{o(1)}$, the additional factor in (EX47.1) is
$B^{o(1)}$, giving the stated fixed-power ceiling. For polynomially
growing adaptive $r$, the factor $r$ cannot simply be dropped: the
actual-selector countermodel in EX39 remains feasible as an abstract
layout model. Its cheap algebraic cover is separately ruled out by
the cut-rank calculation there. As in EX46, the circuit lower
endpoint $4/3$ is available only \emph{if the claimed theorem is
correct}; neither this corollary nor the countermodel determines the
unrestricted circuit-method supremum.


\paragraph{EX48: arbitrary graph-certified seeds reduce to binary seeds.}
\label{ex:graph-seed-reduction}
Let $P$ be a finite labelled Euclidean configuration.  Consider a finite
pointwise convex representation
\[
                    \mathbb E_\omega a_\omega(p)=p
                       \qquad(p\in P).
\]
For each seed map $a=a_\omega$, permit an arbitrary finite graph whose
terminal vertices carry the labels $p\in P$ at positions $a(p)$, with
arbitrarily many auxiliary vertices at arbitrary Euclidean positions.
Positions may coincide.  Let its nonnegative edge loads be $b_e$ and
let $\rho>0$ satisfy
\[
       \sum_{e\in\delta(S)}b_e\ge\rho
       \quad\hbox{whenever $S$ separates the terminal labels}.
       \tag{EX48.1}
\]
Assign this seed the cost
\[
                  C(a)=\rho^{-1}\sum_e b_e\ell_e,
       \qquad \ell_e=\|z_{e,+}-z_{e,-}\|.
\]
The representation has a binary-seed refinement with exactly the same
pointwise mean and expected image diameter at most
$2\mathbb E C(a_\omega)$.  In particular, for
$P_k=\{\pm e_i/2:1\le i\le k\}$,
\[
                         \mathbb E C(a_\omega)\ge\frac{\sqrt k}{2}.
                         \tag{EX48.2}
\]
The number of seed values and auxiliary graph vertices is unrestricted.

\emph{Proof.}
The vector $x_e=b_e/\rho$ is feasible for the undirected Steiner cut
relaxation on this fixed graph.  Goemans and Williamson (1995),
Theorem~2.4, printed pp.~303--304, bounds an integral feasible forest
by at most twice the fractional optimum; their terminal-separating
function is listed in Table~1, p.~297, and discussed in Section~3,
p.~305.  The component containing all terminals is a tree after
unnecessary edges are removed.  Thus there is an embedded tree $T$
containing all terminal labels whose total length $L$ satisfies
\[
                             L\le2C(a).
                             \tag{EX48.3}
\]
The theorem for rational edge costs extends to the present nonnegative
real lengths by approximation on the fixed finite graph.  In particular,
zero-length edges cause no difficulty.

Root $T$ at a vertex of position $z_0$.  For each oriented tree edge
$e$, let $u_e$ be its displacement toward the child and let $S_e$
be the set of terminal labels below it.  Tree telescoping gives
\[
                         a(p)=z_0+\sum_e\mathbf1_{S_e}(p)u_e.
\]
If $L>0$, choose an edge of positive length with probability
$\|u_e\|/L$ and use the binary-image map
\[
               B_e(p)=z_0+\frac{L}{\|u_e\|}\mathbf1_{S_e}(p)u_e.
\]
Then $\mathbb E_e B_e(p)=a(p)$ and
$\operatorname{diam} B_e(P)\le L$.  If $L=0$, the seed is constant
and already has zero binary cost.  Composing this refinement with
$\omega$ proves the first assertion.  The binary images are allowed
at arbitrary Euclidean positions, as in EX-05; they need not stay
inside the original seed's convex hull.  Applying EX-05 proves
(EX48.2).

An equivalent direct lower-bound calculation is useful for checking
the normalization.  Define
\[
             F(a)=\sum_{i=1}^k
                  \langle e_i,a(e_i/2)-a(-e_i/2)\rangle.
\]
The coefficient of each tree displacement in this functional is
$\sum_i\epsilon_i e_i$, with $\epsilon_i\in\{-1,0,1\}$, hence has
norm at most $\sqrt k$.  Therefore
$F(a)\le\sqrt kL\le2\sqrt kC(a)$.  The pointwise mean identity
implies $\mathbb EF(a)=k$, again giving (EX48.2), including coincident
labels.

\emph{Tensor-certificate scope.}
Condition (EX48.1) is necessary whenever the edge loads of the seed
graph certify a GHZ rate $\rho$ on all terminal labels.  More
explicitly, suppose graph-wire tensors with logarithmic edge
dimensions $Nb_e+o(N)$ cover a terminal GHZ tensor of logarithmic
dimension $N\rho+o(N)$ using $2^{o(N)}$ restrictions.  Across any
terminal-separating cut, flattening ranks give
\[
                   N\rho+o(N)\le N\sum_{e\in\delta(S)}b_e+o(N).
\]
Divide by $N$ and take the limit.  This argument also covers exact
restrictions and ordinary asymptotic subrank certificates, over any
field; no converse to the cut condition is assumed.

Consequently, an argument that charges a selector using these
pointwise-averaged graph certificates cannot improve the existing
binary-certificate power exponent merely by allowing the seed size
or graph size to grow.  Indeed, if $C_{\rm bin}\le2C_{\rm graph}$
is the refined cost, a prescribed debit $akC_{\rm graph}$ dominates
$(a/2)kC_{\rm bin}$, and likewise with the additional factor
$\log_2r$ in the many-hub models.  The proved binary-model barriers
therefore apply with a changed constant.  This is a statement about
these certificates, \emph{not} a lower bound for the actual spectral
quantity $\Gamma_k$, and not a proof that an arbitrary actual selector
must pay the graph-certificate cost.

\subsection{Why exact graph restrictions do not yet settle the spectral gate}
\label{sec:ex-robustness}

\paragraph{EX-22: equal-size exact obstructions can disappear asymptotically.}
For every field and integer $q>1$, the matrix multiplication tensor
$\mathrm{MM}_q$ does not restrict exactly to the three-party unit tensor
$\langle q^2\rangle$. Both have all three flattening ranks equal to
$q^2$, so any such restriction would have invertible local maps.
However, a nonzero slice of $\mathrm{MM}_q$ has matrix rank
$q\,\operatorname{rank}(X)\geq q$: contracting the first factor with a
matrix $X$ leaves the bilinear form $(Y,Z)\mapsto\operatorname{tr}(XYZ)$.
The unit tensor has rank-one slices. This property is invariant under
invertible local maps, proving the exact obstruction. In particular it
persists for every equal-size tensor power, since
$\mathrm{MM}_q^{\otimes n}\cong\mathrm{MM}_{q^n}$.

Nevertheless, Wigderson and Zuiddam (2025), Theorem~4.22, p.~58,
record Strassen's theorem
\[
 \widetilde{Q}(\mathrm{MM}_q)=q^2.
\]
Consequently $\langle q^{2n}\rangle$ is a restriction of
$\mathrm{MM}_q^{\otimes(n+o(n))}$. To see the overhead explicitly, write
$\log_{q^2}Q(\mathrm{MM}_q^{\otimes m})=m-e_m$ with
$e_m=o(m)$, and choose $m=n+o(n)$ so that $m-e_m\geq n$.
Coordinate projection supplies the desired unit tensor. Thus an exact
obstruction at every equal-size power, by itself, does not exclude the
asymptotic restriction needed for a price separation.

There is also a direct warning against replacing this missing robust
argument by local flattening-rank continuity. Write
\[
 \mathrm{MM}_2=\sum_{i,j,k\in\{1,2\}}a_{ij}b_{jk}c_{ki}.
\]
Give each diagonal $a$ or $b$ variable weight zero, each off-diagonal
$a$ or $b$ variable weight one, each diagonal $c$ variable weight one,
and each off-diagonal $c$ variable weight zero. The smallest supported
total weight is one. The corresponding degeneration is
$tT+t^3U$, where
\begin{align*}
 T={}&a_{11}b_{11}c_{11}+a_{22}b_{22}c_{22}
      +a_{12}b_{22}c_{21}+a_{21}b_{11}c_{12}\\
    &+a_{11}b_{12}c_{21}+a_{22}b_{21}c_{12},\\
 U={}&a_{12}b_{21}c_{11}+a_{21}b_{12}c_{22}.
\end{align*}
Every local flattening of $T$ has rank four, as is seen from the
linearly independent coefficient slices. Its $c_{11}$ slice is
$a_{11}b_{11}$, of rank one, so $T$ is not in the restriction image of
$\mathrm{MM}_2$. This is the critical-triangle nonclosedness phenomenon
studied by Landsberg, Qi, and Ye (2012), Theorem~3.0.3, pp.~4--5.
The displayed coefficient calculation is valid in every field.

These observations do not disprove a border strengthening of SH-19.
In SH-19 the target lies in the same matrix-multiplication orbit after
blocking matched parties, a stronger hypothesis than conciseness.
They identify exactly why equal-size exact restrictions and preserved
flattening ranks alone are insufficient.

The inspected tensor-network canonical-form theorems do not close this
gap. Molnar et al. (2018), Theorems~2 and~4, pp.~13 and~16--17,
give exact gauge conclusions for equal injective or normal PEPS.
Acuaviva et al. (2022), Theorem~4.11, p.~23, treats equality of
uniform local tensors on all contraction graphs; Proposition~4.20,
p.~30, closes a local gauge orbit when its closure contains a normal
tensor. These are not closedness statements for the image of one
fixed nonuniform $K_{3,3}$ contraction. Section~7, p.~45, explicitly
lists canonical forms for nonuniform PEPS as future work. No border
or $1+o(1)$-overhead obstruction for SH-19 is inferred from them.

\subsection{A failed five-point gate and its exact geometric explanation}

\paragraph{EX-25: every square-root metric is pentagonal.}
Let $\rho$ be any pseudometric and $d=\sqrt\rho$. For any five
points, with repetitions allowed,
\[
 d(x_1,x_2)+\sum_{1\le i<j\le3}d(y_i,y_j)
 \le \sum_{i=1}^2\sum_{j=1}^3d(x_i,y_j).
 \tag{EX25.1}
\]
Here is a direct proof, independent of the computation below.
If $a=d(u,v)\le b=d(v,w)$, the triangle inequality for $\rho$
and elementary algebra give
\[
 d(u,w)\le\sqrt{a^2+b^2}\le b+\frac a2.
 \tag{EX25.2}
\]
Indeed $(b+a/2)^2-(a^2+b^2)=a(b-3a/4)\ge0$.
Relabel the two $x$-points and three $y$-points so that
$a=d(x_1,y_1)$ is a minimum among all six cross distances,
and then so that $b=d(x_2,y_2)\le d(x_2,y_3)$.
In particular $a\le b$. Applying (EX25.2) through the
minimum pair gives
\[
 d(x_1,x_2)+d(y_1,y_2)
 \le d(x_2,y_1)+d(x_1,y_2)+a
 \le d(x_2,y_1)+d(x_1,y_2)+\frac{a+b}{2}.
\]
Applying it to the two remaining pairs gives
\[
 d(y_1,y_3)+d(y_2,y_3)
 \le d(x_1,y_3)+d(x_2,y_3)+\frac{a+b}{2}.
\]
Adding these inequalities proves (EX25.1), including all
zero-distance and repeated-point cases.

This is a special case of Deza and Maehara (1990),
\emph{Metric transforms and Euclidean embeddings},
\emph{Transactions of the American Mathematical Society},
317(2), 661--671,
\href{https://doi.org/10.1090/S0002-9947-1990-0974513-6}
{doi:10.1090/S0002-9947-1990-0974513-6}.
Their Theorem~1, printed pp.~664--665, gives all polygonal
inequalities on at most six terms for $\rho^c$ when
$0<c\le\log_2(3/2)$. More strongly, Corollary~2, printed
p.~667, states that every six-point metric transformed by
such an exponent embeds isometrically in $L_1$.
The primary article text was read through a
\href{https://www.academia.edu/7373668/transactions_of_the_american_mathematical_society_METRIC_TRANSFORMS_AND_EUCLIDEAN_EMBEDDINGS}
{full-text mirror}; the publisher PDF returned HTTP~403.
Only the strict interior exponent $c=1/2$ is needed here, so
the OCR treatment of endpoint inequality signs is immaterial.

In particular, for matrices over any field, every submetric
on at most six points of
$d(A,B)=\sqrt{\operatorname{rank}(A-B)}$ is an $L_1$ metric.
Thus a five-party $K_{3,2}$ asymptotic-cover inequality cannot
exclude this rank square-root metric in any matrix dimension.
The seven-point odd-hypermetric witness in SH-29, obtained by
deleting one negative-character point from SH-20, is minimal in point
count among all non-$L_1$ submetrics of this rank square-root
metric. This assertion concerns $L_1$ inequalities; it does not
say that seven points are necessary for every conceivable
spectral or CAT(0) obstruction.

Before locating the general theorem, an exhaustive integer
certificate was computed for $\operatorname{Mat}_3(\mathbb F_2)$.
Translation sends the two-point part of a pentagon to $(0,Z)$,
and left and right multiplication by invertible matrices sends
$Z$ to $\operatorname{diag}(I_r,0)$ for $r=1,2,3$ while
preserving every rank difference. For each of these three
orbits the program enumerates all $\binom{510}{3}=21978620$
unordered triples disjoint from $\{0,Z\}$, for a total of
$65935860$ triples. Binary Gaussian elimination computes all
ranks exactly. With $S=10^9$, each positive term $\sqrt r$ is
replaced by $\lceil S\sqrt r\rceil/S$ and each negative term
by $\lfloor S\sqrt r\rfloor/S$, using integer square roots.
The maximum computed upper numerators for the objective
``within-part sum minus cross-part sum'' are respectively
\[
 -1171572873,\qquad -343145748,\qquad -632274699.
\]
All are strictly negative. These are certified upper bounds;
they are not asserted to be the exact irrational maxima.
The reproducible program is
\path{programs/ex_search_pentagon_m3.py}, with parameters
\path{inputs/ex_pentagon_m3_parameters.json} and output
\path{outputs/ex_pentagon_m3_results.json}.
Repeated-point cases are already covered by the direct proof.
The computation is retained as an independently checkable failed
approach, rather than used to claim any unsearched global result.
